%% tcpsock.tex
%%
%% Copyright (C) 2000-2018 Simone Piccardi.  Permission is granted to
%% copy, distribute and/or modify this document under the terms of the GNU Free
%% Documentation License, Version 1.1 or any later version published by the
%% Free Software Foundation; with the Invariant Sections being "Un preambolo",
%% with no Front-Cover Texts, and with no Back-Cover Texts.  A copy of the
%% license is included in the section entitled "GNU Free Documentation
%% License".
%%

\chapter{I socket TCP}
\label{cha:TCP_socket}

In questo capitolo tratteremo le basi dei socket TCP, iniziando con una
descrizione delle principali caratteristiche del funzionamento di una
connessione TCP; vedremo poi le varie funzioni che servono alla creazione di
una connessione fra client e server, fornendo alcuni esempi elementari, e
finiremo prendendo in esame l'uso dell'\textit{I/O multiplexing}.


\section{Il funzionamento di una connessione TCP}
\label{sec:TCP_connession}

Prima di entrare nei dettagli delle singole funzioni usate nelle applicazioni
che utilizzano i socket TCP, è fondamentale spiegare alcune delle basi del
funzionamento del protocollo, poiché questa conoscenza è essenziale per
comprendere il comportamento di dette funzioni per questo tipo di socket, ed
il relativo modello di programmazione.

Si ricordi che il protocollo TCP serve a creare degli \textit{stream socket},
cioè una forma di canale di comunicazione che stabilisce una connessione
stabile fra due macchine in rete, in modo che queste possano scambiarsi dei
dati. In questa sezione ci concentreremo sulle modalità con le quali il
protocollo dà inizio e conclude una connessione e faremo inoltre un breve
accenno al significato di alcuni dei vari \textsl{stati} ad essa associati.


\subsection{La creazione della connessione: il \textit{three way handshake}}
\label{sec:TCP_conn_cre}

\itindbeg{three~way~handshake} 

Il processo che porta a creare una connessione TCP viene chiamato
\textit{three way handshake}; la successione tipica degli eventi (e dei
\textsl{segmenti}\footnote{si ricordi che il \textsl{segmento} è l'unità
  elementare di dati trasmessa dal protocollo TCP al livello successivo; tutti
  i segmenti hanno un'intestazione che contiene le informazioni che servono
  allo \textit{stack TCP} (così viene di solito chiamata la parte del kernel
  che realizza il protocollo) per effettuare la comunicazione, fra questi dati
  ci sono una serie di flag usati per gestire la connessione, come SYN, ACK,
  URG, FIN, alcuni di essi, come SYN (che sta per \textit{syncronize})
  corrispondono a funzioni particolari del protocollo e danno il nome al
  segmento, (per maggiori dettagli si veda sez.~\ref{sec:tcp_protocol}).}  di
dati che vengono scambiati) che porta alla creazione di una connessione è la
seguente:
 
\begin{enumerate*}
\item Il server deve essere preparato per accettare le connessioni in arrivo;
  il procedimento si chiama \textsl{apertura passiva} del socket (in inglese
  \textit{passive open}). Questo viene fatto chiamando la sequenza di funzioni
  \func{socket}, \func{bind} e \func{listen}. Completata l'apertura passiva il
  server chiama la funzione \func{accept} e il processo si blocca in attesa di
  connessioni.
  
\item Il client richiede l'inizio della connessione usando la funzione
  \func{connect}, attraverso un procedimento che viene chiamato
  \textsl{apertura attiva}, dall'inglese \textit{active open}. La chiamata di
  \func{connect} blocca il processo e causa l'invio da parte del client di un
  segmento SYN, in sostanza viene inviato al server un pacchetto IP che
  contiene solo le intestazioni di IP e TCP (con il numero di sequenza
  iniziale e il flag SYN) e le opzioni di TCP.
  
\item Il server deve dare ricevuto (l'\textit{acknowledge}) del SYN del
  client, inoltre anche il server deve inviare il suo SYN al client (e
  trasmettere il suo numero di sequenza iniziale) questo viene fatto
  ritrasmettendo un singolo segmento in cui sono impostati entrambi i flag SYN
  e ACK.
  
\item Una volta che il client ha ricevuto l'\textit{acknowledge} dal server la
  funzione \func{connect} ritorna, l'ultimo passo è dare il ricevuto del SYN
  del server inviando un ACK. Alla ricezione di quest'ultimo la funzione
  \func{accept} del server ritorna e la connessione è stabilita.
\end{enumerate*}

Il procedimento viene chiamato \textit{three way handshake} dato che per
realizzarlo devono essere scambiati tre segmenti. In fig.~\ref{fig:TCP_TWH}
si è rappresentata graficamente la sequenza di scambio dei segmenti che
stabilisce la connessione.

% Una analogia citata da R. Stevens per la connessione TCP è quella con il
% sistema del telefono. La funzione \func{socket} può essere considerata
% l'equivalente di avere un telefono. La funzione \func{bind} è analoga al
% dire alle altre persone qual è il proprio numero di telefono perché possano
% chiamare. La funzione \func{listen} è accendere il campanello del telefono
% per sentire le chiamate in arrivo.  La funzione \func{connect} richiede di
% conoscere il numero di chi si vuole chiamare. La funzione \func{accept} è
% quando si risponde al telefono.

\begin{figure}[!htb]
  \centering \includegraphics[width=10cm]{img/three_way_handshake}  
  \caption{Il \textit{three way handshake} del TCP.}
  \label{fig:TCP_TWH}
\end{figure}

\index{numeri~di~sequenza|(}

Si è accennato in precedenza ai \textsl{numeri di sequenza}, che sono anche
riportati in fig.~\ref{fig:TCP_TWH}: per gestire una connessione affidabile
infatti il protocollo TCP prevede nell'intestazione la presenza di un numero a
32 bit (chiamato appunto \textit{sequence number}) che identifica a quale byte
nella sequenza del flusso corrisponde il primo byte della sezione dati
contenuta nel segmento.

Il numero di sequenza di ciascun segmento viene calcolato a partire da un
\textsl{numero di sequenza iniziale} generato in maniera casuale del kernel
all'inizio della connessione e trasmesso con il SYN;
l'\textit{acknowledgement} di ciascun segmento viene effettuato dall'altro
capo della connessione impostando il flag ACK e restituendo nell'apposito
campo dell'intestazione un \textit{acknowledge number}) pari al numero di
sequenza che il ricevente si aspetta di ricevere con il pacchetto successivo;
dato che il primo pacchetto SYN consuma un byte, nel \textit{three way
  handshake} il numero di \textit{acknowledge} è sempre pari al numero di
sequenza iniziale incrementato di uno; lo stesso varrà anche (vedi
fig.~\ref{fig:TCP_close}) per l'\textit{acknowledgement} di un segmento FIN.

\index{numeri~di~sequenza|)}
\itindend{three~way~handshake}


\subsection{Le opzioni TCP.}
\label{sec:TCP_TCP_opt}

Ciascun segmento SYN contiene in genere delle opzioni per il protocollo TCP,
le cosiddette \textit{TCP options}, da non confondere con le opzioni dei
socket TCP che tratteremo in sez.~\ref{sec:sock_tcp_udp_options}. In questo
caso infatti si tratta delle opzioni che vengono trasmesse come parte di un
pacchetto TCP, e non delle funzioni che consentono di impostare i relativi
valori. Queste opzioni vengono inserite fra l'intestazione ed i dati, e
servono a comunicare all'altro capo una serie di parametri utili a regolare la
connessione. Normalmente vengono usate le seguenti opzioni:

\begin{itemize*}
\item \textit{MSS option}, con questa opzione ciascun capo della connessione
  annuncia all'altro il massimo ammontare di dati (MMS sta appunto per
  \textit{Maximum Segment Size}, vedi sez.~\ref{sec:tcp_protocol}) che
  vorrebbe accettare per ciascun segmento nella connessione corrente. È
  possibile leggere e scrivere questo valore attraverso l'opzione del socket
  \const{TCP\_MAXSEG} (vedi sez.~\ref{sec:sock_tcp_udp_options}).
  
\item \textit{window scale option}, il protocollo TCP realizza il controllo di
  flusso attraverso una \textit{advertised window} (la ``\textsl{finestra
    annunciata}'', vedi sez.~\ref{sec:tcp_protocol}) con la quale ciascun capo
  della comunicazione dichiara quanto spazio disponibile ha in memoria per i
  dati. Questo è un numero a 16 bit dell'header, che così può indicare un
  massimo di 65535 byte\footnote{in Linux il massimo è 32767 per evitare
    problemi con alcune realizzazione dello \textit{stack TCP} che usano
    l'aritmetica con segno.} ma alcuni tipi di connessione come quelle ad alta
  velocità (sopra i 45Mbit/sec) e quelle che hanno grandi ritardi nel cammino
  dei pacchetti (come i satelliti) richiedono una finestra più grande per
  poter ottenere il massimo dalla trasmissione.

  Per questo esiste un'altra opzione che indica un fattore di scala da
  applicare al valore della finestra annunciata per la connessione corrente, 
  espresso come numero di bit cui spostare a sinistra il valore della
  finestra annunciata inserito nel pacchetto. Essendo una nuova opzione per
  garantire la compatibilità con delle vecchie realizzazione del protocollo la
  procedura che la attiva prevede come negoziazione che l'altro capo della
  connessione riconosca esplicitamente l'opzione inserendola anche lui nel suo
  SYN di risposta dell'apertura della connessione.

  Con Linux è possibile indicare al kernel di far negoziare il fattore di
  scala in fase di creazione di una connessione tramite la \textit{sysctl}
  \texttt{tcp\_window\_scaling} (vedi sez.~\ref{sec:sock_ipv4_sysctl}). Per
  poter usare questa funzionalità è comunque necessario ampliare le dimensioni
  dei buffer di ricezione e spedizione, cosa che può essere fatta sia a
  livello di sistema con le opportune \textit{sysctl} (vedi
  sez.~\ref{sec:sock_ipv4_sysctl}) che a livello di singoli socket con le
  relative opzioni (vedi sez.~\ref{sec:sock_tcp_udp_options}).

\item \textit{timestamp option}, è anche questa una nuova opzione necessaria
  per le connessioni ad alta velocità per evitare possibili corruzioni di dati
  dovute a pacchetti perduti che riappaiono; anche questa viene negoziata
  all'inizio della connessione come la precedente.
\end{itemize*}

La \textit{MSS option} è generalmente supportata da quasi tutte le
realizzazioni del protocollo, le ultime due opzioni (trattate
nell'\href{http://www.ietf.org/rfc/rfc1323.txt}{RFC~1323}) sono meno comuni;
vengono anche dette \textit{long fat pipe options} dato che questo è il nome
che viene dato alle connessioni caratterizzate da alta velocità o da ritardi
elevati. In ogni caso Linux supporta pienamente entrambe queste opzioni
aggiuntive.


\subsection{La terminazione della connessione}
\label{sec:TCP_conn_term}

Mentre per la creazione di una connessione occorre un interscambio di tre
segmenti, la procedura di chiusura ne richiede normalmente quattro. In questo
caso la successione degli eventi è la seguente:

\begin{enumerate*}
\item Un processo ad uno dei due capi chiama la funzione \func{close}, dando
  l'avvio a quella che viene chiamata \textsl{chiusura attiva} (o
  \textit{active close}). Questo comporta l'emissione di un segmento FIN, che
  serve ad indicare che si è finito con l'invio dei dati sulla connessione.
  
\item L'altro capo della connessione riceve il segmento FIN e dovrà eseguire
  la \textsl{chiusura passiva} (o \textit{passive close}). Al FIN, come ad
  ogni altro pacchetto, viene risposto con un ACK, inoltre il ricevimento del
  FIN viene segnalato al processo che ha aperto il socket (dopo che ogni altro
  eventuale dato rimasto in coda è stato ricevuto) come un
  \textit{end-of-file} sulla lettura: questo perché il ricevimento di un
  segmento FIN significa che non si riceveranno altri dati sulla connessione.
  
\item Una volta rilevata l'\textit{end-of-file} anche il secondo processo
  chiamerà la funzione \func{close} sul proprio socket, causando l'emissione
  di un altro segmento FIN.

\item L'altro capo della connessione riceverà il segmento FIN conclusivo e
  risponderà con un ACK.
\end{enumerate*}

Dato che in questo caso sono richiesti un FIN ed un ACK per ciascuna direzione
normalmente i segmenti scambiati sono quattro.  Questo non è vero sempre
giacché in alcune situazioni il FIN del passo 1 è inviato insieme a dei dati.
Inoltre è possibile che i segmenti inviati nei passi 2 e 3 dal capo che
effettua la chiusura passiva, siano accorpati in un singolo segmento. Come per
il SYN anche il FIN occupa un byte nel numero di sequenza, per cui l'ACK
riporterà un \textit{acknowledge number} incrementato di uno. In
fig.~\ref{fig:TCP_close} si è rappresentata graficamente la sequenza di
scambio dei segmenti che conclude la connessione.

\begin{figure}[!htb]
  \centering \includegraphics[width=10cm]{img/tcp_close}  
  \caption{La chiusura di una connessione TCP.}
  \label{fig:TCP_close}
\end{figure}

Si noti che, nella sequenza di chiusura, fra i passi 2 e 3, è in teoria
possibile che si mantenga un flusso di dati dal capo della connessione che
deve ancora eseguire la chiusura passiva a quello che sta eseguendo la
chiusura attiva.  Nella sequenza indicata i dati verrebbero persi, dato che si
è chiuso il socket dal lato che esegue la chiusura attiva; esistono tuttavia
situazioni in cui si vuole poter sfruttare questa possibilità, usando una
procedura che è chiamata \textit{half-close}; torneremo su questo aspetto e su
come utilizzarlo in sez.~\ref{sec:TCP_shutdown}, quando parleremo della
funzione \func{shutdown}.

La emissione del segmento FIN avviene quando il socket viene chiuso, questo
però non avviene solo per la chiamata esplicita della funzione \func{close},
ma anche alla terminazione di un processo, quando tutti i file vengono chiusi.
Questo comporta ad esempio che se un processo viene terminato da un segnale
tutte le connessioni aperte verranno chiuse.

Infine occorre sottolineare che, benché nella figura (e nell'esempio che
vedremo più avanti in sez.~\ref{sec:TCP_echo}) sia stato il client ad eseguire
la chiusura attiva, nella realtà questa può essere eseguita da uno qualunque
dei due capi della comunicazione (come nell'esempio di
fig.~\ref{fig:TCP_daytime_iter_server_code}), e che anche se il caso più
comune resta quello del client, ci sono alcuni servizi, il più noto dei quali
è l'HTTP 1.0 (con le versioni successive il default è cambiato) per i quali è
il server ad effettuare la chiusura attiva.


\subsection{Un esempio di connessione}
\label{sec:TCP_conn_dia}

Come abbiamo visto le operazioni del TCP nella creazione e conclusione di una
connessione sono piuttosto complesse, ed abbiamo esaminato soltanto quelle
relative ad un andamento normale.  In sez.~\ref{sec:TCP_states} vedremo con
maggiori dettagli che una connessione può assumere vari stati, che ne
caratterizzano il funzionamento, e che sono quelli che vengono riportati dal
comando \cmd{netstat}, per ciascun socket TCP aperto, nel campo
\textit{State}.

Non possiamo affrontare qui una descrizione completa del funzionamento del
protocollo; un approfondimento sugli aspetti principali si trova in
sez.~\ref{sec:tcp_protocol}, ma per una trattazione completa il miglior
riferimento resta \cite{TCPIll1}. Qui ci limiteremo a descrivere brevemente un
semplice esempio di connessione e le transizioni che avvengono nei due casi
appena citati (creazione e terminazione della connessione).

In assenza di connessione lo stato del TCP è \texttt{CLOSED}; quando una
applicazione esegue una apertura attiva il TCP emette un SYN e lo stato
diventa \texttt{SYN\_SENT}; quando il TCP riceve la risposta del SYN$+$ACK
emette un ACK e passa allo stato \texttt{ESTABLISHED}; questo è lo stato
finale in cui avviene la gran parte del trasferimento dei dati.

Dal lato server in genere invece il passaggio che si opera con l'apertura
passiva è quello di portare il socket dallo stato \texttt{CLOSED} allo
stato \texttt{LISTEN} in cui vengono accettate le connessioni.

Dallo stato \texttt{ESTABLISHED} si può uscire in due modi; se un'applicazione
chiama la funzione \func{close} prima di aver ricevuto un
\textit{end-of-file} (chiusura attiva) la transizione è verso lo stato
\texttt{FIN\_WAIT\_1}; se invece l'applicazione riceve un FIN nello stato
\texttt{ESTABLISHED} (chiusura passiva) la transizione è verso lo stato
\texttt{CLOSE\_WAIT}.

In fig.~\ref{fig:TCP_conn_example} è riportato lo schema dello scambio dei
pacchetti che avviene per una un esempio di connessione, insieme ai vari stati
che il protocollo viene ad assumere per i due lati, server e client.

\begin{figure}[!htb]
  \centering \includegraphics[width=9cm]{img/tcp_connection}  
  \caption{Schema dello scambio di pacchetti per un esempio di connessione.}
  \label{fig:TCP_conn_example}
\end{figure}

La connessione viene iniziata dal client che annuncia una \textit{Maximum
  Segment Size} di 1460, un valore tipico con Linux per IPv4 su Ethernet, il
server risponde con lo stesso valore (ma potrebbe essere anche un valore
diverso).

Una volta che la connessione è stabilita il client scrive al server una
richiesta (che assumiamo stare in un singolo segmento, cioè essere minore dei
1460 byte annunciati dal server), quest'ultimo riceve la richiesta e
restituisce una risposta (che di nuovo supponiamo stare in un singolo
segmento). Si noti che l'\textit{acknowledge} della richiesta è mandato
insieme alla risposta: questo viene chiamato \textit{piggybacking} ed avviene
tutte le volte che il server è sufficientemente veloce a costruire la
risposta; in caso contrario si avrebbe prima l'emissione di un ACK e poi
l'invio della risposta.

Infine si ha lo scambio dei quattro segmenti che terminano la connessione
secondo quanto visto in sez.~\ref{sec:TCP_conn_term}; si noti che il capo della
connessione che esegue la chiusura attiva entra nello stato
\texttt{TIME\_WAIT}, sul cui significato torneremo fra poco.

È da notare come per effettuare uno scambio di due pacchetti (uno di richiesta
e uno di risposta) TCP necessiti di ulteriori otto segmenti, se invece si
fosse usato UDP sarebbero stati sufficienti due soli pacchetti. Questo è il
costo che occorre pagare per avere l'affidabilità garantita da TCP, se si
fosse usato UDP si sarebbe dovuto trasferire la gestione di tutta una serie di
dettagli (come la verifica della ricezione dei pacchetti) dal livello del
trasporto all'interno dell'applicazione.

Quello che è bene sempre tenere presente è allora quali sono le esigenze che
si hanno in una applicazione di rete, perché non è detto che TCP sia la
miglior scelta in tutti i casi (ad esempio se si devono solo scambiare dati
già organizzati in piccoli pacchetti l'overhead aggiunto può essere eccessivo)
per questo esistono applicazioni che usano UDP e lo fanno perché nel caso
specifico le sue caratteristiche di velocità e compattezza nello scambio dei
dati rispondono meglio alle esigenze che devono essere affrontate.


\subsection{Lo stato \texttt{TIME\_WAIT}}
\label{sec:TCP_time_wait}

Come riportato da Stevens in \cite{UNP1} lo stato \texttt{TIME\_WAIT} è
probabilmente uno degli aspetti meno compresi del protocollo TCP, è infatti
comune trovare domande su come sia possibile evitare che un'applicazione resti
in questo stato lasciando attiva una connessione ormai conclusa; la risposta è
che non deve essere fatto, ed il motivo cercheremo di spiegarlo adesso.

\itindbeg{Maximum~Segment~Lifetime~(MSL)}
Come si è visto nell'esempio precedente (vedi fig.~\ref{fig:TCP_conn_example})
\texttt{TIME\_WAIT} è lo stato finale in cui il capo di una connessione che
esegue la chiusura attiva resta prima di passare alla chiusura definitiva
della connessione. Il tempo in cui l'applicazione resta in questo stato deve
essere due volte la \textit{Maximum Segment Lifetime} (da qui in avanti
abbreviata in MSL).

La \textit{Maximum Segment Lifetime} è la stima del massimo periodo di tempo
in secondi che un pacchetto IP può vivere sulla rete. Questo tempo è limitato
perché ogni pacchetto IP può essere ritrasmesso dai router un numero massimo
di volte (detto \textit{hop limit}).  Il numero di ritrasmissioni consentito è
indicato dal campo TTL dell'intestazione di IP (per maggiori dettagli vedi
sez.~\ref{sec:ip_protocol}), e viene decrementato ad ogni passaggio da un
router; quando si annulla il pacchetto viene scartato.  Siccome il numero è ad
8 bit il numero massimo di ``\textsl{salti}'' è di 255, pertanto anche se il
TTL (da \textit{time to live}) non è propriamente un limite sul tempo, sulla
sua base si può stimare che un pacchetto IP non possa restare nella rete per
più un certo numero di secondi, che costituisce la MSL.
\itindend{Maximum~Segment~Lifetime~(MSL)}

Ogni realizzazione del TCP deve scegliere un valore per la MSL;
l'\href{http://www.ietf.org/rfc/rfc1122.txt}{RFC~1122} raccomanda 2 minuti,
mentre Linux usa 30 secondi, questo comporta una durata dello stato
\texttt{TIME\_WAIT} che a seconda delle realizzazioni può variare fra 1 a 4
minuti.  Lo stato \texttt{TIME\_WAIT} viene utilizzato dal protocollo per due
motivi principali:
\begin{enumerate*}
\item effettuare in maniera affidabile la terminazione della connessione
  in entrambe le direzioni.
\item consentire l'eliminazione dei segmenti duplicati dalla rete. 
\end{enumerate*}

Il punto è che entrambe le ragioni sono importanti, anche se spesso si fa
riferimento solo alla prima; ma è solo se si tiene conto della seconda che si
capisce il perché della scelta di un tempo pari al doppio della MSL come
durata di questo stato.

Il primo dei due motivi precedenti si può capire tornando a
fig.~\ref{fig:TCP_conn_example}: assumendo che l'ultimo ACK della sequenza
(quello del capo che ha eseguito la chiusura attiva) venga perso, chi esegue
la chiusura passiva non ricevendo risposta rimanderà un ulteriore FIN, per
questo motivo chi esegue la chiusura attiva deve mantenere lo stato della
connessione per essere in grado di reinviare l'ACK e chiuderla correttamente.
Se non fosse così la risposta sarebbe un RST (un altro tipo si segmento) che
verrebbe interpretato come un errore.

Se il TCP deve poter chiudere in maniera pulita entrambe le direzioni della
connessione allora deve essere in grado di affrontare la perdita di uno
qualunque dei quattro segmenti che costituiscono la chiusura. Per questo
motivo un socket deve rimanere attivo nello stato \texttt{TIME\_WAIT} anche
dopo l'invio dell'ultimo ACK, per potere essere in grado di gestirne
l'eventuale ritrasmissione, in caso esso venga perduto.

Il secondo motivo è più complesso da capire, e necessita di una spiegazione
degli scenari in cui può accadere che i pacchetti TCP si possano perdere nella
rete o restare intrappolati, per poi riemergere in un secondo tempo. 

Il caso più comune in cui questo avviene è quello di anomalie
nell'instradamento; può accadere cioè che un router smetta di funzionare o che
una connessione fra due router si interrompa. In questo caso i protocolli di
instradamento dei pacchetti possono impiegare diverso tempo (anche dell'ordine
dei minuti) prima di trovare e stabilire un percorso alternativo per i
pacchetti. Nel frattempo possono accadere casi in cui un router manda i
pacchetti verso un altro e quest'ultimo li rispedisce indietro, o li manda ad
un terzo router che li rispedisce al primo, si creano cioè dei circoli (i
cosiddetti \textit{routing loop}) in cui restano intrappolati i pacchetti.

Se uno di questi pacchetti intrappolati è un segmento TCP, chi l'ha inviato,
non ricevendo un ACK in risposta, provvederà alla ritrasmissione e se nel
frattempo sarà stata stabilita una strada alternativa il pacchetto ritrasmesso
giungerà a destinazione.

Ma se dopo un po' di tempo (che non supera il limite dell'MSL, dato che
altrimenti verrebbe ecceduto il TTL) l'anomalia viene a cessare, il circolo di
instradamento viene spezzato i pacchetti intrappolati potranno essere inviati
alla destinazione finale, con la conseguenza di avere dei pacchetti duplicati;
questo è un caso che il TCP deve essere in grado di gestire.

Allora per capire la seconda ragione per l'esistenza dello stato
\texttt{TIME\_WAIT} si consideri il caso seguente: si supponga di avere una
connessione fra l'IP \texttt{195.110.112.236} porta 1550 e l'IP
\texttt{192.84.145.100} porta 22 (affronteremo il significato delle porte
nella prossima sezione), che questa venga chiusa e che poco dopo si
ristabilisca la stessa connessione fra gli stessi IP sulle stesse porte
(quella che viene detta, essendo gli stessi porte e numeri IP, una nuova
\textsl{incarnazione} della connessione precedente); in questo caso ci si
potrebbe trovare con dei pacchetti duplicati relativi alla precedente
connessione che riappaiono nella nuova.

Ma fintanto che il socket non è chiuso una nuova incarnazione non può essere
creata: per questo un socket TCP resta sempre nello stato \texttt{TIME\_WAIT}
per un periodo di 2 volte la MSL, prima per attendere MSL secondi per essere
sicuri che tutti i pacchetti duplicati in arrivo siano stati ricevuti e
scartati, o che nel frattempo siano stati eliminati dalla rete, poi altri MSL
secondi per essere sicuri che lo stesso avvenga per le risposte nella
direzione opposta.

In questo modo, prima che venga creata una nuova connessione, il protocollo
TCP si assicura che tutti gli eventuali segmenti residui di una precedente
connessione, che potrebbero causare disturbi, siano stati eliminati dalla
rete.


\subsection{I numeri di porta}
\label{sec:TCP_port_num}

In un ambiente multitasking in un dato momento più processi devono poter usare
sia UDP che TCP, e ci devono poter essere più connessioni in contemporanea.
Per poter tenere distinte le diverse connessioni entrambi i protocolli usano i
\textsl{numeri di porta}, che fanno parte, come si può vedere in
sez.~\ref{sec:sock_sa_ipv4} e sez.~\ref{sec:sock_sa_ipv6} pure delle strutture
degli indirizzi del socket.

Quando un client contatta una macchina server deve poter identificare con
quale dei vari possibili programmi server attivi intende parlare. Sia TCP che
UDP definiscono un gruppo di \textsl{porte conosciute} (le cosiddette
\textit{well-known port}) che identificano una serie di servizi noti (ad
esempio la porta 22 identifica il servizio SSH) effettuati da appositi
programmi server che rispondono alle connessioni verso tali porte.

D'altra parte un client non ha necessità di usare dalla sua parte un numero di
porta specifico, per cui in genere vengono usate le cosiddette \textsl{porte
  effimere} (o \textit{ephemeral ports}) cioè porte a cui non è assegnato
nessun servizio noto e che vengono assegnate automaticamente dal kernel alla
creazione della connessione. Queste sono dette effimere in quanto vengono
usate solo per la durata della connessione, e l'unico requisito che deve
essere soddisfatto è che ognuna di esse sia assegnata in maniera univoca.

La lista delle porte conosciute è definita
dall'\href{http://www.ietf.org/rfc/rfc1700.txt}{RFC~1700} che contiene
l'elenco delle porte assegnate dalla IANA (la \textit{Internet Assigned Number
  Authority}) ma l'elenco viene costantemente aggiornato e pubblicato su
Internet (una versione aggiornata si può trovare all'indirizzo
\url{http://www.iana.org/assignments/port-numbers}); inoltre in un sistema
unix-like un analogo elenco viene mantenuto nel file \conffile{/etc/services},
con la corrispondenza fra i vari numeri di porta ed il nome simbolico del
servizio.  I numeri sono divisi in tre intervalli:

\begin{enumerate*}
\item \textsl{le porte note}. I numeri da 0 a 1023. Queste sono controllate e
  assegnate dalla IANA. Se è possibile la stessa porta è assegnata allo stesso
  servizio sia su UDP che su TCP (ad esempio la porta 22 è assegnata a SSH su
  entrambi i protocolli, anche se viene usata solo dal TCP).
  
\item \textsl{le porte registrate}. I numeri da 1024 a 49151. Queste porte non
  sono controllate dalla IANA, che però registra ed elenca chi usa queste
  porte come servizio agli utenti. Come per le precedenti si assegna una porta
  ad un servizio sia per TCP che UDP anche se poi il servizio è effettuato
  solo su TCP. Ad esempio \textit{X Window} usa le porte TCP e UDP dal 6000 al
  6063 anche se il protocollo viene usato solo con TCP.
  
\item \textsl{le porte private} o \textsl{dinamiche}. I numeri da 49152 a
  65535. La IANA non dice nulla riguardo a queste porte che pertanto
  sono i candidati naturali ad essere usate come porte effimere.
\end{enumerate*}

In realtà rispetto a quanto indicato
nell'\href{http://www.ietf.org/rfc/rfc1700.txt}{RFC~1700} i vari sistemi hanno
fatto scelte diverse per le porte effimere, in particolare in
fig.~\ref{fig:TCP_port_alloc} sono riportate quelle di BSD e Linux.

\begin{figure}[!htb]
  \centering \includegraphics[width=13cm]{img/port_alloc}  
  \caption{Allocazione dei numeri di porta.}
  \label{fig:TCP_port_alloc}
\end{figure}

I sistemi Unix hanno inoltre il concetto di \textsl{porte riservate}, che
corrispondono alle porte con numero minore di 1024 e coincidono quindi con le
\textsl{porte note}. La loro caratteristica è che possono essere assegnate a
un socket solo da un processo con i privilegi di amministratore, per far sì
che solo l'amministratore possa allocare queste porte per far partire i
relativi servizi.

Le \textsl{glibc} definiscono in \headfile{netinet/in.h} le costanti
\constd{IPPORT\_RESERVED} e \constd{IPPORT\_USERRESERVED}, in cui la prima
(che vale 1024) indica il limite superiore delle porte riservate, e la seconda
(che vale 5000) il limite inferiore delle porte a disposizione degli utenti.
La convenzione vorrebbe che le porte \textsl{effimere} siano allocate fra
questi due valori. Nel caso di Linux questo è vero solo in uno dei due casi di
fig.~\ref{fig:TCP_port_alloc}, e la scelta fra i due possibili intervalli
viene fatta dinamicamente dal kernel a seconda della memoria disponibile per
la gestione delle relative tabelle.

Si tenga conto poi che ci sono alcuni client, in particolare \cmd{rsh} e
\cmd{rlogin}, che richiedono una connessione su una porta riservata anche dal
lato client come parte dell'autenticazione, contando appunto sul fatto che
solo l'amministratore può usare queste porte. Data l'assoluta inconsistenza in
termini di sicurezza di un tale metodo, al giorno d'oggi esso è in completo
disuso.

Data una connessione TCP, ma la cosa vale anche per altri protocolli del
livello di trasporto come UDP, si suole chiamare \itindex{socket~pair}
\textit{socket pair}\footnote{da non confondere con la coppia di socket della
  omonima funzione \func{socketpair} di sez.~\ref{sec:ipc_socketpair} che
  fanno riferimento ad una coppia di socket sulla stessa macchina, non ai capi
  di una connessione TCP.} la combinazione dei quattro numeri che definiscono
i due capi della connessione e cioè l'indirizzo IP locale e la porta TCP
locale, e l'indirizzo IP remoto e la porta TCP remota.  Questa combinazione,
che scriveremo usando una notazione del tipo (\texttt{195.110.112.152:22},
\texttt{192.84.146.100:20100}), identifica univocamente una connessione su
Internet.  Questo concetto viene di solito esteso anche a UDP, benché in
questo caso non abbia senso parlare di connessione. L'utilizzo del programma
\cmd{netstat} permette di visualizzare queste informazioni nei campi
\textit{Local Address} e \textit{Foreing Address}.

Per capire meglio l'uso delle porte e come vengono utilizzate quando si ha a
che fare con un'applicazione client/server (come quelle che descriveremo in
sez.~\ref{sec:TCP_daytime_application} e sez.~\ref{sec:TCP_echo_application})
esamineremo cosa accade con le connessioni nel caso di un server TCP che deve
gestire connessioni multiple.

Se eseguiamo un \cmd{netstat} su una macchina di prova (il cui indirizzo sia
\texttt{195.110.112.152}) potremo avere un risultato del tipo:
\begin{Terminal}
Active Internet connections (servers and established)
Proto Recv-Q Send-Q Local Address           Foreign Address         State      
tcp        0      0 0.0.0.0:22              0.0.0.0:*               LISTEN
tcp        0      0 0.0.0.0:25              0.0.0.0:*               LISTEN
tcp        0      0 127.0.0.1:53            0.0.0.0:*               LISTEN
\end{Terminal}
essendo presenti e attivi un server SSH, un server di posta e un DNS per il
caching locale.

Questo ci mostra ad esempio che il server SSH ha compiuto un'apertura passiva,
mettendosi in ascolto sulla porta 22 riservata a questo servizio, e che si è
posto in ascolto per connessioni provenienti da uno qualunque degli indirizzi
associati alle interfacce locali. La notazione \texttt{0.0.0.0} usata da
\cmd{netstat} è equivalente all'asterisco utilizzato per il numero di porta,
indica il valore generico, e corrisponde al valore \const{INADDR\_ANY}
definito in \headfiled{arpa/inet.h} (vedi \ref{tab:TCP_ipv4_addr}).

Inoltre si noti come la porta e l'indirizzo di ogni eventuale connessione
esterna non sono specificati; in questo caso la \textit{socket pair} associata
al socket potrebbe essere indicata come (\texttt{*:22}, \texttt{*:*}), usando
anche per gli indirizzi l'asterisco come carattere che indica il valore
generico.

Dato che in genere una macchina è associata ad un solo indirizzo IP, ci si può
chiedere che senso abbia l'utilizzo dell'indirizzo generico per specificare
l'indirizzo locale; ma a parte il caso di macchine che hanno più di un
indirizzo IP (il cosiddetto \textit{multihoming}) esiste sempre anche
l'indirizzo di \textit{loopback}, per cui con l'uso dell'indirizzo generico si
possono accettare connessioni indirizzate verso uno qualunque degli indirizzi
IP presenti. Ma, come si può vedere nell'esempio con il DNS che è in ascolto
sulla porta 53, è possibile anche restringere l'accesso ad uno specifico
indirizzo, cosa che nel caso è fatta accettando solo connessioni che arrivino
sull'interfaccia di \textit{loopback}.

Una volta che ci si vorrà collegare a questa macchina da un'altra, per esempio
quella con l'indirizzo \texttt{192.84.146.100}, si dovrà lanciare su
quest'ultima un client \cmd{ssh} per creare una connessione, e il kernel gli
assocerà una porta effimera (ad esempio la 21100), per cui la connessione sarà
espressa dalla \textit{socket pair} (\texttt{192.84.146.100:21100},
\texttt{195.110.112.152:22}).

Alla ricezione della richiesta dal client il server creerà un processo figlio
per gestire la connessione, se a questo punto eseguiamo nuovamente il
programma \cmd{netstat} otteniamo come risultato:
\begin{Terminal}
Active Internet connections (servers and established)
Proto Recv-Q Send-Q Local Address           Foreign Address         State      
tcp        0      0 0.0.0.0:22              0.0.0.0:*               LISTEN
tcp        0      0 0.0.0.0:25              0.0.0.0:*               LISTEN
tcp        0      0 127.0.0.1:53            0.0.0.0:*               LISTEN
tcp        0      0 195.110.112.152:22      192.84.146.100:21100    ESTABLISHED
\end{Terminal}

Come si può notare il server è ancora in ascolto sulla porta 22, però adesso
c'è un nuovo socket (con lo stato \texttt{ESTABLISHED}) che utilizza anch'esso
la porta 22, ed ha specificato l'indirizzo locale, questo è il socket con cui
il processo figlio gestisce la connessione mentre il padre resta in ascolto
sul socket originale.

Se a questo punto lanciamo un'altra volta il client \cmd{ssh} per una seconda
connessione quello che otterremo usando \cmd{netstat} sarà qualcosa del
genere:
\begin{Terminal}
Active Internet connections (servers and established)
Proto Recv-Q Send-Q Local Address           Foreign Address         State      
tcp        0      0 0.0.0.0:22              0.0.0.0:*               LISTEN
tcp        0      0 0.0.0.0:25              0.0.0.0:*               LISTEN
tcp        0      0 127.0.0.1:53            0.0.0.0:*               LISTEN
tcp        0      0 195.110.112.152:22      192.84.146.100:21100    ESTABLISHED
tcp        0      0 195.110.112.152:22      192.84.146.100:21101    ESTABLISHED
\end{Terminal}
cioè il client effettuerà la connessione usando un'altra porta effimera: con
questa sarà aperta la connessione, ed il server creerà un altro processo
figlio per gestirla.

Tutto ciò mostra come il TCP, per poter gestire le connessioni con un server
concorrente, non può suddividere i pacchetti solo sulla base della porta di
destinazione, ma deve usare tutta l'informazione contenuta nella
\textit{socket pair}, compresa la porta dell'indirizzo remoto.  E se andassimo
a vedere quali sono i processi (ad esempio con il comando \cmd{fuser}, o con
\cmd{lsof}, o usando l'opzione \texttt{-p}) a cui fanno riferimento i vari
socket vedremmo che i pacchetti che arrivano dalla porta remota 21100 vanno al
primo figlio e quelli che arrivano alla porta 21101 al secondo.


\section{Le funzioni di base per la gestione dei socket}
\label{sec:TCP_functions}

In questa sezione descriveremo in maggior dettaglio le varie funzioni che
vengono usate per la gestione di base dei socket TCP, non torneremo però sulla
funzione \func{socket}, che è già stata esaminata accuratamente nel capitolo
precedente in sez.~\ref{sec:sock_creation}. Pur trattandole principalmente dal
punto di vista dei socket TCP, daremo brevemente conto del loro comportamento
anche per altri tipi di socket.


\subsection{La funzione \func{bind}}
\label{sec:TCP_func_bind}

La funzione \funcd{bind} assegna un indirizzo locale ad un
socket,\footnote{nel nostro caso la utilizzeremo per socket TCP, ma la
  funzione è generica e deve essere usata per qualunque tipo di socket
  \const{SOCK\_STREAM} prima che questo possa accettare connessioni.} ed è
usata cioè per specificare la prima parte dalla \textit{socket pair}.  Viene
usata sul lato server per specificare la porta (e gli eventuali indirizzi
locali) su cui poi ci si porrà in ascolto. Il prototipo della funzione è il
seguente:


\begin{funcproto}{
\fhead{sys/socket.h}
\fdecl{int bind(int sockfd, const struct sockaddr *serv\_addr, socklen\_t
  addrlen)} 
\fdesc{Assegna un indirizzo ad un socket.} 
}

{La funzione ritorna $0$ in caso di successo e $-1$ per un errore, nel qual
  caso \var{errno} assumerà uno dei valori: 
  \begin{errlist}
  \item[\errcode{EACCES}] si è cercato di usare una porta riservata senza
    sufficienti privilegi.
  \item[\errcode{EADDRINUSE}] qualche altro socket sta già usando l'indirizzo
    richiesto oppure quando non si è richiesta un porta specifica per avere
    una porta effimera non ve ne sono di disponibili nell'intervallo ad esse
    riservato correntemente in uso.
  \item[\errcode{EBADF}] il file descriptor non è valido.
  \item[\errcode{EINVAL}] il socket ha già un indirizzo assegnato,
    o \param{addrlen} è sbagliata, o \param{serv\_addr} non è valido per il
    dominio.
  \item[\errcode{ENOTSOCK}] il file descriptor non è associato ad un socket.
  \end{errlist}
  ed \errval{EFAULT} nel suo significato generico; inoltre per i socket di
  tipo \const{AF\_UNIX}:
  \begin{errlist}
  \item[\errcode{EADDRNOTAVAIL}] il tipo di indirizzo specificato non è
    disponibile.
  \end{errlist}
  ed \errval{ELOOP}, \errval{ENAMETOOLONG}, \errval{ENOENT}, \errval{ENOMEM}, 
  \errval{ENOTDIR} e \errval{EROFS} nel loro significato generico.}
\end{funcproto}

Il primo argomento è un file descriptor ottenuto da una precedente chiamata a
\func{socket}, mentre il secondo e terzo argomento sono rispettivamente
l'indirizzo (locale) del socket e la dimensione della struttura che lo
contiene, secondo quanto già trattato in sez.~\ref{sec:sock_sockaddr}. 

Con i socket TCP la chiamata \func{bind} permette di specificare l'indirizzo,
la porta, entrambi o nessuno dei due. In genere i server utilizzano una porta
nota che assegnano all'avvio, se questo non viene fatto è il kernel a
scegliere una porta effimera quando vengono eseguite la funzioni
\func{connect} o \func{listen}, ma se questo è normale per il client non lo è
per il server\footnote{un'eccezione a tutto ciò sono i server che usano RPC.
  In questo caso viene fatta assegnare dal kernel una porta effimera che poi
  viene registrata presso il \textit{portmapper}; quest'ultimo è un altro
  demone che deve essere contattato dai client per ottenere la porta effimera
  su cui si trova il server.} che in genere viene identificato dalla porta su
cui risponde (l'elenco di queste porte, e dei relativi servizi, è in
\conffile{/etc/services}).

Con \func{bind} si può assegnare un indirizzo IP specifico ad un socket,
purché questo appartenga ad una interfaccia della macchina. Per un client TCP
questo diventerà l'indirizzo sorgente usato per i tutti i pacchetti inviati
sul socket, mentre per un server TCP questo restringerà l'accesso al socket
solo alle connessioni che arrivano verso tale indirizzo.

Normalmente un client non specifica mai l'indirizzo di un socket, ed il kernel
sceglie l'indirizzo di origine quando viene effettuata la connessione, sulla
base dell'interfaccia usata per trasmettere i pacchetti, che dipenderà dalle
regole di instradamento usate per raggiungere il server (è comunque possibile
impostarlo in maniera specifica con i comandi di gestione avanzata del
routing, vedi sez.~7.3.4 di \cite{AGL}). Se un server non specifica il suo
indirizzo locale il kernel userà come indirizzo di origine l'indirizzo di
destinazione specificato dal SYN del client.

Per specificare un indirizzo generico, con IPv4 si usa il valore
\const{INADDR\_ANY}, il cui valore, come accennato in
sez.~\ref{sec:sock_sa_ipv4}, è pari a zero; nell'esempio
fig.~\ref{fig:TCP_daytime_iter_server_code} si è usata un'assegnazione
immediata del tipo: \includecodesnip{listati/serv_addr_sin_addr.c}

Si noti che si è usato \func{htonl} per assegnare il valore
\const{INADDR\_ANY}, anche se, essendo questo nullo, il riordinamento è
inutile.  Si tenga presente comunque che tutte le costanti \val{INADDR\_}
(riportate in tab.~\ref{tab:TCP_ipv4_addr}) sono definite secondo
l'\textit{endianness} della macchina, ed anche se esse possono essere
invarianti rispetto all'ordinamento dei bit, è comunque buona norma usare
sempre la funzione \func{htonl}.

\begin{table}[htb]
  \centering
  \footnotesize
  \begin{tabular}[c]{|l|l|}
    \hline
    \textbf{Costante} & \textbf{Significato} \\
    \hline
    \hline
    \constd{INADDR\_ANY}      & Indirizzo generico (\texttt{0.0.0.0})\\
    \constd{INADDR\_BROADCAST}& Indirizzo di \textit{broadcast}.\\ 
    \constd{INADDR\_LOOPBACK} & Indirizzo di \textit{loopback}
                                (\texttt{127.0.0.1}).\\ 
    \constd{INADDR\_NONE}     & Indirizzo errato.\\
    \hline    
  \end{tabular}
  \caption{Costanti di definizione di alcuni indirizzi generici per IPv4.}
  \label{tab:TCP_ipv4_addr}
\end{table}

L'esempio precedente funziona correttamente con IPv4 poiché che l'indirizzo è
rappresentabile anche con un intero a 32 bit; non si può usare lo stesso
metodo con IPv6, in cui l'indirizzo deve necessariamente essere specificato
con una struttura, perché il linguaggio C non consente l'uso di una struttura
costante come operando a destra in una assegnazione.

Per questo motivo nell'header \headfile{netinet/in.h} è definita una variabile
\var{in6addr\_any} (dichiarata come \dirct{extern}, ed inizializzata dal
sistema al valore \constd{IN6ADRR\_ANY\_INIT}) che permette di effettuare una
assegnazione del tipo: \includecodesnip{listati/serv_addr_sin6_addr.c} in
maniera analoga si può utilizzare la variabile \var{in6addr\_loopback} per
indicare l'indirizzo di \textit{loopback}, che a sua volta viene inizializzata
staticamente a \constd{IN6ADRR\_LOOPBACK\_INIT}.


\subsection{La funzione \func{connect}}
\label{sec:TCP_func_connect}

La funzione \funcd{connect} è usata da un client TCP per stabilire la
connessione con un server TCP, ma la funzione è generica e supporta vari tipi
di socket.  La differenza è che per socket senza connessione come quelli di
tipo \const{SOCK\_DGRAM} la sua chiamata si limiterà ad impostare l'indirizzo
dal quale e verso il quale saranno inviati e ricevuti i pacchetti, mentre per
socket di tipo \const{SOCK\_STREAM} o \const{SOCK\_SEQPACKET} essa attiverà
la procedura di avvio della connessione (nel caso del TCP il \textit{three way
  handshake}). Il suo prototipo è il seguente:


\begin{funcproto}{
\fhead{sys/socket.h}
\fdecl{int connect(int sockfd, const struct sockaddr *servaddr, socklen\_t
    addrlen)}
\fdesc{Stabilisce una connessione fra due socket.} 
}

{La funzione ritorna $0$ in caso di successo e $-1$ per un errore, nel qual
  caso \var{errno} assumerà uno dei valori: 
  \begin{errlist}
  \item[\errcode{EACCES}, \errcode{EPERM}] si è tentato di eseguire una
    connessione ad un indirizzo \textit{broadcast} senza che il socket fosse
    stato abilitato per il \textit{broadcast}.
  \item[\errcode{EAFNOSUPPORT}] l'indirizzo non ha una famiglia di indirizzi
    corretta nel relativo campo.
  \item[\errcode{EAGAIN}] non ci sono più porte locali libere. 
  \item[\errcode{EALREADY}] il socket è non bloccante (vedi
    sez.~\ref{sec:file_noblocking}) e un tentativo precedente di connessione
    non si è ancora concluso.
  \item[\errcode{ECONNREFUSED}] non c'è nessuno in ascolto sull'indirizzo
    remoto.
  \item[\errcode{EINPROGRESS}] il socket è non bloccante (vedi
    sez.~\ref{sec:file_noblocking}) e la connessione non può essere conclusa
    immediatamente.
  \item[\errcode{ENETUNREACH}] la rete non è raggiungibile.
  \item[\errcode{ETIMEDOUT}] si è avuto timeout durante il tentativo di
    connessione.
  \end{errlist}
  ed inoltre \errval{EADDRINUSE}, \errval{EBADF}, \errval{EFAULT},
  \errval{EINTR}, \errval{EISCONN} e \errval{ENOTSOCK} e nel loro significato
  generico.}
\end{funcproto}

Il primo argomento è un file descriptor ottenuto da una precedente chiamata a
\func{socket}, mentre il secondo e terzo argomento sono rispettivamente
l'indirizzo e la dimensione della struttura che contiene l'indirizzo del
socket, già descritta in sez.~\ref{sec:sock_sockaddr}.

La struttura dell'indirizzo deve essere inizializzata con l'indirizzo IP e il
numero di porta del server a cui ci si vuole connettere usando le funzioni
illustrate in sez.~\ref{sec:sock_addr_func} come mostrato nell'esempio che
vedremo in sez.~\ref{sec:TCP_daytime_client}.

Nel caso di socket TCP la funzione \func{connect} avvia il \textit{three way
  handshake}, e ritorna solo quando la connessione è stabilita o si è
verificato un errore. Le possibili cause di errore sono molteplici (ed i
relativi codici riportati sopra), quelle che però dipendono dalla situazione
della rete e non da errori o problemi nella chiamata della funzione sono le
seguenti:
\begin{enumerate}
\item Il client non riceve risposta al SYN: l'errore restituito è
  \errcode{ETIMEDOUT}. Stevens riporta che BSD invia un primo SYN alla
  chiamata di \func{connect}, un altro dopo 6 secondi, un terzo dopo 24
  secondi, se dopo 75 secondi non ha ricevuto risposta viene ritornato
  l'errore. Linux invece ripete l'emissione del SYN ad intervalli di 30
  secondi per un numero di volte che può essere stabilito dall'utente. Questo
  può essere fatto a livello globale con una opportuna \func{sysctl} (o più
  semplicemente scrivendo il valore voluto in
  \sysctlfile{net/ipv4/tcp\_syn\_retries}, vedi
  sez.~\ref{sec:sock_ipv4_sysctl}) e a livello di singolo socket con l'opzione
  \const{TCP\_SYNCNT} (vedi sez.~\ref{sec:sock_tcp_udp_options}). Il valore
  predefinito per la ripetizione dell'invio è di 5 volte, che comporta un
  timeout dopo circa 180 secondi.

\item Il client riceve come risposta al SYN un RST significa che non c'è
  nessun programma in ascolto per la connessione sulla porta specificata (il
  che vuol dire probabilmente che o si è sbagliato il numero della porta o che
  non è stato avviato il server), questo è un errore fatale e la funzione
  ritorna non appena il RST viene ricevuto riportando un errore
  \errcode{ECONNREFUSED}.
  
  Il flag RST sta per \textit{reset} ed è un segmento inviato direttamente
  dal TCP quando qualcosa non va. Tre condizioni che generano un RST sono:
  quando arriva un SYN per una porta che non ha nessun server in ascolto,
  quando il TCP abortisce una connessione in corso, quando TCP riceve un
  segmento per una connessione che non esiste.
  
\item Il SYN del client provoca l'emissione di un messaggio ICMP di
  destinazione non raggiungibile. In questo caso dato che il messaggio può
  essere dovuto ad una condizione transitoria si ripete l'emissione dei SYN
  come nel caso precedente, fino al timeout, e solo allora si restituisce il
  codice di errore dovuto al messaggio ICMP, che da luogo ad un
  \errcode{ENETUNREACH}.
   
\end{enumerate}

Se si fa riferimento al diagramma degli stati del TCP riportato in
fig.~\ref{fig:TCP_state_diag} la funzione \func{connect} porta un socket dallo
stato \texttt{CLOSED} (lo stato iniziale in cui si trova un socket appena
creato) prima allo stato \texttt{SYN\_SENT} e poi, al ricevimento dell'ACK,
nello stato \texttt{ESTABLISHED}. Se invece la connessione fallisce il socket
non è più utilizzabile e deve essere chiuso.

Si noti infine che con la funzione \func{connect} si è specificato solo
indirizzo e porta del server, quindi solo una metà della \textit{socket pair};
essendo questa funzione usata nei client l'altra metà contenente indirizzo e
porta locale viene lasciata all'assegnazione automatica del kernel, e non è
necessario effettuare una \func{bind}.


\subsection{La funzione \func{listen}}
\label{sec:TCP_func_listen}

La funzione \funcd{listen} serve ad usare un socket in modalità passiva, cioè,
come dice il nome, per metterlo in ascolto di eventuali
connessioni;\footnote{questa funzione può essere usata con socket che
  supportino le connessioni, cioè di tipo \const{SOCK\_STREAM} o
  \const{SOCK\_SEQPACKET}.} in sostanza l'effetto della funzione è di portare
il socket dallo stato \texttt{CLOSED} a quello \texttt{LISTEN}. In genere si
chiama la funzione in un server dopo le chiamate a \func{socket} e \func{bind}
e prima della chiamata ad \func{accept}. Il prototipo della funzione, come
definito dalla pagina di manuale, è:

\begin{funcproto}{
\fhead{sys/socket.h}
\fdecl{int listen(int sockfd, int backlog)}
\fdesc{Pone un socket in attesa di una connessione.} 
}

{La funzione ritorna $0$ in caso di successo e $-1$ per un errore, nel qual
  caso \var{errno} assumerà uno dei valori: 
  \begin{errlist}
  \item[\errcode{EADDRINUSE}] qualche altro socket sta già usando l'indirizzo.
  \item[\errcode{EBADF}] l'argomento \param{sockfd} non è un file descriptor
    valido.
  \item[\errcode{ENOTSOCK}] l'argomento \param{sockfd} non è un socket.
  \item[\errcode{EOPNOTSUPP}] il socket è di un tipo che non supporta questa
    operazione.
  \end{errlist}
}
\end{funcproto}

La funzione pone il socket specificato da \param{sockfd} in modalità passiva e
predispone una coda per le connessioni in arrivo di lunghezza pari a
\param{backlog}. La funzione si può applicare solo a socket di tipo
\const{SOCK\_STREAM} o \const{SOCK\_SEQPACKET}. L'argomento \param{backlog}
indica il numero massimo di connessioni pendenti accettate; se esso viene
ecceduto il client al momento della richiesta della connessione riceverà un
errore di tipo \errcode{ECONNREFUSED}, o se il protocollo, come accade nel
caso del TCP, supporta la ritrasmissione, la richiesta sarà ignorata in modo
che la connessione possa venire ritentata.

Per capire meglio il significato di tutto ciò occorre approfondire la modalità
con cui il kernel tratta le connessioni in arrivo. Per ogni socket in ascolto
infatti vengono mantenute due code:
\begin{enumerate}
\item La coda delle connessioni incomplete (\textit{incomplete connection
    queue}) che contiene una voce per ciascun socket per il quale è arrivato
  un SYN ma il \textit{three way handshake} non si è ancora concluso.  Questi
  socket sono tutti nello stato \texttt{SYN\_RECV}.
\item La coda delle connessioni complete (\textit{complete connection queue})
  che contiene una voce per ciascun socket per il quale il \textit{three way
    handshake} è stato completato ma ancora \func{accept} non è ritornata.
  Questi socket sono tutti nello stato \texttt{ESTABLISHED}.
\end{enumerate}

Lo schema di funzionamento è descritto in fig.~\ref{fig:TCP_listen_backlog}:
quando arriva un segmento SYN da un client il kernel crea una voce nella coda
delle connessioni incomplete e risponde con il segmento SYN$+$ACK. La voce
resterà nella coda delle connessioni incomplete fino al ricevimento del
segmento ACK dal client o fino ad un timeout.

Nel caso di completamento del \textit{three way handshake} la voce viene
spostata nella coda delle connessioni complete.  Quando il processo chiama la
funzione \func{accept} (vedi sez.~\ref{sec:TCP_func_accept}) gli viene passata
la prima voce nella coda delle connessioni complete, oppure, se la coda è
vuota, il processo viene posto in attesa in stato di \textit{sleep} e
risvegliato all'arrivo della prima connessione completa.

\begin{figure}[!htb]
  \centering \includegraphics[width=12cm]{img/tcp_listen_backlog}  
  \caption{Schema di funzionamento delle code delle connessioni complete ed
    incomplete.}
  \label{fig:TCP_listen_backlog}
\end{figure}

Storicamente il valore dell'argomento \param{backlog} era corrispondente al
massimo valore della somma del numero di voci possibili per ciascuna delle due
code. Stevens in \cite{UNP1} riporta che BSD ha sempre applicato un fattore di
1.5 a detto valore, e fornisce una tabella con i risultati ottenuti con vari
kernel, compreso anche il vecchio Linux 2.0, che mostrano le differenze fra
diverse realizzazioni.

In Linux il significato di questo valore è cambiato a partire dal kernel 2.2
per prevenire il \textit{denial of service} chiamato \itindex{SYN~flood}
\textit{SYN flood}. Questo attacco si basa sull'emissione da parte
dell'attaccante di un grande numero di segmenti SYN indirizzati verso una
porta, forgiati con indirizzo IP fasullo (con la tecnica che viene detta
\textit{ip spoofing}); in questo modo i segmenti SYN$+$ACK di risposta vanno
perduti e la coda delle connessioni incomplete viene saturata, impedendo di
fatto ulteriori connessioni.

Per ovviare a questo problema il significato del \param{backlog} è stato
cambiato e adesso indica la lunghezza della coda delle connessioni
complete. La lunghezza della coda delle connessioni incomplete può essere
ancora controllata ma occorre usare esplicitamente la funzione \func{sysctl}
con il parametro \constd{NET\_TCP\_MAX\_SYN\_BACKLOG} o scrivere il valore 
direttamente sul file \sysctlfile{net/ipv4/tcp\_max\_syn\_backlog}.

Quando si attiva la protezione dei \textit{syncookies} però (con l'opzione da
compilare nel kernel e da attivare usando \func{sysctl} o scrivendo nel file
\sysctlfile{net/ipv4/tcp\_syncookies}, vedi sez.~\ref{sec:sock_ipv4_sysctl})
questo valore viene ignorato e non esiste più un valore massimo.  In ogni caso
in Linux il valore di \param{backlog} viene sempre troncato ad un massimo di
\const{SOMAXCONN} se è superiore a detta costante (che di default vale 128);
per i kernel precedenti il 2.4.25 questo valore era fisso e non modificabile,
nelle versioni successive può essere controllato con un parametro di
\func{sysctl}, o scrivendo nel file \sysctlfile{net/core/somaxconn} 
(vedi sez.~\ref{sec:sock_ioctl_IP}).

La scelta storica per il valore assegnato a questo argomento era di 5, e
alcuni vecchi kernel non supportavano neanche valori superiori, ma la
situazione corrente è molto cambiata per via della presenza di server web che
devono gestire un gran numero di connessioni per cui un tale valore non è più
adeguato. Non esiste comunque una risposta univoca per la scelta del valore,
per questo non conviene specificarlo con una costante (il cui cambiamento
richiederebbe la ricompilazione del server) ma usare piuttosto una variabile
di ambiente (vedi sez.~\ref{sec:proc_environ}). 

Stevens tratta accuratamente questo argomento in \cite{UNP1}, con esempi presi
da casistiche reali trovate su dei server web, ed in particolare evidenzia
come non sia più vero che il compito principale della coda sia quello di
gestire il caso in cui il server è occupato fra chiamate successive alla
\func{accept} (per cui la coda più occupata sarebbe quella delle connessioni
completate), ma piuttosto quello di gestire la presenza di un gran numero di
SYN in attesa di concludere il \textit{three way handshake}.

Infine va messo in evidenza che, nel caso di socket TCP, quando un segmento
SYN arriva con tutte le code piene, il pacchetto verrà semplicemente
ignorato. Questo avviene perché la condizione in cui le code sono piene è
ovviamente transitoria, per cui se il client ritrasmette in seguito un
segmento SYN, come previsto dal protocollo, è probabile che essendo passato un
po' di tempo esso possa trovare nella coda lo spazio per una nuova
connessione.

Se al contrario si rispondesse immediatamente con un segmento RST, per
indicare che è impossibile effettuare la connessione, la chiamata a
\func{connect} eseguita dal client fallirebbe ritornando una condizione di
errore. In questo modo si sarebbe costretti ad inserire nell'applicazione la
gestione dei tentativi di riconnessione, che invece grazie a questa modalità
di funzionamento viene effettuata in maniera trasparente dal protocollo
TCP.


\subsection{La funzione \func{accept}}
\label{sec:TCP_func_accept}

La funzione \funcd{accept} è chiamata da un server per gestire la connessione,
nel caso di TCP una volta che sia stato completato il \textit{three way
  handshake};\footnote{come le precedenti, la funzione è generica ed è
  utilizzabile su socket di tipo \const{SOCK\_STREAM}, \const{SOCK\_SEQPACKET}
  e \const{SOCK\_RDM}, ma qui la tratteremo solo per gli aspetti riguardanti
  le connessioni con TCP.} la funzione restituisce un nuovo socket descriptor
su cui si potrà operare per effettuare la comunicazione. Se non ci sono
connessioni completate il processo viene messo in attesa. Il prototipo della
funzione è il seguente:

\begin{funcproto}{
\fhead{sys/socket.h}
\fdecl{int accept(int sockfd, struct sockaddr *addr, socklen\_t *addrlen)} 
\fdesc{Accetta una connessione sul socket specificato.} 
}

{La funzione ritorna un numero di socket descriptor positivo in caso di
  successo e $-1$ per un errore, nel qual caso \var{errno} assumerà uno dei
  valori:
  \begin{errlist}
  \item[\errcode{EAGAIN} o \errcode{EWOULDBLOCK}] il socket è stato impostato
    come non bloccante (vedi sez.~\ref{sec:file_noblocking}), e non ci sono
    connessioni in attesa di essere accettate. In generale possono essere
    restituiti entrambi i valori, per cui se si ha a cuore la portabilità
    occorre controllare entrambi.
  \item[\errcode{EBADF}] l'argomento \param{sockfd} non è un file descriptor
    valido.
  \item[\errcode{ECONNABORTED}] la connessione è stata abortita.
  \item[\errcode{EINTR}] la funzione è stata interrotta da un segnale.
  \item[\errcode{EINVAL}] il socket non è in ascolto o \param{addrlen} non ha
    un valore valido. 
  \item[\errcode{ENOBUFS}, \errcode{ENOMEM}] questo spesso significa che
    l'allocazione della memoria è limitata dai limiti sui buffer dei socket,
    non dalla memoria di sistema.
  \item[\errcode{ENOTSOCK}] l'argomento \param{sockfd} non è un socket.
  \item[\errcode{EOPNOTSUPP}] il socket è di un tipo che non supporta questa
    operazione.
  \item[\errcode{EPERM}] le regole del firewall non consentono la connessione.
  \end{errlist}
  ed inoltre nel loro significato generico: \errval{EFAULT}, \errval{EMFILE},
  \errval{ENFILE}; infine a seconda del protocollo e del kernel possono essere
  restituiti errori di rete relativi al nuovo socket come: \errval{ENOSR},
  \errval{ESOCKTNOSUPPORT}, \errval{EPROTONOSUPPORT}, \errval{ETIMEDOUT},
  \errval{ERESTARTSYS}.}
\end{funcproto}

La funzione estrae la prima connessione relativa al socket \param{sockfd} in
attesa sulla coda delle connessioni complete, che associa ad nuovo socket con
le stesse caratteristiche di \param{sockfd}.  Il socket originale non viene
toccato e resta nello stato di \texttt{LISTEN}, mentre il nuovo socket viene
posto nello stato \texttt{ESTABLISHED}. 

I due argomenti \param{addr} e \param{addrlen} (si noti che quest'ultimo è un
\textit{valure-result argument} passato con un puntatore per riavere indietro
il valore) sono usati rispettivamente per ottenere l'indirizzo del client da
cui proviene la connessione e la lunghezza dello stesso; la dimensione dipende
da quale famiglia di indirizzi si sta utilizzando.  Prima della
chiamata \param{addrlen} deve essere inizializzato alle dimensioni della
struttura degli indirizzi cui punta \param{addr} (un numero positivo); al
ritorno della funzione \param{addrlen} conterrà il numero di byte scritti
dentro \param{addr}. Se questa informazione non interessa basterà
inizializzare a \val{NULL} detti puntatori.

Se la funzione ha successo restituisce il descrittore di un nuovo socket
creato dal kernel (detto \textit{connected socket}) a cui viene associata la
prima connessione completa (estratta dalla relativa coda, vedi
sez.~\ref{sec:TCP_func_listen}) che il client ha effettuato verso il socket
\param{sockfd}. Quest'ultimo (detto \textit{listening socket}) è quello creato
all'inizio e messo in ascolto con \func{listen}, e non viene toccato dalla
funzione.  Se non ci sono connessioni pendenti da accettare la funzione mette
in attesa il processo\footnote{a meno che non si sia impostato il socket per
  essere non bloccante (vedi sez.~\ref{sec:file_noblocking}), nel qual caso
  ritorna con l'errore \errcode{EAGAIN}.  Torneremo su questa modalità di
  operazione in sez.~\ref{sec:TCP_sock_multiplexing}.}  fintanto che non ne
arriva una.

La funzione può essere usata solo con socket che supportino la connessione
(cioè di tipo \const{SOCK\_STREAM}, \const{SOCK\_SEQPACKET} o
\const{SOCK\_RDM}). Per alcuni protocolli che richiedono una conferma
esplicita della connessione,\footnote{attualmente in Linux solo DECnet ha
  questo comportamento.} la funzione opera solo l'estrazione dalla coda delle
connessioni, la conferma della connessione viene eseguita implicitamente dalla
prima chiamata ad una \func{read} o una \func{write}, mentre il rifiuto della
connessione viene eseguito con la funzione \func{close}. 

Si tenga presente che con Linux, seguendo POSIX.1, è sufficiente includere
\texttt{sys/socket.h}, ma alcune implementazioni di altri sistemi possono
richiedere l'inclusione di \texttt{sys/types.h}, per cui dovendo curare la
portabilità può essere il caso di includere anche questo file.  Inoltre Linux
presenta un comportamento diverso nella gestione degli errori rispetto ad
altre realizzazioni dei socket BSD, infatti la funzione \func{accept} passa
gli errori di rete pendenti sul nuovo socket come codici di errore per
\func{accept}, per cui l'applicazione deve tenerne conto ed eventualmente
ripetere la chiamata alla funzione come per l'errore di \errcode{EAGAIN}
(torneremo su questo in sez.~\ref{sec:TCP_echo_critical}).

Un'altra differenza con BSD è che la funzione non fa ereditare al nuovo socket
i flag del socket originale, come \const{O\_NONBLOCK},\footnote{ed in generale
  tutti quelli che si possono impostare con \func{fcntl}, vedi
  sez.~\ref{sec:file_fcntl_ioctl}.} che devono essere rispecificati ogni
volta. Tutto questo deve essere tenuto in conto se si devono scrivere
programmi portabili. Per poter effettuare questa impostazione in maniera
atomica, senza dover ricorrere ad ulteriori chiamate a \func{fcntl} su Linux è
disponibile anche la funzione \funcd{accept4}, il cui prototipo è:\footnote{la
  funzione è utilizzabile solo se si è definito la macro \macro{\_GNU\_SOURCE}
  ed ovviamente non è portabile.}

\begin{funcproto}{
\fhead{sys/socket.h}
\fdecl{int accept4(int sockfd, struct sockaddr *addr, socklen\_t *addrlen, int
  flags)} 
\fdesc{Accetta una connessione sul socket specificato.} 
}

{La funzione ritorna $0$ in caso di successo e $-1$ per un errore, nel qual
  caso \var{errno} assumerà gli stessi valori di \func{accept}.}
\end{funcproto}

La funzione aggiunge un quarto argomento \param{flags} usato come maschera
binaria, e se questo è nullo il suo comportamento è identico a quello di
\func{accept}. Con \param{flags} si possono impostare contestualmente
all'esecuzione sul file descriptor restituito i due flag di
\const{O\_NONBLOCK} e \const{O\_CLOEXEC}, fornendo un valore che sia un OR
aritmetico delle costanti in tab.\ref{tab:accept4_flags_arg}.

\begin{table}[htb]
  \centering
  \footnotesize
  \begin{tabular}[c]{|l|l|}
    \hline
    \textbf{Costante} & \textbf{Significato} \\
    \hline
    \hline
    \constd{SOCK\_NONBLOCK} & imposta sul file descriptor restituito il flag
                              di \const{O\_NONBLOCK}\\
    \constd{SOCK\_NOXEC}    & imposta sul file descriptor restituito il flag
                              di \const{O\_CLOEXEC}\\ 
    \hline    
  \end{tabular}
  \caption{Costanti per i possibili valori dell'argomento \param{flags} di
    \func{accept4}.}
  \label{tab:accept4_flags_arg}
\end{table}

Il meccanismo di funzionamento di \func{accept} è essenziale per capire il
funzionamento di un server: in generale infatti c'è sempre un solo socket in
ascolto, detto per questo \textit{listening socket}, che resta per tutto il
tempo nello stato \texttt{LISTEN}, mentre le connessioni vengono gestite dai
nuovi socket, detti \textit{connected socket}, ritornati da \func{accept}, che
si trovano automaticamente nello stato \texttt{ESTABLISHED}, e vengono
utilizzati per lo scambio dei dati, che avviene su di essi, fino alla chiusura
della connessione.  Si può riconoscere questo schema anche nell'esempio
elementare di fig.~\ref{fig:TCP_daytime_iter_server_code}, dove per ogni
connessione il socket creato da \func{accept} viene chiuso dopo l'invio dei
dati.


\subsection{Le funzioni \func{getsockname} e \func{getpeername}}
\label{sec:TCP_get_names}

Oltre a tutte quelle viste finora, dedicate all'utilizzo dei socket, esistono
alcune funzioni ausiliarie che possono essere usate per recuperare alcune
informazioni relative ai socket ed alle connessioni ad essi associate. Le due
più elementari sono le seguenti, usate per ottenere i dati relativi alla
\textit{socket pair} associata ad un certo socket.  La prima è
\funcd{getsockname} e serve ad ottenere l'indirizzo locale associato ad un
socket; il suo prototipo è:

\begin{funcproto}{
\fhead{sys/socket.h}
\fdecl{int getsockname(int sockfd, struct sockaddr *name, socklen\_t *namelen)}
\fdesc{Legge l'indirizzo locale di un socket.} 
}

{La funzione ritorna $0$ in caso di successo e $-1$ per un errore, nel qual
  caso \var{errno} assumerà uno dei valori: 
  \begin{errlist}
  \item[\errcode{EBADF}] l'argomento \param{sockfd} non è un file descriptor
    valido.
  \item[\errcode{ENOBUFS}] non ci sono risorse sufficienti nel sistema per
  \item[\errcode{ENOTSOCK}] l'argomento \param{sockfd} non è un socket.
    eseguire l'operazione.
  \end{errlist}
  ed \errval{EFAULT} nel suo significato generico.}
\end{funcproto}

La funzione restituisce la struttura degli indirizzi del socket \param{sockfd}
nella struttura indicata dal puntatore \param{name} la cui lunghezza è
specificata tramite l'argomento \param{namlen}. Quest'ultimo è un
\textit{value result argument} e pertanto viene passato come indirizzo per
avere indietro anche il numero di byte effettivamente scritti nella struttura
puntata da \param{name}. Si tenga presente che se si è utilizzato un buffer
troppo piccolo per \param{name} l'indirizzo risulterà troncato.

La funzione si usa tutte le volte che si vuole avere l'indirizzo locale di un
socket; ad esempio può essere usata da un client (che usualmente non chiama
\func{bind}) per ottenere l'indirizzo IP e la porta locale associati al socket
restituito da una \func{connect}, o da un server che ha chiamato \func{bind}
su un socket usando $0$ come porta locale per ottenere il numero di porta
effimera assegnato dal kernel.

Inoltre quando un server esegue una \func{bind} su un indirizzo generico, se
chiamata dopo il completamento di una connessione sul socket restituito da
\func{accept}, restituisce l'indirizzo locale che il kernel ha assegnato a
quella connessione. Invece tutte le volte che si vuole avere l'indirizzo
remoto di un socket si usa la funzione \funcd{getpeername}, il cui prototipo
è:

\begin{funcproto}{
\fhead{sys/socket.h}
\fdecl{int getpeername(int sockfd, struct sockaddr * name, socklen\_t *
  namelen)} 
\fdesc{Legge l'indirizzo remoto di un socket.} 
}

{La funzione ritorna $0$ in caso di successo e $-1$ per un errore, nel qual
  caso \var{errno} assumerà uno dei valori: 
  \begin{errlist}
  \item[\errcode{EBADF}] l'argomento \param{sockfd} non è un file descriptor
    valido.
  \item[\errcode{ENOBUFS}] non ci sono risorse sufficienti nel sistema per
    eseguire l'operazione.
  \item[\errcode{ENOTCONN}] il socket non è connesso.
  \item[\errcode{ENOTSOCK}] l'argomento \param{sockfd} non è un socket.
  \end{errlist}
  ed \errval{EFAULT} nel suo significato generico.}
\end{funcproto}


La funzione è identica a \func{getsockname}, ed usa la stessa sintassi, ma
restituisce l'indirizzo remoto del socket, cioè quello associato all'altro
capo della connessione.  Ci si può chiedere a cosa serva questa funzione dato
che dal lato client l'indirizzo remoto è sempre noto quando si esegue la
\func{connect} mentre dal lato server si possono usare, come vedremo in
fig.~\ref{fig:TCP_daytime_cunc_server_code}, i valori di ritorno di
\func{accept}.

Il fatto è che in generale quest'ultimo caso non è sempre possibile.  In
particolare questo avviene quando il server, invece di gestire la connessione
direttamente in un processo figlio, come vedremo nell'esempio di server
concorrente di sez.~\ref{sec:TCP_daytime_cunc_server}, lancia per ciascuna
connessione un altro programma, usando \func{exec}. Questa ad esempio è la
modalità con cui opera il \textsl{super-server} \cmd{xinetd}, che può gestire
tutta una serie di servizi diversi, eseguendo su ogni connessione ricevuta
sulle porte tenute sotto controllo, il relativo server.

In questo caso benché il processo figlio abbia una immagine della memoria che
è copia di quella del processo padre (e contiene quindi anche la struttura
ritornata da \func{accept}), all'esecuzione di \func{exec} verrà caricata in
memoria l'immagine del programma eseguito, che a questo punto perde ogni
riferimento ai valori tornati da \func{accept}.  Il socket descriptor però
resta aperto, e se si è seguita una opportuna convenzione per rendere noto al
programma eseguito qual è il socket connesso (ad esempio il solito
\cmd{xinetd} fa sempre in modo che i file descriptor 0, 1 e 2 corrispondano al
socket connesso) quest'ultimo potrà usare la funzione \func{getpeername} per
determinare l'indirizzo remoto del client.

Infine è da chiarire (si legga la pagina di manuale) che, come per
\func{accept}, il terzo argomento, che è specificato dallo standard POSIX.1g
come di tipo \code{socklen\_t *} in realtà deve sempre corrispondere ad un
\ctyp{int *} come prima dello standard, perché tutte le realizzazioni dei
socket BSD fanno questa assunzione.


\subsection{La funzione \func{close}}
\label{sec:TCP_func_close}

La funzione standard Unix \func{close} (vedi sez.~\ref{sec:file_open_close})
che si usa sui file può essere usata con lo stesso effetto anche sui file
descriptor associati ad un socket.

L'azione di questa funzione quando applicata a socket è di marcarlo come
chiuso e ritornare immediatamente al processo. Una volta chiamata il socket
descriptor non è più utilizzabile dal processo e non può essere usato come
argomento per una \func{write} o una \func{read} (anche se l'altro capo della
connessione non avesse chiuso la sua parte).  Il kernel invierà comunque tutti
i dati che ha in coda prima di iniziare la sequenza di chiusura.  Vedremo più
avanti in sez.~\ref{sec:sock_generic_options} come sia possibile cambiare
questo comportamento, e cosa può essere fatto perché il processo possa
assicurarsi che l'altro capo abbia ricevuto tutti i dati.

Come per tutti i file descriptor anche per i socket viene mantenuto un numero
di riferimenti, per cui se più di un processo ha lo stesso socket aperto
l'emissione del FIN e la sequenza di chiusura di TCP non viene innescata
fintanto che il numero di riferimenti non si annulla, questo si applica, come
visto in sez.~\ref{sec:file_shared_access}, sia ai file descriptor duplicati
che a quelli ereditati dagli eventuali processi figli, ed è il comportamento
che ci si aspetta in una qualunque applicazione client/server.

Per attivare immediatamente l'emissione del FIN e la sequenza di chiusura
descritta in sez.~\ref{sec:TCP_conn_term}, si può invece usare la funzione
\func{shutdown} su cui torneremo in seguito (vedi
sez.~\ref{sec:TCP_shutdown}).


\section{Un esempio elementare: il servizio \textit{daytime}}
\label{sec:TCP_daytime_application}

Avendo introdotto le funzioni di base per la gestione dei socket, potremo
vedere in questa sezione un primo esempio di applicazione elementare che
realizza il servizio \textit{daytime} su TCP, secondo quanto specificato
dall'\href{http://www.ietf.org/rfc/rfc867.txt}{RFC~867}.  Prima di passare
agli esempi del client e del server, inizieremo riesaminando con maggiori
dettagli una peculiarità delle funzioni di I/O, già accennata in
sez.~\ref{sec:file_read} e sez.~\ref{sec:file_write}, che nel caso dei socket
è particolarmente rilevante.  Passeremo poi ad illustrare gli esempi della
realizzazione, sia dal lato client, che dal lato server, che si è effettuata
sia in forma iterativa che concorrente.


\subsection{Il comportamento delle funzioni di I/O}
\label{sec:sock_io_behav}

Una cosa che si tende a dimenticare quando si ha a che fare con i socket è che
le funzioni di input/output non sempre hanno lo stesso comportamento che
avrebbero con i normali file di dati (in particolare questo accade per i
socket di tipo stream).

Infatti con i socket è comune che funzioni come \func{read} o \func{write}
possano restituire in input o scrivere in output un numero di byte minore di
quello richiesto. Come già accennato in sez.~\ref{sec:file_read} questo è un
comportamento normale per le funzioni di I/O, ma con i normali file di dati il
problema si avverte solo in lettura, quando si incontra la fine del file. In
generale non è così, e con i socket questo è particolarmente evidente.


\begin{figure}[!htbp]
  \footnotesize \centering
  \begin{minipage}[c]{\codesamplewidth}
    \includecodesample{listati/FullRead.c}
  \end{minipage} 
  \normalsize
  \caption{La funzione \func{FullRead}, che legge esattamente \var{count} byte
    da un file descriptor, iterando opportunamente le letture.}
  \label{fig:sock_FullRead_code}
\end{figure}

Quando ci si trova ad affrontare questo comportamento tutto quello che si deve
fare è semplicemente ripetere la lettura (o la scrittura) per la quantità di
byte restanti, tenendo conto che le funzioni si possono bloccare se i dati non
sono disponibili: è lo stesso comportamento che si può avere scrivendo più di
\const{PIPE\_BUF} byte in una \textit{pipe} (si riveda quanto detto in
sez.~\ref{sec:ipc_pipes}).

Per questo motivo, seguendo l'esempio di R. W. Stevens in \cite{UNP1}, si sono
definite due funzioni, \func{FullRead} e \func{FullWrite}, che eseguono
lettura e scrittura tenendo conto di questa caratteristica, ed in grado di
ritornare solo dopo avere letto o scritto esattamente il numero di byte
specificato; il sorgente è riportato rispettivamente in
fig.~\ref{fig:sock_FullRead_code} e fig.~\ref{fig:sock_FullWrite_code} ed è
disponibile fra i sorgenti allegati alla guida nei file \file{FullRead.c} e
\file{FullWrite.c}.

\begin{figure}[!htbp]
  \centering
  \footnotesize \centering
  \begin{minipage}[c]{\codesamplewidth}
    \includecodesample{listati/FullWrite.c}
  \end{minipage} 
  \normalsize
  \caption{La funzione \func{FullWrite}, che scrive esattamente \var{count}
    byte su un file descriptor, iterando opportunamente le scritture.}
  \label{fig:sock_FullWrite_code}
\end{figure}

Come si può notare le due funzioni ripetono la lettura/scrittura in un ciclo
fino all'esaurimento del numero di byte richiesti, in caso di errore viene
controllato se questo è \errcode{EINTR} (cioè un'interruzione della
\textit{system call} dovuta ad un segnale), nel qual caso l'accesso viene
ripetuto, altrimenti l'errore viene ritornato al programma chiamante,
interrompendo il ciclo.

Nel caso della lettura, se il numero di byte letti è zero, significa che si è
arrivati alla fine del file (per i socket questo significa in genere che
l'altro capo è stato chiuso, e quindi non sarà più possibile leggere niente) e
pertanto si ritorna senza aver concluso la lettura di tutti i byte
richiesti. Entrambe le funzioni restituiscono 0 in caso di successo, ed un
valore negativo in caso di errore, \func{FullRead} restituisce il numero di
byte non letti in caso di \textit{end-of-file} prematuro.


\subsection{Il client \textit{daytime}}
\label{sec:TCP_daytime_client}

Il primo esempio di applicazione delle funzioni di base illustrate in
sez.~\ref{sec:TCP_functions} è relativo alla creazione di un client elementare
per il servizio \textit{daytime}, un servizio elementare, definito
nell'\href{http://www.ietf.org/rfc/rfc867.txt}{RFC~867}, che restituisce
l'ora locale della macchina a cui si effettua la richiesta, e che è assegnato
alla porta 13.

In fig.~\ref{fig:TCP_daytime_client_code} è riportata la sezione principale
del codice del nostro client. Il sorgente completo del programma
(\texttt{TCP\_daytime.c}, che comprende il trattamento delle opzioni ed una
funzione per stampare un messaggio di aiuto) è allegato alla guida nella
sezione dei codici sorgente e può essere compilato su una qualunque macchina
GNU/Linux.

\begin{figure}[!htbp]
  \footnotesize \centering
  \begin{minipage}[c]{\codesamplewidth}
    \includecodesample{listati/TCP_daytime.c}
  \end{minipage} 
  \normalsize
  \caption{Esempio di codice di un client elementare per il servizio
    \textit{daytime}.} 
  \label{fig:TCP_daytime_client_code}
\end{figure}

Il programma anzitutto (\texttt{\small 1--5}) include gli header necessari;
dopo la dichiarazione delle variabili (\texttt{\small 9--12}) si è omessa
tutta la parte relativa al trattamento degli argomenti passati dalla linea di
comando (effettuata con le apposite funzioni illustrate in
sez.~\ref{sec:proc_opt_handling}).

Il primo passo (\texttt{\small 14--18}) è creare un socket TCP (quindi di tipo
\const{SOCK\_STREAM} e di famiglia \const{AF\_INET}). La funzione
\func{socket} ritorna il descrittore che viene usato per identificare il
socket in tutte le chiamate successive. Nel caso la chiamata fallisca si
stampa un errore (\texttt{\small 16}) con la funzione \func{perror} e si esce
(\texttt{\small 17}) con un codice di errore.

Il passo seguente (\texttt{\small 19--27}) è quello di costruire un'apposita
struttura \struct{sockaddr\_in} in cui sarà inserito l'indirizzo del server ed
il numero della porta del servizio. Il primo passo (\texttt{\small 20}) è
inizializzare tutto a zero, per poi inserire il tipo di indirizzo
(\texttt{\small 21}) e la porta (\texttt{\small 22}), usando per quest'ultima
la funzione \func{htons} per convertire il formato dell'intero usato dal
computer a quello usato nella rete, infine (\texttt{\small 23--27}) si può
utilizzare la funzione \func{inet\_pton} per convertire l'indirizzo numerico
passato dalla linea di comando.

A questo punto (\texttt{\small 28--32}) usando la funzione \func{connect} sul
socket creato in precedenza (\texttt{\small 29}) si può stabilire la
connessione con il server. Per questo si deve utilizzare come secondo
argomento la struttura preparata in precedenza con il relativo indirizzo; si
noti come, esistendo diversi tipi di socket, si sia dovuto effettuare un cast.
Un valore di ritorno della funzione negativo implica il fallimento della
connessione, nel qual caso si stampa un errore (\texttt{\small 30}) e si
ritorna (\texttt{\small 31}).

Completata con successo la connessione il passo successivo (\texttt{\small
  34--40}) è leggere la data dal socket; il protocollo prevede che il server
invii sempre una stringa alfanumerica, il formato della stringa non è
specificato dallo standard, per cui noi useremo il formato usato dalla
funzione \func{ctime}, seguito dai caratteri di terminazione \verb|\r\n|, cioè
qualcosa del tipo:
\begin{Verbatim}
Wed Apr 4 00:53:00 2001\r\n
\end{Verbatim}
questa viene letta dal socket (\texttt{\small 34}) con la funzione \func{read}
in un buffer temporaneo; la stringa poi deve essere terminata (\texttt{\small
  35}) con il solito carattere nullo per poter essere stampata (\texttt{\small
  36}) sullo \textit{standard output} con l'uso di \func{fputs}.

Come si è già spiegato in sez.~\ref{sec:sock_io_behav} la risposta dal socket
potrà arrivare in un unico pacchetto di 26 byte (come avverrà senz'altro nel
caso in questione) ma potrebbe anche arrivare in 26 pacchetti di un byte.  Per
questo nel caso generale non si può mai assumere che tutti i dati arrivino con
una singola lettura, pertanto quest'ultima deve essere effettuata in un ciclo
in cui si continui a leggere fintanto che la funzione \func{read} non ritorni
uno zero (che significa che l'altro capo ha chiuso la connessione) o un numero
minore di zero (che significa un errore nella connessione).

Si noti come in questo caso la fine dei dati sia specificata dal server che
chiude la connessione (anche questo è quanto richiesto dal protocollo); questa
è una delle tecniche possibili (è quella usata pure dal protocollo HTTP), ma
ce ne possono essere altre, ad esempio FTP marca la conclusione di un blocco
di dati con la sequenza ASCII \verb|\r\n| (carriage return e line feed),
mentre il DNS mette la lunghezza in testa ad ogni blocco che trasmette. Il
punto essenziale è che TCP non provvede nessuna indicazione che permetta di
marcare dei blocchi di dati, per cui se questo è necessario deve provvedere il
programma stesso.

Se abilitiamo il servizio \textit{daytime}\footnote{in genere questo viene
  fornito direttamente dal \textsl{superdemone} \cmd{xinetd}, pertanto basta
  assicurarsi che esso sia abilitato nel relativo file di configurazione.}
possiamo verificare il funzionamento del nostro client, avremo allora:
\begin{Console}
[piccardi@gont sources]$ \textbf{./daytime 127.0.0.1}
Mon Apr 21 20:46:11 2003
\end{Console}
%$
e come si vede tutto funziona regolarmente.


\subsection{Un server \textit{daytime} iterativo}
\label{sec:TCP_daytime_iter_server}

Dopo aver illustrato il client daremo anche un esempio di un server
elementare, che sia anche in grado di rispondere al precedente client. Come
primo esempio realizzeremo un server iterativo, in grado di fornire una sola
risposta alla volta. Il codice del programma è nuovamente mostrato in
fig.~\ref{fig:TCP_daytime_iter_server_code}, il sorgente completo
(\texttt{TCP\_iter\_daytimed.c}) è allegato insieme agli altri file degli
esempi.

\begin{figure}[!htbp]
  \footnotesize \centering
  \begin{minipage}[c]{\codesamplewidth}
    \includecodesample{listati/TCP_iter_daytimed.c}
  \end{minipage} 
  \normalsize
  \caption{Esempio di codice di un semplice server per il servizio daytime.}
  \label{fig:TCP_daytime_iter_server_code}
\end{figure}

Come per il client si includono (\texttt{\small 1--9}) gli header necessari a
cui è aggiunto quello per trattare i tempi, e si definiscono (\texttt{\small
  14--18}) alcune costanti e le variabili necessarie in seguito. Come nel caso
precedente si sono omesse le parti relative al trattamento delle opzioni da
riga di comando.

La creazione del socket (\texttt{\small 20--24}) è analoga al caso precedente,
come pure l'inizializzazione (\texttt{\small 25--29}) della struttura
\struct{sockaddr\_in}.  Anche in questo caso (\texttt{\small 28}) si usa la
porta standard del servizio \textit{daytime}, ma come indirizzo IP si usa
(\texttt{\small 29}) il valore predefinito \const{INET\_ANY}, che corrisponde
all'indirizzo generico.

Si effettua poi (\texttt{\small 30--34}) la chiamata alla funzione \func{bind}
che permette di associare la precedente struttura al socket, in modo che
quest'ultimo possa essere usato per accettare connessioni su una qualunque
delle interfacce di rete locali. In caso di errore si stampa (\texttt{\small
  31}) un messaggio, e si termina (\texttt{\small 32}) immediatamente il
programma.

Il passo successivo (\texttt{\small 35--39}) è quello di mettere ``\textsl{in
  ascolto}'' il socket; questo viene fatto (\texttt{\small 36}) con la
funzione \func{listen} che dice al kernel di accettare connessioni per il
socket che abbiamo creato; la funzione indica inoltre, con il secondo
argomento, il numero massimo di connessioni che il kernel accetterà di mettere
in coda per il suddetto socket. Di nuovo in caso di errore si stampa
(\texttt{\small 37}) un messaggio, e si esce (\texttt{\small 38})
immediatamente.

La chiamata a \func{listen} completa la preparazione del socket per l'ascolto
(che viene chiamato anche \textit{listening descriptor}) a questo punto si può
procedere con il ciclo principale (\texttt{\small 40--53}) che viene eseguito
indefinitamente. Il primo passo (\texttt{\small 42}) è porsi in attesa di
connessioni con la chiamata alla funzione \func{accept}, come in precedenza in
caso di errore si stampa (\texttt{\small 43}) un messaggio, e si esce
(\texttt{\small 44}).

Il processo resterà in stato di \textit{sleep} fin quando non arriva e viene
accettata una connessione da un client; quando questo avviene \func{accept}
ritorna, restituendo un secondo descrittore, che viene chiamato
\textit{connected descriptor}, e che è quello che verrà usato dalla successiva
chiamata alla \func{write} per scrivere la risposta al client.

Il ciclo quindi proseguirà determinando (\texttt{\small 46}) il tempo corrente
con una chiamata a \texttt{time}, con il quale si potrà opportunamente
costruire (\texttt{\small 47}) la stringa con la data da trasmettere
(\texttt{\small 48}) con la chiamata a \func{write}. Completata la
trasmissione il nuovo socket viene chiuso (\texttt{\small 52}).  A questo
punto il ciclo si chiude ricominciando da capo in modo da poter ripetere
l'invio della data in risposta ad una successiva connessione.

È importante notare che questo server è estremamente elementare, infatti, a
parte il fatto di poter essere usato solo con indirizzi IPv4, esso è in grado
di rispondere ad un solo un client alla volta: è cioè, come dicevamo, un
\textsl{server iterativo}. Inoltre è scritto per essere lanciato da linea di
comando, se lo si volesse utilizzare come demone occorrerebbero le opportune
modifiche (come una chiamata a \func{daemon} prima dell'inizio del ciclo
principale) per tener conto di quanto illustrato in
sez.~\ref{sec:sess_daemon}. Si noti anche che non si è inserita nessuna forma
di gestione della terminazione del processo, dato che tutti i file descriptor
vengono chiusi automaticamente alla sua uscita, e che, non generando figli,
non è necessario preoccuparsi di gestire la loro terminazione.


\subsection{Un server \textit{daytime} concorrente}
\label{sec:TCP_daytime_cunc_server}

Il server \texttt{daytime} dell'esempio in
sez.~\ref{sec:TCP_daytime_iter_server} è un tipico esempio di server iterativo,
in cui viene servita una richiesta alla volta; in generale però, specie se il
servizio è più complesso e comporta uno scambio di dati più sostanzioso di
quello in questione, non è opportuno bloccare un server nel servizio di un
client per volta; per questo si ricorre alle capacità di multitasking del
sistema.

Come spiegato in sez.~\ref{sec:proc_fork} una delle modalità più comuni di
funzionamento da parte dei server è quella di usare la funzione \func{fork}
per creare, ad ogni richiesta da parte di un client, un processo figlio che si
incarichi della gestione della comunicazione.  Si è allora riscritto il server
\textit{daytime} dell'esempio precedente in forma concorrente, inserendo anche
una opzione per la stampa degli indirizzi delle connessioni ricevute.

In fig.~\ref{fig:TCP_daytime_cunc_server_code} è mostrato un estratto del
codice, in cui si sono tralasciati il trattamento delle opzioni e le parti
rimaste invariate rispetto al precedente esempio (cioè tutta la parte
riguardante l'apertura passiva del socket). Al solito il sorgente completo del
server, nel file \texttt{TCP\_cunc\_daytimed.c}, è allegato insieme ai
sorgenti degli altri esempi.

\begin{figure}[!htbp]
  \footnotesize \centering
  \begin{minipage}[c]{\codesamplewidth}
    \includecodesample{listati/TCP_cunc_daytimed.c}
  \end{minipage} 
  \normalsize
  \caption{Esempio di codice di un server concorrente elementare per il 
    servizio daytime.}
  \label{fig:TCP_daytime_cunc_server_code}
\end{figure}

Stavolta (\texttt{\small 21--26}) la funzione \func{accept} è chiamata
fornendo una struttura di indirizzi in cui saranno ritornati l'indirizzo IP e
la porta da cui il client effettua la connessione, che in un secondo tempo,
(\texttt{\small 40--44}), se il logging è abilitato, stamperemo sullo
\textit{standard output}.

Quando \func{accept} ritorna il server chiama la funzione \func{fork}
(\texttt{\small 27--31}) per creare il processo figlio che effettuerà
(\texttt{\small 32--46}) tutte le operazioni relative a quella connessione,
mentre il padre proseguirà l'esecuzione del ciclo principale in attesa di
ulteriori connessioni.

Si noti come il figlio operi solo sul socket connesso, chiudendo
immediatamente (\texttt{\small 33}) il socket \var{list\_fd}; mentre il padre
continua ad operare solo sul socket in ascolto chiudendo (\texttt{\small 48})
\var{conn\_fd} al ritorno dalla \func{fork}. Per quanto abbiamo detto in
sez.~\ref{sec:TCP_func_close} nessuna delle due chiamate a \func{close} causa
l'innesco della sequenza di chiusura perché il numero di riferimenti al file
descriptor non si è annullato.

Infatti subito dopo la creazione del socket \var{list\_fd} ha una referenza, e
lo stesso vale per \var{conn\_fd} dopo il ritorno di \func{accept}, ma dopo la
\func{fork} i descrittori vengono duplicati nel padre e nel figlio per cui
entrambi i socket si trovano con due referenze. Questo fa sì che quando il
padre chiude \var{sock\_fd} esso resta con una referenza da parte del figlio,
e sarà definitivamente chiuso solo quando quest'ultimo, dopo aver completato
le sue operazioni, chiamerà (\texttt{\small 45}) la funzione \func{close}.

In realtà per il figlio non sarebbe necessaria nessuna chiamata a
\func{close}, in quanto con la \func{exit} finale (\texttt{\small 45}) tutti i
file descriptor, quindi anche quelli associati ai socket, vengono
automaticamente chiusi.  Tuttavia si è preferito effettuare esplicitamente le
chiusure per avere una maggiore chiarezza del codice, e per evitare eventuali
errori, prevenendo ad esempio un uso involontario del \textit{listening
  descriptor}.

Si noti invece come sia essenziale che il padre chiuda ogni volta il socket
connesso dopo la \func{fork}; se così non fosse nessuno di questi socket
sarebbe effettivamente chiuso dato che alla chiusura da parte del figlio
resterebbe ancora un riferimento nel padre. Si avrebbero così due effetti: il
padre potrebbe esaurire i descrittori disponibili (che sono un numero limitato
per ogni processo) e soprattutto nessuna delle connessioni con i client
verrebbe chiusa.

Come per ogni server iterativo il lavoro di risposta viene eseguito
interamente dal processo figlio. Questo si incarica (\texttt{\small 34}) di
chiamare \func{time} per leggere il tempo corrente, e di stamparlo
(\texttt{\small 35}) sulla stringa contenuta in \var{buffer} con l'uso di
\func{snprintf} e \func{ctime}. Poi la stringa viene scritta (\texttt{\small
  36--39}) sul socket, controllando che non ci siano errori. Anche in questo
caso si è evitato il ricorso a \func{FullWrite} in quanto la stringa è
estremamente breve e verrà senz'altro scritta in un singolo segmento.

Inoltre nel caso sia stato abilitato il \textit{logging} delle connessioni, si
provvede anche (\texttt{\small 40--43}) a stampare sullo \textit{standard
  output} l'indirizzo e la porta da cui il client ha effettuato la
connessione, usando i valori contenuti nelle strutture restituite da
\func{accept}, eseguendo le opportune conversioni con \func{inet\_ntop} e
\func{ntohs}.

Ancora una volta l'esempio è estremamente semplificato, si noti come di nuovo
non si sia gestita né la terminazione del processo né il suo uso come demone,
che tra l'altro sarebbe stato incompatibile con l'uso della opzione di logging
che stampa gli indirizzi delle connessioni sullo standard output. Un altro
aspetto tralasciato è la gestione della terminazione dei processi figli,
torneremo su questo più avanti quando tratteremo alcuni esempi di server più
complessi.


\section{Un esempio più completo: il servizio \textit{echo}}
\label{sec:TCP_echo_application}

L'esempio precedente, basato sul servizio \textit{daytime}, è un esempio molto
elementare, in cui il flusso dei dati va solo nella direzione dal server al
client. In questa sezione esamineremo un esempio di applicazione client/server
un po' più complessa, che usi i socket TCP per una comunicazione in entrambe
le direzioni.

Ci limiteremo a fornire una realizzazione elementare, che usi solo le funzioni
di base viste finora, ma prenderemo in esame, oltre al comportamento in
condizioni normali, anche tutti i possibili scenari particolari (errori,
sconnessione della rete, crash del client o del server durante la connessione)
che possono avere luogo durante l'impiego di un'applicazione di rete, partendo
da una versione primitiva che dovrà essere rimaneggiata di volta in volta per
poter tenere conto di tutte le evenienze che si possono manifestare nella vita
reale di un'applicazione di rete, fino ad arrivare ad una realizzazione
completa.


\subsection{Il servizio \textit{echo}}
\label{sec:TCP_echo}

Nella ricerca di un servizio che potesse fare da esempio per una comunicazione
bidirezionale, si è deciso, seguendo la scelta di Stevens in \cite{UNP1}, di
usare il servizio \textit{echo}, che si limita a restituire in uscita quanto
immesso in ingresso. Infatti, nonostante la sua estrema semplicità, questo
servizio costituisce il prototipo ideale per una generica applicazione di rete
in cui un server risponde alle richieste di un client.  Nel caso di una
applicazione più complessa quello che si potrà avere in più è una elaborazione
dell'input del client, che in molti casi viene interpretato come un comando,
da parte di un server che risponde fornendo altri dati in uscita.

Il servizio \textit{echo} è uno dei servizi standard solitamente provvisti
direttamente dal superserver \cmd{inetd}, ed è definito
dall'\href{http://www.ietf.org/rfc/rfc862.txt}{RFC~862}. Come dice il nome il
servizio deve riscrivere indietro sul socket i dati che gli vengono inviati in
ingresso. L'RFC descrive le specifiche del servizio sia per TCP che UDP, e per
qil primo stabilisce che una volta stabilita la connessione ogni dato in
ingresso deve essere rimandato in uscita fintanto che il chiamante non ha
chiude la connessione. Al servizio è assegnata la porta riservata 7.

Nel nostro caso l'esempio sarà costituito da un client che legge una linea di
caratteri dallo \textit{standard input} e la scrive sul server. A sua volta il
server leggerà la linea dalla connessione e la riscriverà immutata
all'indietro. Sarà compito del client leggere la risposta del server e
stamparla sullo \textit{standard output}.


\subsection{Il client \textit{echo}: prima versione}
\label{sec:TCP_echo_client}

Il codice della prima versione del client per il servizio \textit{echo} (file
\texttt{TCP\_echo\_first.c}) dei sorgenti allegati alla guida) è riportato in
fig.~\ref{fig:TCP_echo_client_1}. Esso ricalca la struttura del precedente
client per il servizio \textit{daytime} (vedi
sez.~\ref{sec:TCP_daytime_client}), e la prima parte (\texttt{\small 10--27})
è sostanzialmente identica, a parte l'uso di una porta diversa.

\begin{figure}[!htb]
  \footnotesize \centering
  \begin{minipage}[c]{\codesamplewidth}
    \includecodesample{listati/TCP_echo_first.c}
  \end{minipage} 
  \normalsize
  \caption{Codice della prima versione del client \textit{echo}.}
  \label{fig:TCP_echo_client_1}
\end{figure}

Al solito si è tralasciata la sezione relativa alla gestione delle opzioni a
riga di comando.  Una volta dichiarate le variabili, si prosegue
(\texttt{\small 10--13}) con la creazione del socket con l'usuale controllo
degli errori, alla preparazione (\texttt{\small 14--17}) della struttura
dell'indirizzo, che stavolta usa la porta 7 riservata al servizio
\textit{echo}, infine si converte (\texttt{\small 18--22}) l'indirizzo
specificato a riga di comando.  A questo punto (\texttt{\small 23--27}) si può
eseguire la connessione al server secondo la stessa modalità usata in
sez.~\ref{sec:TCP_daytime_client}.

Completata la connessione, per gestire il funzionamento del protocollo si usa
la funzione \code{ClientEcho}, il cui codice si è riportato a parte in
fig.~\ref{fig:TCP_client_echo_sub}. Questa si preoccupa di gestire tutta la
comunicazione, leggendo una riga alla volta dallo \textit{standard input}
\file{stdin}, scrivendola sul socket e ristampando su \file{stdout} quanto
ricevuto in risposta dal server. Al ritorno dalla funzione (\texttt{\small
  30--31}) anche il programma termina.

La funzione \code{ClientEcho} utilizza due buffer (\texttt{\small 3}) per
gestire i dati inviati e letti sul socket.  La comunicazione viene gestita
all'interno di un ciclo (\texttt{\small 5--10}), i dati da inviare sulla
connessione vengono presi dallo \file{stdin} usando la funzione \func{fgets},
che legge una linea di testo (terminata da un \texttt{CR} e fino al massimo di
\const{MAXLINE} caratteri) e la salva sul buffer di invio.

Si usa poi (\texttt{\small 6}) la funzione \func{FullWrite}, vista in
sez.~\ref{sec:sock_io_behav}, per scrivere i dati sul socket, gestendo
automaticamente l'invio multiplo qualora una singola \func{write} non sia
sufficiente.  I dati vengono riletti indietro (\texttt{\small 7}) con una
\func{read}\footnote{si è fatta l'assunzione implicita che i dati siano
  contenuti tutti in un solo segmento, così che la chiamata a \func{read} li
  restituisca sempre tutti; avendo scelto una dimensione ridotta per il buffer
  questo sarà sempre vero, vedremo più avanti come superare il problema di
  rileggere indietro tutti e soli i dati disponibili, senza bloccarsi.} sul
buffer di ricezione e viene inserita (\texttt{\small 8}) la terminazione della
stringa e per poter usare (\texttt{\small 9}) la funzione \func{fputs} per
scriverli su \file{stdout}.

\begin{figure}[!htbp]
  \footnotesize \centering
  \begin{minipage}[c]{\codesamplewidth}
    \includecodesample{listati/ClientEcho_first.c}
  \end{minipage} 
  \normalsize
  \caption{Codice della prima versione della funzione \texttt{ClientEcho} per 
    la gestione del servizio \textit{echo}.}
  \label{fig:TCP_client_echo_sub}
\end{figure}

Quando si concluderà l'invio di dati mandando un \textit{end-of-file} sullo
\textit{standard input} si avrà il ritorno di \func{fgets} con un puntatore
nullo (si riveda quanto spiegato in sez.~\ref{sec:file_line_io}) e la
conseguente uscita dal ciclo; al che la subroutine ritorna ed il nostro
programma client termina.

Si può effettuare una verifica del funzionamento del client abilitando il
servizio \textit{echo} nella configurazione di \cmd{xinetd} sulla propria
macchina ed usandolo direttamente verso di esso in locale, vedremo in
dettaglio più avanti (in sez.~\ref{sec:TCP_echo_startup}) il funzionamento del
programma, usato però con la nostra versione del server \textit{echo}, che
illustriamo immediatamente.


\subsection{Il server \textit{echo}: prima versione}
\label{sec:TCPsimp_server_main}

La prima versione del server, contenuta nel file \texttt{TCP\_echod\_first.c},
è riportata in fig.~\ref{fig:TCP_echo_server_first_code} e
fig.~\ref{fig:TCP_echo_server_first_main}. Come abbiamo fatto per il client
anche il server è stato diviso in un corpo principale, costituito dalla
funzione \code{main}, che è molto simile a quello visto nel precedente esempio
per il server del servizio \textit{daytime} di
sez.~\ref{sec:TCP_daytime_cunc_server}, e da una funzione ausiliaria
\code{ServEcho} che si cura della gestione del servizio.

\begin{figure}[!htb]
  \footnotesize \centering
  \begin{minipage}[c]{\codesamplewidth}
    \includecodesample{listati/TCP_echod_first_init.c}
  \end{minipage} 
  \normalsize
  \caption{Codice di inizializzatione della prima versione del server
    per il servizio \textit{echo}.}
  \label{fig:TCP_echo_server_first_code}
\end{figure}

In questo caso però, rispetto a quanto visto nell'esempio di
fig.~\ref{fig:TCP_daytime_cunc_server_code} si è preferito scrivere il server
curando maggiormente alcuni dettagli, per tenere conto anche di alcune
esigenze generali (che non riguardano direttamente la rete), come la
possibilità di lanciare il server anche in modalità interattiva e la cessione
dei privilegi di amministratore non appena questi non sono più necessari.

La sezione iniziale del programma (\texttt{\small 8--21}) è la stessa del
server di sez.~\ref{sec:TCP_daytime_cunc_server}, ed ivi descritta in
dettaglio: crea il socket, inizializza l'indirizzo e esegue \func{bind}; dato
che quest'ultima funzione viene usata su una porta riservata, il server dovrà
essere eseguito da un processo con i privilegi di amministratore, pena il
fallimento della chiamata.

Una volta eseguita la funzione \func{bind} però i privilegi di amministratore
non sono più necessari, per questo è sempre opportuno rilasciarli, in modo da
evitare problemi in caso di eventuali vulnerabilità del server.  Per questo
prima (\texttt{\small 22--26}) si esegue \func{setgid} per assegnare il
processo ad un gruppo senza privilegi,\footnote{si è usato il valore 65534,
  ovvero -1 per il formato \ctyp{short}, che di norma in tutte le
  distribuzioni viene usato per identificare il gruppo \texttt{nogroup} e
  l'utente \texttt{nobody}, usati appunto per eseguire programmi che non
  richiedono nessun privilegio particolare.} e poi si ripete (\texttt{\small
  27--30}) l'operazione usando \func{setuid} per cambiare anche
l'utente.\footnote{si tenga presente che l'ordine in cui si eseguono queste
  due operazioni è importante, infatti solo avendo i privilegi di
  amministratore si può cambiare il gruppo di un processo ad un altro di cui
  non si fa parte, per cui chiamare prima \func{setuid} farebbe fallire una
  successiva chiamata a \func{setgid}.  Inoltre si ricordi (si riveda quanto
  esposto in sez.~\ref{sec:proc_perms}) che usando queste due funzioni il
  rilascio dei privilegi è irreversibile.}  Infine (\texttt{\small 30--36}),
qualora sia impostata la variabile \var{demonize}, prima (\texttt{\small 31})
si apre il sistema di logging per la stampa degli errori, e poi
(\texttt{\small 32--35}) si invoca \func{daemon} per eseguire in background il
processo come demone.


\begin{figure}[!htb]
  \footnotesize \centering
  \begin{minipage}[c]{\codesamplewidth}
    \includecodesample{listati/TCP_echod_first_main.c}
  \end{minipage} 
  \normalsize
  \caption{Codice di inizializzatione della prima versione del server
    per il servizio \textit{echo}.}
  \label{fig:TCP_echo_server_first_main}
\end{figure}

A questo punto il programma prosegue nel ciclo principale, illustrato in
fig.~\ref{fig:TCP_echo_server_first_main}, usando lo schema già visto in
precedenza per server per il servizio \textit{daytime}, con l'unica differenza
della chiamata alla funzione \code{PrintErr}, riportata in
fig.~\ref{fig:TCP_PrintErr}, al posto di \func{perror} per la stampa degli
errori.

Si inizia con il porre (\texttt{\small 3--6}) in ascolto il socket, e poi si
esegue indefinitamente il ciclo principale (\texttt{\small 7--26}).
All'interno di questo si ricevono (\texttt{\small 9--12}) le connessioni,
creando (\texttt{\small 13--16}) un processo figlio per ciascuna di esse.
Quest'ultimo (\texttt{\small 17--21}), chiuso (\texttt{\small 18}) il
\textit{listening socket}, esegue (\texttt{\small 19}) la funzione di gestione
del servizio \code{ServEcho}, ed al ritorno di questa esce (\texttt{\small
  20}).

Il padre invece si limita (\texttt{\small 22}) a chiudere il \textit{connected
  socket} per ricominciare da capo il ciclo in attesa di nuove connessioni. In
questo modo si ha un server concorrente. La terminazione del padre non è
gestita esplicitamente, e deve essere effettuata inviando un segnale al
processo.

Avendo trattato direttamente la gestione del programma come demone, si è
dovuto anche provvedere alla necessità di poter stampare eventuali messaggi di
errore attraverso il sistema del \textit{syslog} trattato in
sez.~\ref{sec:sess_daemon}. Come accennato questo è stato fatto utilizzando
come \textit{wrapper} la funzione \code{PrintErr}, il cui codice è riportato
in fig.~\ref{fig:TCP_PrintErr}.

\begin{figure}[!htb]
  \footnotesize \centering
  \begin{minipage}[c]{\codesamplewidth}
    \includecodesample{listati/PrintErr.c}
  \end{minipage} 
  \normalsize
  \caption{Codice della funzione \code{PrintErr} per la generalizzazione della
    stampa degli errori sullo \textit{standard input} o attraverso il
    \texttt{syslog}.}
  \label{fig:TCP_PrintErr}
\end{figure}


In essa ci si limita a controllare (\texttt{\small 2}) se è stato impostato
(valore attivo per default) l'uso come demone, nel qual caso (\texttt{\small
  3}) si usa \func{syslog} (vedi sez.~\ref{sec:sess_daemon}) per stampare il
messaggio di errore fornito come argomento sui log di sistema. Se invece si è
in modalità interattiva (attivabile con l'opzione \texttt{-i}) si usa
(\texttt{\small 5}) semplicemente la funzione \func{perror} per stampare sullo
\textit{standard error}.

La gestione del servizio \textit{echo} viene effettuata interamente nella
funzione \code{ServEcho}, il cui codice è mostrato in
fig.~\ref{fig:TCP_ServEcho_first}, e la comunicazione viene gestita all'interno
di un ciclo (\texttt{\small 6--13}).  I dati inviati dal client vengono letti
(\texttt{\small 6}) dal socket con una semplice \func{read}, di cui non si
controlla lo stato di uscita, assumendo che ritorni solo in presenza di dati
in arrivo. 

La riscrittura (\texttt{\small 7}) viene invece gestita dalla funzione
\func{FullWrite} (descritta in fig.~\ref{fig:sock_FullWrite_code}) che si
incarica di tenere conto automaticamente della possibilità che non tutti i
dati di cui è richiesta la scrittura vengano trasmessi con una singola
\func{write}.

\begin{figure}[!htb] 
  \footnotesize \centering
  \begin{minipage}[c]{\codesamplewidth}
    \includecodesample{listati/ServEcho_first.c}
  \end{minipage} 
  \normalsize
  \caption{Codice della prima versione della funzione \code{ServEcho} per la
    gestione del servizio \textit{echo}.}
  \label{fig:TCP_ServEcho_first}
\end{figure}

In caso di errore di scrittura (si ricordi che \func{FullWrite} restituisce un
valore nullo in caso di successo) si provvede (\texttt{\small 8--10}) a
stampare il relativo messaggio con \func{PrintErr}.  Quando il client chiude
la connessione il ricevimento del segmento FIN fa ritornare la \func{read} con
un numero di byte letti pari a zero, il che causa l'uscita dal ciclo e il
ritorno (\texttt{\small 12}) della funzione, che a sua volta causa la
terminazione del processo figlio.


\subsection{L'avvio e il funzionamento ordinario}
\label{sec:TCP_echo_startup}

Benché il codice dell'esempio precedente sia molto ridotto, esso ci permetterà
di considerare in dettaglio le varie problematiche che si possono incontrare
nello scrivere un'applicazione di rete. Infatti attraverso l'esame delle sue
modalità di funzionamento ordinarie, all'avvio e alla terminazione, e di
quello che avviene nelle varie situazioni limite, da una parte potremo
approfondire la comprensione del protocollo TCP/IP e dall'altra ricavare le
indicazioni necessarie per essere in grado di scrivere applicazioni robuste,
in grado di gestire anche i casi limite.

Il primo passo è compilare e lanciare il server (usando l'utente di
amministrazione, per poter usare la porta 7 che è riservata). All'avvio esso
eseguirà l'apertura passiva con la sequenza delle chiamate a \func{socket},
\func{bind}, \func{listen} e poi si bloccherà nella \func{accept}. A questo
punto si potrà controllarne lo stato con \cmd{netstat}:
\begin{Console}
[piccardi@roke piccardi]$ \textbf{netstat -at}
Active Internet connections (servers and established)
Proto Recv-Q Send-Q Local Address           Foreign Address         State 
...
tcp        0      0 *:echo                  *:*                     LISTEN
...
\end{Console}
%$
che ci mostra come il socket sia in ascolto sulla porta richiesta, accettando
connessioni da qualunque indirizzo e da qualunque porta e su qualunque
interfaccia locale.

A questo punto si può lanciare il client, esso chiamerà \func{socket} e
\func{connect}; una volta completato il \textit{three way handshake} la
connessione è stabilita; la \func{connect} ritornerà nel client\footnote{si
  noti che è sempre la \func{connect} del client a ritornare per prima, in
  quanto questo avviene alla ricezione del secondo segmento (l'ACK del server)
  del \textit{three way handshake}, la \func{accept} del server ritorna solo
  dopo un altro mezzo RTT quando il terzo segmento (l'ACK del client) viene
  ricevuto.}  e la \func{accept} nel server, ed usando di nuovo \cmd{netstat}
otterremmo che:
\begin{Terminal}
Active Internet connections (servers and established)
Proto Recv-Q Send-Q Local Address           Foreign Address         State
tcp        0      0 *:echo                  *:*                     LISTEN
tcp        0      0 roke:echo               gont:32981              ESTABLISHED
\end{Terminal}
mentre per quanto riguarda l'esecuzione dei programmi avremo che:
\begin{itemize}
\item il client chiama la funzione \code{ClientEcho} che si blocca sulla
  \func{fgets} dato che non si è ancora scritto nulla sul terminale.
\item il server eseguirà una \func{fork} facendo chiamare al processo figlio
  la funzione \code{ServEcho}, quest'ultima si bloccherà sulla \func{read}
  dal socket sul quale ancora non sono presenti dati.
\item il processo padre del server chiamerà di nuovo \func{accept}
  bloccandosi fino all'arrivo di un'altra connessione.
\end{itemize}
e se usiamo il comando \cmd{ps} per esaminare lo stato dei processi otterremo
un risultato del tipo:
\begin{Console}
[piccardi@roke piccardi]$ \textbf{ps ax}
  PID TTY      STAT   TIME COMMAND
 ...  ...      ...    ...  ...
 2356 pts/0    S      0:00 ./echod
 2358 pts/1    S      0:00 ./echo 127.0.0.1
 2359 pts/0    S      0:00 ./echod
\end{Console} 
%$
(dove si sono cancellate le righe inutili) da cui si evidenzia la presenza di
tre processi, tutti in stato di \textit{sleep} (vedi
tab.~\ref{tab:proc_proc_states}).

Se a questo punto si inizia a scrivere qualcosa sul client non sarà trasmesso
niente fintanto che non si preme il tasto di a capo (si ricordi quanto detto
in sez.~\ref{sec:file_line_io} a proposito dell'I/O su terminale).  Solo
allora \func{fgets} ritornerà ed il client scriverà quanto immesso dal
terminale sul socket, per poi passare a rileggere quanto gli viene inviato
all'indietro dal server, che a sua volta sarà inviato sullo \textit{standard
  output}, che nel caso ne provoca l'immediata stampa a video.


\subsection{La conclusione normale}
\label{sec:TCP_echo_conclusion}

Tutto quello che scriveremo sul client sarà rimandato indietro dal server e
ristampato a video fintanto che non concluderemo l'immissione dei dati; una
sessione tipica sarà allora del tipo: 
\begin{Console}
[piccardi@roke sources]$ \textbf{./echo 127.0.0.1}
Questa e` una prova
Questa e` una prova
Ho finito
Ho finito
\end{Console} 
%$
che termineremo inviando un EOF dal terminale (usando la combinazione di tasti
ctrl-D, che non compare a schermo); se eseguiamo un \cmd{netstat} a questo
punto avremo:
\begin{Console}
[piccardi@roke piccardi]$ \textbf{netstat -at} 
tcp        0      0 *:echo                  *:*                     LISTEN
tcp        0      0 localhost:33032         localhost:echo          TIME_WAIT
\end{Console} 
%$
con il client che entra in \texttt{TIME\_WAIT}.

Esaminiamo allora in dettaglio la sequenza di eventi che porta alla
terminazione normale della connessione, che ci servirà poi da riferimento
quando affronteremo il comportamento in caso di conclusioni anomale:

\begin{enumerate}
\item inviando un carattere di EOF da terminale la \func{fgets} ritorna
  restituendo un puntatore nullo che causa l'uscita dal ciclo di \code{while},
  così la funzione \code{ClientEcho} ritorna.
\item al ritorno di \code{ClientEcho} ritorna anche la funzione \code{main}, e
  come parte del processo di terminazione tutti i file descriptor vengono chiusi
  (si ricordi quanto detto in sez.~\ref{sec:proc_term_conclusion}); questo
  causa la chiusura del socket di comunicazione; il client allora invierà un
  FIN al server a cui questo risponderà con un ACK.  A questo punto il client
  verrà a trovarsi nello stato \texttt{FIN\_WAIT\_2} ed il server nello stato
  \texttt{CLOSE\_WAIT} (si riveda quanto spiegato in
  sez.~\ref{sec:TCP_conn_term}).
\item quando il server riceve il FIN la \func{read} del processo figlio che
  gestisce la connessione ritorna restituendo 0 causando così l'uscita dal
  ciclo e il ritorno di \code{ServEcho}, a questo punto il processo figlio
  termina chiamando \func{exit}.
\item all'uscita del figlio tutti i file descriptor vengono chiusi, la
  chiusura del socket connesso fa sì che venga effettuata la sequenza finale
  di chiusura della connessione, viene emesso un FIN dal server che riceverà
  un ACK dal client, a questo punto la connessione è conclusa e il client
  resta nello stato \texttt{TIME\_WAIT}.
\end{enumerate}


\subsection{La gestione dei processi figli}
\label{sec:TCP_child_hand}

Tutto questo riguarda la connessione, c'è però da tenere conto dell'effetto
del procedimento di chiusura del processo figlio nel server (si veda quanto
esaminato in sez.~\ref{sec:proc_termination}). In questo caso avremo l'invio
del segnale \signal{SIGCHLD} al padre, ma dato che non si è installato un
gestore e che l'azione predefinita per questo segnale è quella di essere
ignorato, non avendo predisposto la ricezione dello stato di terminazione,
otterremo che il processo figlio entrerà nello stato di \textit{zombie} (si
riveda quanto illustrato in sez.~\ref{sec:sig_sigchld}), come risulterà
ripetendo il comando \cmd{ps}:
\begin{Terminal}
 2356 pts/0    S      0:00 ./echod
 2359 pts/0    Z      0:00 [echod <defunct>]
\end{Terminal}

Dato che non è il caso di lasciare processi \textit{zombie}, occorrerà
ricevere opportunamente lo stato di terminazione del processo (si veda
sez.~\ref{sec:proc_wait}), cosa che faremo utilizzando \signal{SIGCHLD}
secondo quanto illustrato in sez.~\ref{sec:sig_sigchld}. Una prima modifica al
nostro server è pertanto quella di inserire la gestione della terminazione dei
processi figli attraverso l'uso di un gestore.  Per questo useremo la funzione
\code{Signal} (che abbiamo illustrato in fig.~\ref{fig:sig_Signal_code}), per
installare il gestore che riceve i segnali dei processi figli terminati già
visto in fig.~\ref{fig:sig_sigchld_handl}.  Basterà allora aggiungere il
seguente codice: \includecodesnip{listati/sigchildhand.c}
\noindent
all'esempio illustrato in fig.~\ref{fig:TCP_echo_server_first_code}.

In questo modo però si introduce un altro problema. Si ricordi infatti che,
come spiegato in sez.~\ref{sec:sig_gen_beha}, quando un programma si trova in
stato di \texttt{sleep} durante l'esecuzione di una \textit{system call},
questa viene interrotta alla ricezione di un segnale. Per questo motivo, alla
fine dell'esecuzione del gestore del segnale, se questo ritorna, il programma
riprenderà l'esecuzione ritornando dalla \textit{system call} interrotta con
un errore di \errcode{EINTR}.

Vediamo allora cosa comporta tutto questo nel nostro caso: quando si chiude il
client, il processo figlio che gestisce la connessione terminerà, ed il padre,
per evitare la creazione di \textit{zombie}, riceverà il segnale
\signal{SIGCHLD} eseguendo il relativo gestore. Al ritorno del gestore però
l'esecuzione nel padre ripartirà subito con il ritorno della funzione
\func{accept} (a meno di un caso fortuito in cui il segnale arriva durante
l'esecuzione del programma in risposta ad una connessione) con un errore di
\errcode{EINTR}. Non avendo previsto questa eventualità il programma considera
questo un errore fatale terminando a sua volta con un messaggio del tipo:
\begin{Console}
[root@gont sources]# \textbf{./echod -i}
accept error: Interrupted system call
\end{Console}

Come accennato in sez.~\ref{sec:sig_gen_beha} le conseguenze di questo
comportamento delle \textit{system call} possono essere superate in due modi
diversi, il più semplice è quello di modificare il codice di \func{Signal} per
richiedere il riavvio automatico delle \textit{system call} interrotte secondo
la semantica di BSD, usando l'opzione \const{SA\_RESTART} di \func{sigaction};
rispetto a quanto visto in fig.~\ref{fig:sig_Signal_code}. Definiremo allora
la nuova funzione \func{SignalRestart}\footnote{anche questa è definita,
  insieme alle altre funzioni riguardanti la gestione dei segnali, nel file
  \file{SigHand.c}, il cui contento completo può essere trovato negli esempi
  allegati.} come mostrato in fig.~\ref{fig:sig_SignalRestart_code}, ed
installeremo il gestore usando quest'ultima.

\begin{figure}[!htbp]
  \footnotesize \centering
  \begin{minipage}[c]{\codesamplewidth}
    \includecodesample{listati/SignalRestart.c}
  \end{minipage}  
  \normalsize 
  \caption{La funzione \func{SignalRestart}, che installa un gestore di
    segnali in semantica BSD per il riavvio automatico delle \textit{system
      call} interrotte.}
  \label{fig:sig_SignalRestart_code}
\end{figure}

Come si può notare questa funzione è identica alla precedente \func{Signal},
illustrata in fig.~\ref{fig:sig_Signal_code}, solo che in questo caso invece
di inizializzare a zero il campo \var{sa\_flags} di \struct{sigaction}, lo si
inizializza (\texttt{\small 5}) al valore \const{SA\_RESTART}. Usando questa
funzione al posto di \func{Signal} nel server non è necessaria nessuna altra
modifica: le \textit{system call} interrotte saranno automaticamente
riavviate, e l'errore \errcode{EINTR} non si manifesterà più.

La seconda soluzione è più invasiva e richiede di controllare tutte le volte
l'errore restituito dalle varie \textit{system call}, ripetendo la chiamata
qualora questo corrisponda ad \errcode{EINTR}. Questa soluzione ha però il
pregio della portabilità, infatti lo standard POSIX dice che la funzionalità
di riavvio automatico delle \textit{system call}, fornita da
\const{SA\_RESTART}, è opzionale, per cui non è detto che essa sia disponibile
su qualunque sistema.  Inoltre in certi casi,\footnote{Stevens in \cite{UNP1}
  accenna che la maggior parte degli Unix derivati da BSD non fanno ripartire
  \func{select}; altri non riavviano neanche \func{accept} e \func{recvfrom},
  cosa che invece nel caso di Linux viene sempre fatta.} anche quando questa è
presente, non è detto possa essere usata con \func{accept}.

La portabilità nella gestione dei segnali però viene al costo di una
riscrittura parziale del server, la nuova versione di questo, in cui si sono
introdotte una serie di nuove opzioni che ci saranno utili per il debug, è
mostrata in fig.~\ref{fig:TCP_echo_server_code_second}, dove si sono riportate
la sezioni di codice modificate nella seconda versione del programma, il
codice completo di quest'ultimo si trova nel file
\texttt{TCP\_echod\_second.c} dei sorgenti allegati alla guida.

La prima modifica effettuata è stata quella di introdurre una nuova opzione a
riga di comando, \texttt{-c}, che permette di richiedere il comportamento
compatibile nella gestione di \signal{SIGCHLD} al posto della semantica BSD
impostando la variabile \var{compat} ad un valore non nullo. Questa è
preimpostata al valore nullo, cosicché se non si usa questa opzione il
comportamento di default del server è di usare la semantica BSD. 

Una seconda opzione aggiunta è quella di inserire un tempo di attesa fisso
specificato in secondi fra il ritorno della funzione \func{listen} e la
chiamata di \func{accept}, specificabile con l'opzione \texttt{-w}, che
permette di impostare la variabile \var{waiting}.  Infine si è introdotta una
opzione \texttt{-d} per abilitare il debugging che imposta ad un valore non
nullo la variabile \var{debugging}. Al solito si è omessa da
fig.~\ref{fig:TCP_echo_server_code_second} la sezione di codice relativa alla
gestione di tutte queste opzioni, che può essere trovata nel sorgente del
programma.

\begin{figure}[!htb ]
  \footnotesize \centering
  \begin{minipage}[c]{\codesamplewidth}
    \includecodesample{listati/TCP_echod_second.c}
  \end{minipage} 
  \normalsize
  \caption{La sezione nel codice della seconda versione del server
    per il servizio \textit{echo} modificata per tener conto dell'interruzione
    delle \textit{system call}.}
  \label{fig:TCP_echo_server_code_second}
\end{figure}

Vediamo allora come è cambiato il nostro server; una volta definite le
variabili e trattate le opzioni il primo passo (\texttt{\small 9--13}) è
verificare la semantica scelta per la gestione di \signal{SIGCHLD}, a seconda
del valore di \var{compat} (\texttt{\small 9}) si installa il gestore con la
funzione \func{Signal} (\texttt{\small 10}) o con \texttt{SignalRestart}
(\texttt{\small 12}), essendo quest'ultimo il valore di default.

Tutta la sezione seguente, che crea il socket, cede i privilegi di
amministratore ed eventualmente lancia il programma come demone, è rimasta
invariata e pertanto è stata omessa in
fig.~\ref{fig:TCP_echo_server_code_second}; l'unica modifica effettuata prima
dell'entrata nel ciclo principale è stata quella di aver introdotto, subito
dopo la chiamata (\texttt{\small 17--20}) alla funzione \func{listen}, una
eventuale pausa con una condizione (\texttt{\small 21}) sulla variabile
\var{waiting}, che viene inizializzata, con l'opzione \texttt{-w Nsec}, al
numero di secondi da aspettare (il valore preimpostato è nullo).

Si è potuto lasciare inalterata tutta la sezione di creazione del socket
perché nel server l'unica chiamata ad una \textit{system call} lenta, che può
essere interrotta dall'arrivo di \signal{SIGCHLD}, è quella ad \func{accept},
che è l'unica funzione che può mettere il processo padre in stato di sleep nel
periodo in cui un figlio può terminare; si noti infatti come le altre
\textit{system call} lente (si ricordi la distinzione fatta in
sez.~\ref{sec:sig_gen_beha}) o sono chiamate prima di entrare nel ciclo
principale, quando ancora non esistono processi figli, o sono chiamate dai
figli stessi e non risentono di \signal{SIGCHLD}.

Per questo l'unica modifica sostanziale nel ciclo principale (\texttt{\small
  23--42}), rispetto precedente versione di fig.~\ref{fig:TCP_ServEcho_first},
è nella sezione (\texttt{\small 25--31}) in cui si effettua la chiamata di
\func{accept}.  Quest'ultima viene effettuata (\texttt{\small 26--27})
all'interno di un ciclo di \code{while}\footnote{la sintassi del C relativa a
  questo ciclo può non essere del tutto chiara. In questo caso infatti si è
  usato un ciclo vuoto che non esegue nessuna istruzione, in questo modo
  quello che viene ripetuto con il ciclo è soltanto il codice che esprime la
  condizione all'interno del \code{while}.}  che la ripete indefinitamente
qualora in caso di errore il valore di \var{errno} sia \errcode{EINTR}. Negli
altri casi si esce in caso di errore effettivo (\texttt{\small 27--29}),
altrimenti il programma prosegue.

\begin{figure}[!htb]
  \footnotesize \centering
  \begin{minipage}[c]{\codesamplewidth}
    \includecodesample{listati/ServEcho_second.c}
  \end{minipage} 
  \normalsize
  \caption{Codice della seconda versione della funzione \code{ServEcho} per la
    gestione del servizio \textit{echo}.}
  \label{fig:TCP_ServEcho_second}
\end{figure}

Si noti che in questa nuova versione si è aggiunta una ulteriore sezione
(\texttt{\small 32--40}) di aiuto per il debug del programma, che eseguita con
un controllo (\texttt{\small 33}) sul valore della variabile \var{debugging}
impostato dall'opzione \texttt{-d}. Qualora questo sia nullo, come
preimpostato, non accade nulla, altrimenti (\texttt{\small 33}) l'indirizzo
ricevuto da \var{accept} viene convertito in una stringa che poi
(\texttt{\small 34--39}) viene opportunamente stampata o sullo schermo o nei
log.

Infine come ulteriore miglioria si è perfezionata la funzione \code{ServEcho},
sia per tenere conto della nuova funzionalità di debugging, che per effettuare
un controllo in caso di errore; il codice della nuova versione è mostrato in
fig.~\ref{fig:TCP_ServEcho_second}.

 Rispetto alla precedente versione di fig.~\ref{fig:TCP_ServEcho_first} in
questo caso si è provveduto a controllare (\texttt{\small 7--10}) il valore di
ritorno di \func{read} per rilevare un eventuale errore, in modo da stampare
(\texttt{\small 8}) un messaggio di errore e ritornare (\texttt{\small 9})
concludendo la connessione.

Inoltre qualora sia stata attivata la funzionalità di debug (avvalorando
\var{debugging} tramite l'apposita opzione \texttt{-d}) si provvederà a
stampare (\texttt{\small 16--24}) il numero di byte e la stringa letta dal
client, tenendo conto della modalità di invocazione del server, se interattiva
o in forma di demone.


\section{I vari scenari critici}
\label{sec:TCP_echo_critical}

Con le modifiche viste in sez.~\ref{sec:TCP_child_hand} il nostro esempio
diventa in grado di affrontare la gestione ordinaria delle connessioni, ma un
server di rete deve tenere conto che, al contrario di quanto avviene per i
server che operano nei confronti di processi presenti sulla stessa macchina,
la rete è di sua natura inaffidabile, per cui è necessario essere in grado di
gestire tutta una serie di situazioni critiche che non esistono per i processi
locali.


\subsection{La terminazione precoce della connessione}
\label{sec:TCP_conn_early_abort}

La prima situazione critica è quella della terminazione precoce, causata da un
qualche errore sulla rete, della connessione effettuata da un client. Come
accennato in sez.~\ref{sec:TCP_func_accept} la funzione \func{accept} riporta
tutti gli eventuali errori di rete pendenti su una connessione sul
\textit{connected socket}. Di norma questo non è un problema, in quanto non
appena completata la connessione, \func{accept} ritorna, e l'errore sarà
rilevato in seguito dal processo che gestisce la connessione, alla prima
chiamata di una funzione che opera sul socket.

È però possibile, dal punto di vista teorico, incorrere anche in uno scenario
del tipo di quello mostrato in fig.~\ref{fig:TCP_early_abort}, in cui la
connessione viene abortita sul lato client per un qualche errore di rete con
l'invio di un segmento RST, prima che nel server sia stata chiamata la
funzione \func{accept}.

\begin{figure}[!htb]
  \centering \includegraphics[width=10cm]{img/tcp_client_early_abort}  
  \caption{Un possibile caso di terminazione precoce della connessione.}
  \label{fig:TCP_early_abort}
\end{figure}

Benché questo non sia un fatto comune, un evento simile può essere osservato
con dei server molto occupati. In tal caso, con una struttura del server
simile a quella del nostro esempio, in cui la gestione delle singole
connessioni è demandata a processi figli, può accadere che il \textit{three
  way handshake} venga completato e la relativa connessione abortita subito
dopo, prima che il padre, per via del carico della macchina, abbia fatto in
tempo ad eseguire la chiamata ad \func{accept}. Di nuovo si ha una situazione
analoga a quella illustrata in fig.~\ref{fig:TCP_early_abort}, in cui la
connessione viene stabilita, ma subito dopo si ha una condizione di errore che
la chiude prima che essa sia stata accettata dal programma.

Questo significa che oltre alla interruzione da parte di un segnale, che
abbiamo trattato in sez.~\ref{sec:TCP_child_hand} nel caso particolare di
\signal{SIGCHLD}, si possono ricevere altri errori non fatali all'uscita di
\func{accept}; questi, come nel caso precedente, necessitano semplicemente la
ripetizione della chiamata senza uscire dal programma. In questo caso anche la
versione modificata del nostro server non sarebbe adatta, in quanto uno di
questi errori causerebbe la terminazione dello stesso. In Linux i possibili
errori di rete non fatali, riportati sul socket connesso al ritorno di
\func{accept}, sono \errcode{ENETDOWN}, \errcode{EPROTO},
\errcode{ENOPROTOOPT}, \errcode{EHOSTDOWN}, \errcode{ENONET},
\errcode{EHOSTUNREACH}, \errcode{EOPNOTSUPP} e \errcode{ENETUNREACH}.

Si tenga presente che questo tipo di terminazione non è riproducibile
terminando il client prima della chiamata ad \func{accept}, come si potrebbe
fare usando l'opzione \texttt{-w} per introdurre una pausa dopo il lancio del
demone, in modo da poter avere il tempo per lanciare e terminare una
connessione usando il programma client. In tal caso infatti, alla terminazione
del client, il socket associato alla connessione viene semplicemente chiuso,
attraverso la sequenza vista in sez.~\ref{sec:TCP_conn_term}, per cui la
\func{accept} ritornerà senza errori, e si avrà semplicemente un
\textit{end-of-file} al primo accesso al socket. Nel caso di Linux inoltre,
anche qualora si modifichi il client per fargli gestire l'invio di un segmento
RST alla chiusura dal socket (usando l'opzione \const{SO\_LINGER}, vedi
sez.~\ref{sec:sock_options_main}), non si ha comunque nessun errore al ritorno
di \func{accept}, quanto un errore di \errcode{ECONNRESET} al primo tentativo
di accesso al socket.



\subsection{La terminazione precoce del server}
\label{sec:TCP_server_crash}

Un secondo caso critico è quello in cui si ha una terminazione precoce del
server, ad esempio perché il programma ha un crash. In tal caso si suppone che
il processo termini per un errore fatale, cosa che potremo simulare
inviandogli un segnale di terminazione. La conclusione del processo comporta
la chiusura di tutti i file descriptor aperti, compresi tutti i socket
relativi a connessioni stabilite; questo significa che al momento del crollo
del servizio il client riceverà un FIN dal server in corrispondenza della
chiusura del socket.

Vediamo allora cosa succede nel nostro caso, facciamo partire una connessione
con il server e scriviamo una prima riga, poi terminiamo il server con un
\texttt{C-c}. A questo punto scriviamo una seconda riga e poi un'altra riga
ancora. Il risultato finale della sessione è il seguente:
\begin{Console}
[piccardi@gont sources]$ \textbf{./echo 192.168.1.141}
Prima riga
Prima riga
Seconda riga dopo il C-c
Altra riga
[piccardi@gont sources]$
\end{Console}

Come si vede il nostro client, nonostante la connessione sia stata interrotta
prima dell'invio della seconda riga, non solo accetta di inviarla, ma prende
anche un'altra riga prima di terminare senza riportare nessun
errore. 

Per capire meglio cosa è successo conviene analizzare il flusso dei pacchetti
utilizzando un analizzatore di traffico come \cmd{tcpdump}. Il comando
permette di selezionare, nel traffico di rete generato su una macchina, i
pacchetti che interessano, stampando a video (o salvando su disco) il loro
contenuto. Non staremo qui ad entrare nei dettagli dell'uso del programma, che
sono spiegati dalla pagina di manuale;\footnote{per una trattazione di base si
  può consultare sez.~5.2.2 di \cite{SGL}.} per l'uso che vogliamo farne
quello che ci interessa è, posizionandosi sulla macchina che fa da client,
selezionare tutti i pacchetti che sono diretti o provengono dalla macchina che
fa da server. In questo modo (posto che non ci siano altre connessioni col
server, cosa che avremo cura di evitare) tutti i pacchetti rilevati
apparterranno alla nostra sessione di interrogazione del servizio.

Il comando \cmd{tcpdump} permette selezioni molto complesse, basate sulle
interfacce su cui passano i pacchetti, sugli indirizzi IP, sulle porte, sulle
caratteristiche ed il contenuto dei pacchetti stessi, inoltre permette di
combinare fra loro diversi criteri di selezione con degli operatori logici;
quando un pacchetto che corrisponde ai criteri di selezione scelti viene
rilevato i suoi dati vengono stampati sullo schermo (anche questi secondo un
formato configurabile in maniera molto precisa).  

Lanciando il comando prima di ripetere la sessione di lavoro mostrata
nell'esempio precedente potremo allora catturare tutti pacchetti scambiati fra
il client ed il server; i risultati (in realtà si è ridotta la lunghezza
dell'output rispetto al reale tagliando alcuni dati non necessari alla
comprensione del flusso) prodotti in questa occasione da \cmd{tcpdump} sono
allora i seguenti:
\begin{Console}
[root@gont gapil]# \textbf{tcpdump src 192.168.1.141 or dst 192.168.1.141 -N -t}
tcpdump: listening on eth0
gont.34559 > anarres.echo: S 800922320:800922320(0) win 5840 
anarres.echo > gont.34559: S 511689719:511689719(0) ack 800922321 win 5792 
gont.34559 > anarres.echo: . ack 1 win 5840 
gont.34559 > anarres.echo: P 1:12(11) ack 1 win 5840 
anarres.echo > gont.34559: . ack 12 win 5792 
anarres.echo > gont.34559: P 1:12(11) ack 12 win 5792 
gont.34559 > anarres.echo: . ack 12 win 5840 
anarres.echo > gont.34559: F 12:12(0) ack 12 win 5792 
gont.34559 > anarres.echo: . ack 13 win 5840 
gont.34559 > anarres.echo: P 12:37(25) ack 13 win 5840 
anarres.echo > gont.34559: R 511689732:511689732(0) win 0 
\end{Console}

Le prime tre righe vengono prodotte al momento in cui lanciamo il nostro
client, e corrispondono ai tre pacchetti del \textit{three way handshake}.
L'output del comando riporta anche i numeri di sequenza iniziali, mentre la
lettera \texttt{S} indica che per quel pacchetto si aveva il SYN flag attivo.
Si noti come a partire dal secondo pacchetto sia sempre attivo il campo
\texttt{ack}, seguito dal numero di sequenza per il quale si da il ricevuto;
quest'ultimo, a partire dal terzo pacchetto, viene espresso in forma relativa
per maggiore compattezza.  Il campo \texttt{win} in ogni riga indica la
\textit{advertised window} di cui parlavamo in sez.~\ref{sec:TCP_TCP_opt}.

Allora si può verificare dall'output del comando come venga appunto realizzata
la sequenza di pacchetti descritta in sez.~\ref{sec:TCP_conn_cre}: prima viene
inviato dal client un primo pacchetto con il SYN che inizia la connessione, a
cui il server risponde dando il ricevuto con un secondo pacchetto, che a sua
volta porta un SYN, cui il client risponde con un il terzo pacchetto di
ricevuto.

Ritorniamo allora alla nostra sessione con il servizio echo: dopo le tre righe
del \textit{three way handshake} non avremo nulla fin tanto che non scriveremo
una prima riga sul client; al momento in cui facciamo questo si genera una
sequenza di altri quattro pacchetti. Il primo, dal client al server,
contraddistinto da una lettera \texttt{P} che significa che il flag PSH è
impostato, contiene la nostra riga (che è appunto di 11 caratteri), e ad esso
il server risponde immediatamente con un pacchetto vuoto di ricevuto.

Poi tocca al server riscrivere indietro quanto gli è stato inviato, per cui
sarà lui a mandare indietro un terzo pacchetto con lo stesso contenuto appena
ricevuto, e a sua volta riceverà dal client un ACK nel quarto pacchetto.
Questo causerà la ricezione dell'eco nel client che lo stamperà a video.

A questo punto noi procediamo ad interrompere l'esecuzione del server con un
\texttt{C-c} (cioè con l'invio di \signal{SIGTERM}): nel momento in cui
facciamo questo vengono immediatamente generati altri due pacchetti. La
terminazione del processo infatti comporta la chiusura di tutti i suoi file
descriptor, il che comporta, per il socket che avevamo aperto, l'inizio della
sequenza di chiusura illustrata in sez.~\ref{sec:TCP_conn_term}.  Questo
significa che dal server partirà un FIN, che è appunto il primo dei due
pacchetti, contraddistinto dalla lettera \texttt{F}, cui seguirà al solito un
ACK da parte del client.  

A questo punto la connessione dalla parte del server è chiusa, ed infatti se
usiamo \cmd{netstat} per controllarne lo stato otterremo che sul server si ha:
\begin{Console}
anarres:/home/piccardi# \textbf{netstat -ant}
Active Internet connections (servers and established)
Proto Recv-Q Send-Q Local Address           Foreign Address         State
...      ...    ... ...                     ...                     ...
tcp        0      0 192.168.1.141:7         192.168.1.2:34626       FIN_WAIT2
\end{Console}
cioè essa è andata nello stato \texttt{FIN\_WAIT2}, che indica l'avvenuta
emissione del segmento FIN, mentre sul client otterremo che essa è andata
nello stato \texttt{CLOSE\_WAIT}:
\begin{Console}
[root@gont gapil]# \textbf{netstat -ant}
Active Internet connections (servers and established)
Proto Recv-Q Send-Q Local Address           Foreign Address         State
...      ...    ... ...                     ...                     ...
tcp        1      0 192.168.1.2:34582       192.168.1.141:7         CLOSE_WAIT 
\end{Console}

Il problema è che in questo momento il client è bloccato dentro la funzione
\texttt{ClientEcho} nella chiamata a \func{fgets}, e sta attendendo dell'input
dal terminale, per cui non è in grado di accorgersi di nulla. Solo quando
inseriremo la seconda riga il comando uscirà da \func{fgets} e proverà a
scriverla sul socket. Questo comporta la generazione degli ultimi due
pacchetti riportati da \cmd{tcpdump}: il primo, inviato dal client contenente
i 25 caratteri della riga appena letta, e ad esso la macchina server
risponderà, non essendoci più niente in ascolto sulla porta 7, con un segmento
di RST, contraddistinto dalla lettera \texttt{R}, che causa la conclusione
definitiva della connessione anche nel client, dove non comparirà più
nell'output di \cmd{netstat}.

Come abbiamo accennato in sez.~\ref{sec:TCP_conn_term} e come vedremo più
avanti in sez.~\ref{sec:TCP_shutdown} la chiusura di un solo capo di un socket
è una operazione lecita, per cui la nostra scrittura avrà comunque successo
(come si può constatare lanciando usando \cmd{strace}\footnote{il comando
  \cmd{strace} è un comando di debug molto utile che prende come argomento un
  altro comando e ne stampa a video tutte le invocazioni di una \textit{system
    call}, coi relativi argomenti e valori di ritorno, per cui usandolo in
  questo contesto potremo verificare che effettivamente la \func{write} ha
  scritto la riga, che in effetti è stata pure trasmessa via rete.}), in
quanto il nostro programma non ha a questo punto alcun modo di sapere che
dall'altra parte non c'è più nessuno processo in grado di leggere quanto
scriverà. Questo sarà chiaro solo dopo il tentativo di scrittura, e la
ricezione del segmento RST di risposta che indica che dall'altra parte non si
è semplicemente chiuso un capo del socket, ma è completamente terminato il
programma.

Per questo motivo il nostro client proseguirà leggendo dal socket, e dato che
questo è stato chiuso avremo che, come spiegato in
sez.~\ref{sec:TCP_conn_term}, la funzione \func{read} ritorna normalmente con
un valore nullo. Questo comporta che la seguente chiamata a \func{fputs} non
ha effetto (viene stampata una stringa nulla) ed il client si blocca di nuovo
nella successiva chiamata a \func{fgets}. Per questo diventa possibile
inserire una terza riga e solo dopo averlo fatto si avrà la terminazione del
programma.

Per capire come questa avvenga comunque, non avendo inserito nel codice nessun
controllo di errore, occorre ricordare che, a parte la bidirezionalità del
flusso dei dati, dal punto di vista del funzionamento nei confronti delle
funzioni di lettura e scrittura, i socket sono del tutto analoghi a delle
\textit{pipe}. Allora, da quanto illustrato in sez.~\ref{sec:ipc_pipes},
sappiamo che tutte le volte che si cerca di scrivere su una \textit{pipe} il
cui altro capo non è aperto il lettura il processo riceve un segnale di
\signal{SIGPIPE}, e questo è esattamente quello che avviene in questo caso, e
siccome non abbiamo un gestore per questo segnale, viene eseguita l'azione
preimpostata, che è quella di terminare il processo.

Per gestire in maniera più corretta questo tipo di evento dovremo allora
modificare il nostro client perché sia in grado di trattare le varie tipologie
di errore, per questo dovremo riscrivere la funzione \func{ClientEcho}, in
modo da controllare gli stati di uscita delle varie chiamate. Si è riportata
la nuova versione della funzione in fig.~\ref{fig:TCP_ClientEcho_second}.

\begin{figure}[!htb]
  \footnotesize \centering
  \begin{minipage}[c]{\codesamplewidth}
    \includecodesample{listati/ClientEcho_second.c}
  \end{minipage} 
  \normalsize
  \caption{La sezione nel codice della seconda versione della funzione
    \func{ClientEcho} usata dal client per il servizio \textit{echo}
    modificata per tener conto degli eventuali errori.}
  \label{fig:TCP_ClientEcho_second}
\end{figure}

Come si può vedere in questo caso si controlla il valore di ritorno di tutte
le funzioni, ed inoltre si verifica la presenza di un eventuale
\textit{end-of-file} in caso di lettura. Con questa modifica il nostro client
\cmd{echo} diventa in grado di accorgersi della chiusura del socket da parte
del server, per cui ripetendo la sequenza di operazioni precedenti stavolta
otterremo che:
\begin{Console}
[piccardi@gont sources]$ \textbf{./echo 192.168.1.141}
Prima riga
Prima riga
Seconda riga dopo il C-c
EOF sul socket
\end{Console}
%$
ma di nuovo si tenga presente che non c'è modo di accorgersi della chiusura
del socket fin quando non si esegue la scrittura della seconda riga; il
protocollo infatti prevede che ci debba essere una scrittura prima di ricevere
un RST che confermi la chiusura del file, e solo alle successive scritture si
potrà ottenere un errore.

Questa caratteristica dei socket ci mette di fronte ad un altro problema
relativo al nostro client, e che cioè esso non è in grado di accorgersi di
nulla fintanto che è bloccato nella lettura del terminale fatta con
\func{gets}. In questo caso il problema è minimo, ma esso riemergerà più
avanti, ed è quello che si deve affrontare tutte le volte quando si ha a che
fare con la necessità di lavorare con più descrittori, nel qual caso si pone
la questione di come fare a non restare bloccati su un socket quando altri
potrebbero essere liberi. Vedremo come affrontare questa problematica in
sez.~\ref{sec:TCP_sock_multiplexing}.
 

\subsection{Altri scenari di terminazione della connessione}
\label{sec:TCP_conn_crash}

La terminazione del server è solo uno dei possibili scenari di terminazione
della connessione, un altro caso è ad esempio quello in cui si ha
un'interruzione sulla rete, cosa che potremo simulare facilmente staccando il
cavo di rete.  Un'altra condizione è quella di un blocco completo della
macchina su cui gira il server che deve essere riavviata, cosa che potremo
simulare sia eseguendo un reset fisico (un normale shutdown non va bene; in
tal caso infatti il sistema provvede a terminare tutti i processi, per cui la
situazione sarebbe sostanzialmente identica alla precedente) oppure, in
maniera più gentile, riavviando la macchina dopo aver interrotto la
connessione di rete.

Cominciamo ad analizzare il primo caso, l'interruzione del collegamento di
rete. Ripetiamo la nostra sessione di lavoro precedente, lanciamo il client,
scriviamo una prima riga, poi stacchiamo il cavo e scriviamo una seconda
riga. Il risultato che otterremo è:
\begin{Console}
[piccardi@gont sources]$ \textbf{./echo 192.168.1.141}
Prima riga
Prima riga
Seconda riga dopo l'interruzione
Errore in lettura: No route to host
\end{Console}
%$

Quello che succede in questo è che il programma, dopo aver scritto la seconda
riga, resta bloccato per un tempo molto lungo, prima di dare l'errore
\errcode{EHOSTUNREACH}.  Se andiamo ad osservare con \cmd{strace} cosa accade
nel periodo in cui il programma è bloccato vedremo che stavolta, a differenza
del caso precedente, il programma è bloccato nella lettura dal socket.

Se poi, come nel caso precedente, usiamo l'accortezza di analizzare il
traffico di rete fra client e server con \cmd{tcpdump}, otterremo il seguente
risultato:
\begin{Console}
[root@gont sources]# \textbf{tcpdump src 192.168.1.141 or dst 192.168.1.141 -N -t}
tcpdump: listening on eth0
gont.34685 > anarres.echo: S 1943495663:1943495663(0) win 5840
anarres.echo > gont.34685: S 1215783131:1215783131(0) ack 1943495664 win 5792 
gont.34685 > anarres.echo: . ack 1 win 5840
gont.34685 > anarres.echo: P 1:12(11) ack 1 win 5840 
anarres.echo > gont.34685: . ack 12 win 5792 
anarres.echo > gont.34685: P 1:12(11) ack 12 win 5792 
gont.34685 > anarres.echo: . ack 12 win 5840 
gont.34685 > anarres.echo: P 12:45(33) ack 12 win 5840 
gont.34685 > anarres.echo: P 12:45(33) ack 12 win 5840 
gont.34685 > anarres.echo: P 12:45(33) ack 12 win 5840 
gont.34685 > anarres.echo: P 12:45(33) ack 12 win 5840 
gont.34685 > anarres.echo: P 12:45(33) ack 12 win 5840 
gont.34685 > anarres.echo: P 12:45(33) ack 12 win 5840 
gont.34685 > anarres.echo: P 12:45(33) ack 12 win 5840 
gont.34685 > anarres.echo: P 12:45(33) ack 12 win 5840 
gont.34685 > anarres.echo: P 12:45(33) ack 12 win 5840 
arp who-has anarres tell gont
arp who-has anarres tell gont
arp who-has anarres tell gont
arp who-has anarres tell gont
arp who-has anarres tell gont
arp who-has anarres tell gont
...
\end{Console}

In questo caso l'andamento dei primi sette pacchetti è esattamente lo stesso
di prima. Solo che stavolta, non appena inviata la seconda riga, il programma
si bloccherà nella successiva chiamata a \func{read}, non ottenendo nessuna
risposta. Quello che succede è che nel frattempo il kernel provvede, come
richiesto dal protocollo TCP, a tentare la ritrasmissione della nostra riga un
certo numero di volte, con tempi di attesa crescente fra un tentativo ed il
successivo, per tentare di ristabilire la connessione.

Il risultato finale qui dipende dalla realizzazione dello \textit{stack TCP},
e nel caso di Linux anche dall'impostazione di alcuni dei parametri di sistema
che si trovano in \file{/proc/sys/net/ipv4}, che ne controllano il
comportamento: in questo caso in particolare da
\sysctlrelfile{net/ipv4}{tcp\_retries2} (vedi
sez.~\ref{sec:sock_ipv4_sysctl}). Questo parametro infatti specifica il numero
di volte che deve essere ritentata la ritrasmissione di un pacchetto nel mezzo
di una connessione prima di riportare un errore di timeout.  Il valore
preimpostato è pari a 15, il che comporterebbe 15 tentativi di ritrasmissione,
ma nel nostro caso le cose sono andate diversamente, dato che le
ritrasmissioni registrate da \cmd{tcpdump} sono solo 8; inoltre l'errore
riportato all'uscita del client non è stato \errcode{ETIMEDOUT}, come dovrebbe
essere in questo caso, ma \errcode{EHOSTUNREACH}.

Per capire l'accaduto continuiamo ad analizzare l'output di \cmd{tcpdump}:
esso ci mostra che a un certo punto i tentativi di ritrasmissione del
pacchetto sono cessati, per essere sostituiti da una serie di richieste di
protocollo ARP in cui il client richiede l'indirizzo del server.

Come abbiamo accennato in sez.~\ref{sec:net_tcpip_general} ARP è il protocollo
che si incarica di trovare le corrispondenze fra indirizzo IP e indirizzo
hardware sulla scheda di rete. È evidente allora che nel nostro caso, essendo
client e server sulla stessa rete, è scaduta la voce nella \textit{ARP
  cache}\footnote{la \textit{ARP cache} è una tabella mantenuta internamente
  dal kernel che contiene tutte le corrispondenze fra indirizzi IP e indirizzi
  fisici, ottenute appunto attraverso il protocollo ARP; le voci della tabella
  hanno un tempo di vita limitato, passato il quale scadono e devono essere
  nuovamente richieste.} relativa ad \texttt{anarres}, ed il nostro client ha
iniziato ad effettuare richieste ARP sulla rete per sapere l'IP di
quest'ultimo, che essendo scollegato non poteva rispondere. Anche per questo
tipo di richieste esiste un timeout, per cui dopo un certo numero di tentativi
il meccanismo si è interrotto, e l'errore riportato al programma a questo
punto è stato \errcode{EHOSTUNREACH}, in quanto non si era più in grado di
contattare il server.

Un altro errore possibile in questo tipo di situazione, che si può avere
quando la macchina è su una rete remota, è \errcode{ENETUNREACH}; esso viene
riportato alla ricezione di un pacchetto ICMP di \textit{destination
  unreachable} da parte del router che individua l'interruzione della
connessione. Di nuovo anche qui il risultato finale dipende da quale è il
meccanismo più veloce che porta ad accorgersi del problema.

Se però agiamo sui parametri del kernel, e scriviamo in \file{tcp\_retries2}
un valore di tentativi più basso, possiamo evitare la scadenza della
\textit{ARP cache} e vedere cosa succede. Così se ad esempio richiediamo 4
tentativi di ritrasmissione, l'analisi di \cmd{tcpdump} ci riporterà il
seguente scambio di pacchetti:
\begin{Console}
[root@gont gapil]# \textbf{tcpdump src 192.168.1.141 or dst 192.168.1.141 -N -t}
tcpdump: listening on eth0
gont.34752 > anarres.echo: S 3646972152:3646972152(0) win 5840
anarres.echo > gont.34752: S 2735190336:2735190336(0) ack 3646972153 win 5792 
gont.34752 > anarres.echo: . ack 1 win 5840 
gont.34752 > anarres.echo: P 1:12(11) ack 1 win 5840
anarres.echo > gont.34752: . ack 12 win 5792 
anarres.echo > gont.34752: P 1:12(11) ack 12 win 5792 
gont.34752 > anarres.echo: . ack 12 win 5840 
gont.34752 > anarres.echo: P 12:45(33) ack 12 win 5840 
gont.34752 > anarres.echo: P 12:45(33) ack 12 win 5840 
gont.34752 > anarres.echo: P 12:45(33) ack 12 win 5840 
gont.34752 > anarres.echo: P 12:45(33) ack 12 win 5840 
gont.34752 > anarres.echo: P 12:45(33) ack 12 win 5840 
\end{Console}
e come si vede in questo caso i tentativi di ritrasmissione del pacchetto
iniziale sono proprio 4 (per un totale di 5 voci con quello trasmesso la prima
volta), ed in effetti, dopo un tempo molto più breve rispetto a prima ed in
corrispondenza dell'invio dell'ultimo tentativo, quello che otterremo come
errore all'uscita del client sarà diverso, e cioè:
\begin{Console}
[piccardi@gont sources]$ \textbf{./echo 192.168.1.141}
Prima riga
Prima riga
Seconda riga dopo l'interruzione
Errore in lettura: Connection timed out
\end{Console}
%$
che corrisponde appunto, come ci aspettavamo, alla ricezione di un
\errcode{ETIMEDOUT}.

Analizziamo ora il secondo scenario, in cui si ha un crollo della macchina che
fa da server. Al solito lanciamo il nostro client, scriviamo una prima riga
per verificare che sia tutto a posto, poi stacchiamo il cavo e riavviamo il
server. A questo punto, ritornato attivo il server, scriviamo una seconda
riga. Quello che otterremo in questo caso è:
\begin{Console}
[piccardi@gont sources]$ \textbf{./echo 192.168.1.141}
Prima riga
Prima riga
Seconda riga dopo l'interruzione
Errore in lettura Connection reset by peer
\end{Console}
%$
e l'errore ricevuti da \func{read} stavolta è \errcode{ECONNRESET}. Se al
solito riportiamo l'analisi dei pacchetti effettuata con \cmd{tcpdump},
avremo:
\begin{Console}
[root@gont gapil]# \textbf{tcpdump src 192.168.1.141 or dst 192.168.1.141 -N -t}
tcpdump: listening on eth0
gont.34756 > anarres.echo: S 904864257:904864257(0) win 5840 
anarres.echo > gont.34756: S 4254564871:4254564871(0) ack 904864258 win 5792
gont.34756 > anarres.echo: . ack 1 win 5840 
gont.34756 > anarres.echo: P 1:12(11) ack 1 win 5840
anarres.echo > gont.34756: . ack 12 win 5792 
anarres.echo > gont.34756: P 1:12(11) ack 12 win 5792
gont.34756 > anarres.echo: . ack 12 win 5840 
gont.34756 > anarres.echo: P 12:45(33) ack 12 win 5840
anarres.echo > gont.34756: R 4254564883:4254564883(0) win 0 
\end{Console}

Ancora una volta i primi sette pacchetti sono gli stessi; ma in questo caso
quello che succede dopo lo scambio iniziale è che, non avendo inviato nulla
durante il periodo in cui si è riavviato il server, il client è del tutto
ignaro dell'accaduto per cui quando effettuerà una scrittura, dato che la
macchina server è stata riavviata e che tutti gli stati relativi alle
precedenti connessioni sono completamente persi, anche in presenza di una
nuova istanza del server echo non sarà possibile consegnare i dati in arrivo,
per cui alla loro ricezione il kernel risponderà con un segmento di RST.

Il client da parte sua, dato che neanche in questo caso non è stato emesso un
FIN, dopo aver scritto verrà bloccato nella successiva chiamata a \func{read},
che però adesso ritornerà immediatamente alla ricezione del segmento RST,
riportando appunto come errore \errcode{ECONNRESET}. Occorre precisare che se
si vuole che il client sia in grado di accorgersi del crollo del server anche
quando non sta effettuando uno scambio di dati, è possibile usare una
impostazione speciale del socket (ci torneremo in
sez.~\ref{sec:sock_generic_options}) che provvede all'esecuzione di questo
controllo.


\section{L'uso dell'\textit{I/O multiplexing}}
\label{sec:TCP_sock_multiplexing}

Affronteremo in questa sezione l'utilizzo dell'\textit{I/O multiplexing},
affrontato in sez.~\ref{sec:file_multiplexing}, nell'ambito delle applicazioni
di rete. Già in sez.~\ref{sec:TCP_server_crash} era emerso il problema
relativo al client del servizio \textit{echo} che non era in grado di
accorgersi della terminazione precoce del server, essendo bloccato nella
lettura dei dati immessi da tastiera.

Abbiamo visto in sez.~\ref{sec:file_multiplexing} quali sono le funzionalità
del sistema che ci permettono di tenere sotto controllo più file descriptor in
contemporanea; in quella occasione non abbiamo fatto esempi, in quanto quando
si tratta con file normali questa tipologia di I/O normalmente non viene
usata, è invece un caso tipico delle applicazioni di rete quello di dover
gestire varie connessioni da cui possono arrivare dati comuni in maniera
asincrona, per cui riprenderemo l'argomento in questa sezione.


\subsection{Il comportamento della funzione \func{select} con i socket.}
\label{sec:TCP_sock_select}

Iniziamo con la prima delle funzioni usate per l'\textit{I/O multiplexing},
\func{select}; il suo funzionamento è già stato descritto in dettaglio in
sez.~\ref{sec:file_multiplexing} e non staremo a ripetere quanto detto lì;
sappiamo che la funzione ritorna quando uno o più dei file descriptor messi
sotto controllo è pronto per la relativa operazione.

In quell'occasione non abbiamo però definito cosa si intende per pronto,
infatti per dei normali file, o anche per delle \textit{pipe}, la condizione
di essere pronti per la lettura o la scrittura è ovvia; invece lo è molto meno
nel caso dei socket, visto che possono intervenire tutte una serie di
possibili condizioni di errore dovute alla rete. Occorre allora specificare
chiaramente quali sono le condizioni per cui un socket risulta essere
``\textsl{pronto}'' quando viene passato come membro di uno dei tre
\textit{file descriptor set} usati da \func{select}.

Le condizioni che fanno si che la funzione \func{select} ritorni segnalando
che un socket (che sarà riportato nel primo insieme di file descriptor) è
pronto per la lettura sono le seguenti:
\begin{itemize*}
\item nel buffer di ricezione del socket sono arrivati dei dati in quantità
  sufficiente a superare il valore di una \textsl{soglia di basso livello} (il
  cosiddetto \textit{low watermark}). Questo valore è espresso in numero di
  byte e può essere impostato con l'opzione del socket \const{SO\_RCVLOWAT}
  (tratteremo l'uso di questa opzione in sez.~\ref{sec:sock_generic_options});
  il suo valore di default è 1 per i socket TCP e UDP. In questo caso una
  operazione di lettura avrà successo e leggerà un numero di byte maggiore di
  zero.
\item il lato in lettura della connessione è stato chiuso; si è cioè ricevuto
  un segmento FIN (si ricordi quanto illustrato in
  sez.~\ref{sec:TCP_conn_term}) sulla connessione. In questo caso una
  operazione di lettura avrà successo, ma non risulteranno presenti dati (in
  sostanza \func{read} ritornerà con un valore nullo) per indicare la
  condizione di \textit{end-of-file}.
\item c'è stato un errore sul socket. In questo caso una operazione di lettura
  non si bloccherà ma restituirà una condizione di errore (ad esempio
  \func{read} restituirà $-1$) e imposterà la variabile \var{errno} al relativo
  valore. Vedremo in sez.~\ref{sec:sock_generic_options} come sia possibile
  estrarre e cancellare gli errori pendenti su un socket senza usare
  \func{read} usando l'opzione \const{SO\_ERROR}.
\item quando si sta utilizzando un \textit{listening socket} ed ci sono delle
  connessioni completate. In questo caso la funzione \func{accept} non si
  bloccherà.\footnote{in realtà questo non è sempre vero, come accennato in
    sez.~\ref{sec:TCP_conn_early_abort} una connessione può essere abortita
    dalla ricezione di un segmento RST una volta che è stata completata,
    allora se questo avviene dopo che \func{select} è ritornata, ma prima
    della chiamata ad \func{accept}, quest'ultima, in assenza di altre
    connessioni, potrà bloccarsi.}
\end{itemize*}

Le condizioni che fanno si che la funzione \func{select} ritorni segnalando
che un socket (che sarà riportato nel secondo insieme di file descriptor) è
pronto per la scrittura sono le seguenti:
\begin{itemize*}
\item nel buffer di invio è disponibile una quantità di spazio superiore al
  valore della \textsl{soglia di basso livello} in scrittura ed inoltre o il
  socket è già connesso o non necessita (ad esempio è UDP) di connessione.  Il
  valore della soglia è espresso in numero di byte e può essere impostato con
  l'opzione del socket \const{SO\_SNDLOWAT} (trattata in 
  sez.~\ref{sec:sock_generic_options}); il suo valore di default è 2048 per i
  socket TCP e UDP. In questo caso una operazione di scrittura non si
  bloccherà e restituirà un valore positivo pari al numero di byte accettati
  dal livello di trasporto.
\item il lato in scrittura della connessione è stato chiuso. In questo caso
  una operazione di scrittura sul socket genererà il segnale \signal{SIGPIPE}.
\item c'è stato un errore sul socket. In questo caso una operazione di
  scrittura non si bloccherà ma restituirà una condizione di errore ed
  imposterà opportunamente la variabile \var{errno}. Vedremo in
  sez.~\ref{sec:sock_generic_options} come sia possibile estrarre e cancellare
  errori pendenti su un socket usando l'opzione \const{SO\_ERROR}.
\end{itemize*}

Infine c'è una sola condizione che fa sì che \func{select} ritorni segnalando
che un socket (che sarà riportato nel terzo insieme di file descriptor) ha una
condizione di eccezione pendente, e cioè la ricezione sul socket di
\textsl{dati urgenti} (o \textit{out-of-band}), una caratteristica specifica
dei socket TCP su cui torneremo in sez.~\ref{sec:TCP_urgent_data}.

Si noti come nel caso della lettura \func{select} si applichi anche ad
operazioni che non hanno nulla a che fare con l'I/O di dati come il
riconoscimento della presenza di connessioni pronte, in modo da consentire
anche l'utilizzo di \func{accept} in modalità non bloccante. Si noti infine
come in caso di errore un socket venga sempre riportato come pronto sia per la
lettura che per la scrittura.

Lo scopo dei due valori di soglia per i buffer di ricezione e di invio è
quello di consentire maggiore flessibilità nell'uso di \func{select} da parte
dei programmi, se infatti si sa che una applicazione non è in grado di fare
niente fintanto che non può ricevere o inviare una certa quantità di dati, si
possono utilizzare questi valori per far sì che \func{select} ritorni solo
quando c'è la certezza di avere dati a sufficienza.\footnote{questo tipo di
  controllo è utile di norma solo per la lettura, in quanto in genere le
  operazioni di scrittura sono già controllate dall'applicazione, che sa
  sempre quanti dati invia, mentre non è detto possa conoscere la quantità di
  dati in ricezione; per cui, nella situazione in cui si conosce almeno un
  valore minimo, per evitare la penalizzazione dovuta alla ripetizione delle
  operazioni di lettura per accumulare dati sufficienti, si può lasciare al
  kernel il compito di impostare un minimo al di sotto del quale il socket,
  pur avendo disponibili dei dati, non viene dato per pronto in lettura.}



\subsection{Un esempio di \textit{I/O multiplexing}}
\label{sec:TCP_multiplex_example}

Abbiamo incontrato la problematica tipica che conduce all'uso dell'\textit{I/O
  multiplexing} nella nostra analisi degli errori in
sez.~\ref{sec:TCP_conn_early_abort}, quando il nostro client non era in grado
di rendersi conto di errori sulla connessione essendo impegnato nella attesa
di dati in ingresso dallo \textit{standard input}.

In questo caso il problema è quello di dover tenere sotto controllo due
diversi file descriptor, lo \textit{standard input}, da cui viene letto il
testo che vogliamo inviare al server, e il socket connesso con il server su
cui detto testo sarà scritto e dal quale poi si vorrà ricevere la
risposta. L'uso dell'\textit{I/O multiplexing} consente di tenere sotto controllo
entrambi, senza restare bloccati.

Nel nostro caso quello che ci interessa è non essere bloccati in lettura sullo
\textit{standard input} in caso di errori sulla connessione o chiusura della
stessa da parte del server. Entrambi questi casi possono essere rilevati
usando \func{select}, per quanto detto in sez.~\ref{sec:TCP_sock_select},
mettendo sotto osservazione i file descriptor per la condizione di essere
pronti in lettura: sia infatti che si ricevano dati, che la connessione sia
chiusa regolarmente (con la ricezione di un segmento FIN) che si riceva una
condizione di errore (con un segmento RST) il socket connesso sarà pronto in
lettura (nell'ultimo caso anche in scrittura, ma questo non è necessario ai
nostri scopi).

\begin{figure}[!htb]
  \footnotesize \centering
  \begin{minipage}[c]{\codesamplewidth}
    \includecodesample{listati/ClientEcho_third.c}
  \end{minipage} 
  \normalsize
  \caption{La sezione nel codice della terza versione della funzione
    \func{ClientEcho} usata dal client per il servizio \textit{echo}
    modificata per l'uso di \func{select}.}
  \label{fig:TCP_ClientEcho_third}
\end{figure}

Riprendiamo allora il codice del client, modificandolo per l'uso di
\func{select}. Quello che dobbiamo modificare è la funzione \func{ClientEcho}
di fig.~\ref{fig:TCP_ClientEcho_second}, dato che tutto il resto, che riguarda
le modalità in cui viene stabilita la connessione con il server, resta
assolutamente identico. La nostra nuova versione di \func{ClientEcho}, la
terza della serie, è riportata in fig.~\ref{fig:TCP_ClientEcho_third}, il
codice completo si trova nel file \file{TCP\_echo\_third.c} dei sorgenti
allegati alla guida.

In questo caso la funzione comincia (\texttt{\small 8--9}) con l'azzeramento
del \textit{file descriptor set} \var{fset} e l'impostazione del valore
\var{maxfd}, da passare a \func{select} come massimo per il numero di file
descriptor. Per determinare quest'ultimo si usa la macro \code{max} definita
nel nostro file \file{macro.h} che raccoglie una collezione di macro di
preprocessore di varia utilità.

La funzione prosegue poi (\texttt{\small 10--41}) con il ciclo principale, che
viene ripetuto indefinitamente. Per ogni ciclo si reinizializza
(\texttt{\small 11--12}) il \textit{file descriptor set}, impostando i valori
per il file descriptor associato al socket \var{socket} e per lo
\textit{standard input} (il cui valore si recupera con la funzione
\func{fileno}). Questo è necessario in quanto la successiva (\texttt{\small
  13}) chiamata a \func{select} comporta una modifica dei due bit relativia
questi file, che quindi devono essere reimpostati all'inizio di ogni ciclo.

Si noti come la chiamata a \func{select} venga eseguita usando come primo
argomento il valore di \var{maxfd}, precedentemente calcolato, e passando poi
il solo \textit{file descriptor set} per il controllo dell'attività in
lettura, negli altri argomenti vengono passati tutti puntatori nulli, non
interessando in questo caso né il controllo delle altre attività, né
l'impostazione di un valore di timeout.

Al ritorno di \func{select} si provvede a controllare quale dei due file
descriptor presenta attività in lettura, cominciando (\texttt{\small 14--24})
con il file descriptor associato allo \textit{standard input}. In caso di
attività (quando cioè \macro{FD\_ISSET} ritorna una valore diverso da zero) si
esegue (\texttt{\small 15}) una \func{fgets} per leggere gli eventuali dati
presenti; se non ve ne sono (e la funzione restituisce pertanto un puntatore
nullo) si ritorna immediatamente (\texttt{\small 16}) dato che questo
significa che si è chiuso lo \textit{standard input} e quindi concluso l'utilizzo del
client; altrimenti (\texttt{\small 18--22}) si scrivono i dati appena letti
sul socket, prevedendo una uscita immediata in caso di errore di scrittura.

Controllato lo \textit{standard input} si passa a controllare (\texttt{\small
  25--40}) il socket connesso, in caso di attività (\texttt{\small 26}) si
esegue subito una \func{read} di cui si controlla il valore di ritorno; se
questo è negativo si è avuto un errore e pertanto si esce immediatamente
segnalandolo (\texttt{\small 27--30}), se è nullo significa che il server ha
chiuso la connessione, e di nuovo si esce con stampando prima un messaggio di
avviso (\texttt{\small 31--34}), altrimenti (\texttt{\small 35--39}) si
effettua la terminazione della stringa e la si stampa a sullo \textit{standard
  output}, uscendo in caso di errore, per ripetere il ciclo da capo.

Con questo meccanismo il programma invece di essere bloccato in lettura sullo
\textit{standard input} resta bloccato sulla \func{select}, che ritorna
soltanto quando viene rilevata attività su uno dei due file descriptor posti
sotto controllo.  Questo di norma avviene solo quando si è scritto qualcosa
sullo \textit{standard input}, o quando si riceve dal socket la risposta a quanto si
era appena scritto. 

Ma adesso il client diventa capace di accorgersi immediatamente della
terminazione del server; in tal caso infatti il server chiuderà il socket
connesso, ed alla ricezione del FIN la funzione \func{select} ritornerà (come
illustrato in sez.~\ref{sec:TCP_sock_select}) segnalando una condizione di end
of file, per cui il nostro client potrà uscire immediatamente.

Riprendiamo la situazione affrontata in sez.~\ref{sec:TCP_server_crash},
terminando il server durante una connessione, in questo caso quello che
otterremo, una volta scritta una prima riga ed interrotto il server con un
\texttt{C-c}, sarà:
\begin{Console}
[piccardi@gont sources]$ \textbf{./echo 192.168.1.1}
Prima riga
Prima riga
EOF sul socket
\end{Console}
%$
dove l'ultima riga compare immediatamente dopo aver interrotto il server. Il
nostro client infatti è in grado di accorgersi immediatamente che il socket
connesso è stato chiuso ed uscire immediatamente.

Veniamo allora agli altri scenari di terminazione anomala visti in
sez.~\ref{sec:TCP_conn_crash}. Il primo di questi è l'interruzione fisica della
connessione; in questo caso avremo un comportamento analogo al precedente, in
cui si scrive una riga e non si riceve risposta dal server e non succede
niente fino a quando non si riceve un errore di \errcode{EHOSTUNREACH} o
\errcode{ETIMEDOUT} a seconda dei casi.

La differenza è che stavolta potremo scrivere più righe dopo l'interruzione,
in quanto il nostro client dopo aver inviato i dati non si bloccherà più nella
lettura dal socket, ma nella \func{select}; per questo potrà accettare
ulteriore dati che scriverà di nuovo sul socket, fintanto che c'è spazio sul
buffer di uscita (ecceduto il quale si bloccherà in scrittura). 

Si ricordi infatti che il client non ha modo di determinare se la connessione
è attiva o meno (dato che in molte situazioni reali l'inattività può essere
temporanea).  Tra l'altro se si ricollega la rete prima della scadenza del
timeout, potremo anche verificare come tutto quello che si era scritto viene
poi effettivamente trasmesso non appena la connessione ridiventa attiva, per
cui otterremo qualcosa del tipo:
\begin{Console}
[piccardi@gont sources]$ \textbf{./echo 192.168.1.1}
Prima riga
Prima riga
Seconda riga dopo l'interruzione
Terza riga
Quarta riga
Seconda riga dopo l'interruzione
Terza riga
Quarta riga
\end{Console}
%$
in cui, una volta riconnessa la rete, tutto quello che abbiamo scritto durante
il periodo di disconnessione restituito indietro e stampato immediatamente.

Lo stesso comportamento visto in sez.~\ref{sec:TCP_server_crash} si riottiene
nel caso di un crollo completo della macchina su cui sta il server. In questo
caso di nuovo il client non è in grado di accorgersi di niente dato che si
suppone che il programma server non venga terminato correttamente, ma si
blocchi tutto senza la possibilità di avere l'emissione di un segmento FIN che
segnala la terminazione della connessione. 

Di nuovo fintanto che la connessione non si riattiva (con il riavvio della
macchina del server) il client non è in grado di fare altro che accettare
dell'input e tentare di inviarlo. La differenza in questo caso è che non
appena la connessione ridiventa attiva i dati verranno sì trasmessi, ma
essendo state perse tutte le informazioni relative alle precedenti connessioni
ai tentativi di scrittura del client sarà risposto con un segmento RST che
provocherà il ritorno di \func{select} per la ricezione di un errore di
\errcode{ECONNRESET}.


\subsection{La funzione \func{shutdown}}
\label{sec:TCP_shutdown}

Come spiegato in sez.~\ref{sec:TCP_conn_term} il procedimento di chiusura di
un socket TCP prevede che da entrambe le parti venga emesso un segmento FIN. È
pertanto del tutto normale dal punto di vista del protocollo che uno dei due
capi chiuda la connessione quando l'altro capo la lascia aperta; abbiamo
incontrato questa situazione nei vari scenari critici di
sez.~\ref{sec:TCP_echo_critical}.

\itindbeg{half-close}

È pertanto possibile avere una situazione in cui un capo della connessione,
non avendo più nulla da scrivere, possa chiudere il socket, segnalando così
l'avvenuta terminazione della trasmissione (l'altro capo riceverà infatti un
\textit{end-of-file} in lettura) mentre dall'altra parte si potrà proseguire
la trasmissione dei dati scrivendo sul socket che da quel lato è ancora
aperto.  Questa è quella situazione in cui si dice che il socket è
\textsl{mezzo chiuso} (``\textit{half closed}'').

Il problema che si pone è che se la chiusura del socket è effettuata con la
funzione \func{close}, come spiegato in sez.~\ref{sec:TCP_func_close}, si
perde ogni possibilità di poter leggere quanto l'altro capo può star
continuando a scrivere. Per permettere di segnalare che si è finito con la
scrittura, continuando al contempo a leggere quanto può provenire dall'altro
capo del socket, si può usare la funzione \funcd{shutdown}, il cui prototipo
è:

\begin{funcproto}{
\fhead{sys/socket.h}
\fdecl{int shutdown(int sockfd, int how)}
\fdesc{Chiude un lato della connessione fra due socket.} 
}

{La funzione ritorna $0$ in caso di successo e $-1$ per un errore, nel qual
  caso \var{errno} assumerà uno dei valori: 
  \begin{errlist}
  \item[\errcode{EBADF}] \param{sockfd} non è un file descriptor valido.
  \item[\errcode{EINVAL}] il valore di \param{how} non è valido.
  \item[\errcode{ENOTCONN}] il socket non è connesso.
  \item[\errcode{ENOTSOCK}] il file descriptor non corrisponde a un socket.
  \end{errlist}}
\end{funcproto}

La funzione prende come primo argomento il socket \param{sockfd} su cui si
vuole operare e come secondo argomento un valore intero \param{how} che indica
la modalità di chiusura del socket, quest'ultima può prendere soltanto tre
valori: 
\begin{basedescript}{\desclabelwidth{2.2cm}\desclabelstyle{\nextlinelabel}}
\item[\constd{SHUT\_RD}] chiude il lato in lettura del socket, non sarà più
  possibile leggere dati da esso, tutti gli eventuali dati trasmessi
  dall'altro capo del socket saranno automaticamente scartati dal kernel, che,
  in caso di socket TCP, provvederà comunque ad inviare i relativi segmenti di
  ACK.
\item[\constd{SHUT\_WR}] chiude il lato in scrittura del socket, non sarà più
  possibile scrivere dati su di esso. Nel caso di socket TCP la chiamata causa
  l'emissione di un segmento FIN, secondo la procedura chiamata
  \textit{half-close}. Tutti i dati presenti nel buffer di scrittura prima
  della chiamata saranno inviati, seguiti dalla sequenza di chiusura
  illustrata in sez.~\ref{sec:TCP_conn_term}.
\item[\constd{SHUT\_RDWR}] chiude sia il lato in lettura che quello in
  scrittura del socket. È equivalente alla chiamata in sequenza con
  \const{SHUT\_RD} e \const{SHUT\_WR}.
\end{basedescript}

\itindend{half-close}

Ci si può chiedere quale sia l'utilità di avere introdotto \const{SHUT\_RDWR}
quando questa sembra rendere \func{shutdown} del tutto equivalente ad una
\func{close}. In realtà non è così, esiste infatti un'altra differenza con
\func{close}, più sottile. Finora infatti non ci siamo presi la briga di
sottolineare in maniera esplicita che, come per i file e le \textit{fifo},
anche per i socket possono esserci più riferimenti contemporanei ad uno stesso
socket. 

Per cui si avrebbe potuto avere l'impressione che sia una corrispondenza
univoca fra un socket ed il file descriptor con cui vi si accede. Questo non è
assolutamente vero, (e lo abbiamo già visto nel codice del server di
fig.~\ref{fig:TCP_echo_server_first_code}), ed è invece assolutamente normale
che, come per gli altri oggetti, ci possano essere più file descriptor che
fanno riferimento allo stesso socket.

Allora se avviene uno di questi casi quello che succederà è che la chiamata a
\func{close} darà effettivamente avvio alla sequenza di chiusura di un socket
soltanto quando il numero di riferimenti a quest'ultimo diventerà nullo.
Fintanto che ci sono file descriptor che fanno riferimento ad un socket l'uso
di \func{close} si limiterà a deallocare nel processo corrente il file
descriptor utilizzato, ma il socket resterà pienamente accessibile attraverso
tutti gli altri riferimenti. 

Se torniamo all'esempio originale del server di
fig.~\ref{fig:TCP_echo_server_first_code} abbiamo infatti che ci sono due
\func{close}, una sul socket connesso nel padre, ed una sul socket in ascolto
nel figlio, ma queste non effettuano nessuna chiusura reale di detti socket,
dato che restano altri riferimenti attivi, uno al socket connesso nel figlio
ed uno a quello in ascolto nel padre.

Questo non avviene affatto se si usa \func{shutdown} con argomento
\const{SHUT\_RDWR} al posto di \func{close}; in questo caso infatti la
chiusura del socket viene effettuata immediatamente, indipendentemente dalla
presenza di altri riferimenti attivi, e pertanto sarà efficace anche per tutti
gli altri file descriptor con cui, nello stesso o in altri processi, si fa
riferimento allo stesso socket.

Il caso più comune di uso di \func{shutdown} è comunque quello della chiusura
del lato in scrittura, per segnalare all'altro capo della connessione che si è
concluso l'invio dei dati, restando comunque in grado di ricevere quanto
questi potrà ancora inviarci. Questo è ad esempio l'uso che ci serve per
rendere finalmente completo il nostro esempio sul servizio \textit{echo}. 

Il nostro client infatti presenta ancora un problema, che nell'uso che finora
ne abbiamo fatto non è emerso, ma che ci aspetta dietro l'angolo non appena
usciamo dall'uso interattivo e proviamo ad eseguirlo redirigendo
\textit{standard input} e \textit{standard output}. Così se eseguiamo:
\begin{Console}
[piccardi@gont sources]$ \textbf{./echo 192.168.1.1 < ../fileadv.tex  > copia}
\end{Console}
%$
vedremo che il file \texttt{copia} risulta mancare della parte finale.

Per capire cosa avviene in questo caso occorre tenere presente come avviene la
comunicazione via rete; quando redirigiamo lo \textit{standard input} il
nostro client inizierà a leggere il contenuto del file \texttt{../fileadv.tex}
a blocchi di dimensione massima pari a \texttt{MAXLINE} per poi scriverlo,
alla massima velocità consentitagli dalla rete, sul socket. Dato che la
connessione è con una macchina remota, occorre un certo tempo perché i
pacchetti vi arrivino, vengano processati, e poi tornino
indietro. Considerando trascurabile il tempo di processo, questo tempo è
quello impiegato nella trasmissione via rete, che viene detto RTT (dalla
denominazione inglese \itindex{Round~Trip~Time~(RTT)} \textit{Round Trip
  Time}) ed è quello che viene stimato con l'uso del comando \cmd{ping}.

A questo punto, se torniamo al codice mostrato in
fig.~\ref{fig:TCP_ClientEcho_third}, possiamo vedere che mentre i pacchetti
sono in transito sulla rete il client continua a leggere e a scrivere fintanto
che il file in ingresso finisce. Però non appena viene ricevuto un
\textit{end-of-file} in ingresso il nostro client termina. Nel caso
interattivo, in cui si inviavano brevi stringhe una alla volta, c'era sempre
il tempo di eseguire la lettura completa di quanto il server rimandava
indietro. In questo caso invece, quando il client termina, essendo la
comunicazione saturata e a piena velocità, ci saranno ancora pacchetti in
transito sulla rete che devono arrivare al server e poi tornare indietro, ma
siccome il client esce immediatamente dopo la fine del file in ingresso,
questi non faranno a tempo a completare il percorso e verranno persi.

Per evitare questo tipo di problema, invece di uscire una volta completata la
lettura del file in ingresso, occorre usare \func{shutdown} per effettuare la
chiusura del lato in scrittura del socket. In questo modo il client segnalerà
al server la chiusura del flusso dei dati, ma potrà continuare a leggere
quanto il server gli sta ancora inviando indietro, fino a quando anch'esso,
riconosciuta la chiusura del socket in scrittura da parte del client,
effettuerà la chiusura dalla sua parte. Solo alla ricezione della chiusura del
socket da parte del server il client potrà essere sicuro della ricezione di
tutti i dati e della terminazione effettiva della connessione.

\begin{figure}[!htbp]
  \footnotesize \centering
  \begin{minipage}[c]{\codesamplewidth}
    \includecodesample{listati/ClientEcho.c}
  \end{minipage} 
  \normalsize
  \caption{La sezione nel codice della versione finale della funzione
    \func{ClientEcho}, che usa \func{shutdown} per una conclusione corretta
    della connessione.}
  \label{fig:TCP_ClientEcho}
\end{figure}

Si è allora riportato in fig.~\ref{fig:TCP_ClientEcho} la versione finale
della nostra funzione \func{ClientEcho}, in grado di gestire correttamente
l'intero flusso di dati fra client e server. Il codice completo del client,
comprendente la gestione delle opzioni a riga di comando e le istruzioni per
la creazione della connessione, si trova nel file
\texttt{TCP\_echo\_fourth.c}, distribuito coi sorgenti allegati alla guida.

La nuova versione è molto simile alla precedente di
fig.~\ref{fig:TCP_ClientEcho_third}; la prima differenza è l'introduzione
(\texttt{\small 7}) della variabile \var{eof}, inizializzata ad un valore
nullo, che serve a mantenere traccia dell'avvenuta conclusione della lettura
del file in ingresso.

La seconda modifica (\texttt{\small 12--15}) è stata quella di rendere
subordinata ad un valore nullo di \var{eof} l'impostazione del file descriptor
set per l'osservazione dello \textit{standard input}. Se infatti il valore di
\var{eof} è non nullo significa che si è già raggiunta la fine del file in
ingresso ed è pertanto inutile continuare a tenere sotto controllo lo
\textit{standard input} nella successiva (\texttt{\small 16}) chiamata a
\func{select}.

Le maggiori modifiche rispetto alla precedente versione sono invece nella
gestione (\texttt{\small 18--22}) del caso in cui la lettura con \func{fgets}
restituisce un valore nullo, indice della fine del file. Questa nella
precedente versione causava l'immediato ritorno della funzione; in questo caso
prima (\texttt{\small 19}) si imposta opportunamente \var{eof} ad un valore
non nullo, dopo di che (\texttt{\small 20}) si effettua la chiusura del lato
in scrittura del socket con \func{shutdown}. Infine (\texttt{\small 21}) si
usa la macro \macro{FD\_CLR} per togliere lo \textit{standard input} dal
\textit{file descriptor set}.

In questo modo anche se la lettura del file in ingresso è conclusa, la
funzione non esce dal ciclo principale (\texttt{\small 11--50}), ma continua
ad eseguirlo ripetendo la chiamata a \func{select} per tenere sotto controllo
soltanto il socket connesso, dal quale possono arrivare altri dati, che
saranno letti (\texttt{\small 31}) ed opportunamente trascritti
(\texttt{\small 44--48}) sullo \textit{standard output}.

Il ritorno della funzione, e la conseguente terminazione normale del client,
viene invece adesso gestito all'interno (\texttt{\small 30--49}) della lettura
dei dati dal socket; se infatti dalla lettura del socket si riceve una
condizione di \textit{end-of-file}, la si tratterà (\texttt{\small 36--43}) in
maniera diversa a seconda del valore di \var{eof}. Se infatti questa è diversa
da zero (\texttt{\small 37--39}), essendo stata completata la lettura del file
in ingresso, vorrà dire che anche il server ha concluso la trasmissione dei
dati restanti, e si potrà uscire senza errori, altrimenti si stamperà
(\texttt{\small 40--42}) un messaggio di errore per la chiusura precoce della
connessione.


\subsection{Un server basato sull'\textit{I/O multiplexing}}
\label{sec:TCP_serv_select}

Seguendo di nuovo le orme di Stevens in \cite{UNP1} vediamo ora come con
l'utilizzo dell'\textit{I/O multiplexing} diventi possibile riscrivere
completamente il nostro server \textit{echo} con una architettura
completamente diversa, in modo da evitare di dover creare un nuovo processo
tutte le volte che si ha una connessione.\footnote{ne faremo comunque una
  realizzazione diversa rispetto a quella presentata da Stevens in
  \cite{UNP1}.}

La struttura del nuovo server è illustrata in
fig.~\ref{fig:TCP_echo_multiplex}, in questo caso avremo un solo processo che
ad ogni nuova connessione da parte di un client sul socket in ascolto si
limiterà a registrare l'entrata in uso di un nuovo file descriptor ed
utilizzerà \func{select} per rilevare la presenza di dati in arrivo su tutti i
file descriptor attivi, operando direttamente su ciascuno di essi.

\begin{figure}[!htb]
  \centering \includegraphics[width=13cm]{img/TCPechoMult}
  \caption{Schema del nuovo server echo basato sull'\textit{I/O multiplexing}.}
  \label{fig:TCP_echo_multiplex}
\end{figure}

La sezione principale del codice del nuovo server è illustrata in
fig.~\ref{fig:TCP_SelectEchod}. Si è tralasciata al solito la gestione delle
opzioni, che è identica alla versione precedente. Resta invariata anche tutta
la parte relativa alla gestione dei segnali, degli errori, e della cessione
dei privilegi, così come è identica la gestione della creazione del socket (si
può fare riferimento al codice già illustrato in
sez.~\ref{sec:TCPsimp_server_main}); al solito il codice completo del server è
disponibile coi sorgenti allegati nel file \texttt{select\_echod.c}.

\begin{figure}[!htbp]
  \footnotesize \centering
  \begin{minipage}[c]{\codesamplewidth}
    \includecodesample{listati/select_echod.c}
  \end{minipage} 
  \normalsize
  \caption{La sezione principale della nuova versione di server
    \textit{echo} basato sull'uso della funzione \func{select}.}
  \label{fig:TCP_SelectEchod}
\end{figure}

In questo caso, una volta aperto e messo in ascolto il socket, tutto quello
che ci servirà sarà chiamare \func{select} per rilevare la presenza di nuove
connessioni o di dati in arrivo, e processarli immediatamente. Per realizzare
lo schema mostrato in fig.~\ref{fig:TCP_echo_multiplex} il programma usa una
tabella dei socket connessi mantenuta nel vettore \var{fd\_open} dimensionato
al valore di \const{FD\_SETSIZE}, ed una variabile \var{max\_fd} per
registrare il valore più alto dei file descriptor aperti.

Prima di entrare nel ciclo principale (\texttt{\small 5--53}) la nostra
tabella viene inizializzata (\texttt{\small 2}) a zero (valore che
utilizzeremo come indicazione del fatto che il relativo file descriptor non è
aperto), mentre il valore massimo (\texttt{\small 3}) per i file descriptor
aperti viene impostato a quello del socket in ascolto, in quanto esso è
l'unico file aperto, oltre i tre standard, e pertanto avrà il valore più alto,
che verrà anche (\texttt{\small 4}) inserito nella tabella.

La prima sezione (\texttt{\small 6--8}) del ciclo principale esegue la
costruzione del \textit{file descriptor set} \var{fset} in base ai socket
connessi in un certo momento; all'inizio ci sarà soltanto il socket in ascolto
ma nel prosieguo delle operazioni verranno utilizzati anche tutti i socket
connessi registrati nella tabella \var{fd\_open}.  Dato che la chiamata di
\func{select} modifica il valore del \textit{file descriptor set} è necessario
ripetere (\texttt{\small 6}) ogni volta il suo azzeramento per poi procedere
con il ciclo (\texttt{\small 7--8}) in cui si impostano i socket trovati
attivi.

Per far questo si usa la caratteristica dei file descriptor, descritta in
sez.~\ref{sec:file_open_close}, per cui il kernel associa sempre ad ogni nuovo
file il file descriptor con il valore più basso disponibile. Questo fa sì che
si possa eseguire il ciclo (\texttt{\small 7}) a partire da un valore minimo,
che sarà sempre quello del socket in ascolto, mantenuto in \var{list\_fd},
fino al valore massimo di \var{max\_fd} che dovremo aver cura di tenere
aggiornato.  Dopo di che basterà controllare (\texttt{\small 8}) nella nostra
tabella se il file descriptor è in uso o meno,\footnote{si tenga presente che
  benché il kernel assegni sempre il primo valore libero, si potranno sempre
  avere dei \textsl{buchi} nella nostra tabella dato che nelle operazioni i
  socket saranno aperti e chiusi in corrispondenza della creazione e
  conclusione delle connessioni.} e impostare \var{fset} di conseguenza.

Una volta inizializzato con i socket aperti il nostro \textit{file descriptor
  set} potremo chiamare \func{select} per fargli osservare lo stato degli
stessi (in lettura, presumendo che la scrittura sia sempre consentita). Come
per il precedente esempio di sez.~\ref{sec:TCP_child_hand}, essendo questa
l'unica funzione che può bloccarsi ed essere interrotta da un segnale, la
eseguiremo (\texttt{\small 9--10}) all'interno di un ciclo di \code{while},
che la ripete indefinitamente qualora esca con un errore di \errcode{EINTR}.
Nel caso invece di un errore normale si provvede (\texttt{\small 11--14}) ad
uscire dal programma stampando un messaggio di errore.

Infine quando la funzione ritorna normalmente avremo in \var{n} il numero di
socket da controllare. Nello specifico si danno due casi per cui \func{select}
può ritornare: o si è ricevuta una nuova connessione ed è pronto il socket in
ascolto, sul quale si può eseguire \func{accept}, o c'è attività su uno dei
socket connessi, sui quali si può eseguire \func{read}.

Il primo caso viene trattato immediatamente (\texttt{\small 15--24}): si
controlla (\texttt{\small 15}) che il socket in ascolto sia fra quelli attivi,
nel qual caso anzitutto (\texttt{\small 16}) se ne decrementa il numero
mantenuto nella variabile \var{n}. Poi, inizializzata (\texttt{\small 17}) la
lunghezza della struttura degli indirizzi, si esegue \func{accept} per
ottenere il nuovo socket connesso, controllando che non ci siano errori
(\texttt{\small 18--21}). In questo caso non c'è più la necessità di
controllare per interruzioni dovute a segnali, in quanto siamo sicuri che
\func{accept} non si bloccherà. Per completare la trattazione occorre a questo
punto aggiungere (\texttt{\small 22}) il nuovo file descriptor alla tabella di
quelli connessi, ed inoltre, se è il caso, aggiornare (\texttt{\small 23}) il
valore massimo in \var{max\_fd}.

Una volta controllato l'arrivo di nuove connessioni si passa a verificare se
ci sono dati sui socket connessi, per questo si ripete un ciclo
(\texttt{\small 26--52}) fintanto che il numero di socket attivi indicato
dalla variabile \var{n} resta diverso da zero. In questo modo, se l'unico
socket con attività era quello connesso, avendola opportunamente decrementata
in precedenza, essa risulterà nulla, pertanto il ciclo di verifica verrà
saltato e si ritornerà all'inzizio del ciclo principale, ripetendo, dopo
l'inizializzazione del \textit{file descriptor set} con i nuovi valori nella
tabella, la chiamata di \func{select}.

Se il socket attivo non è quello in ascolto, o ce ne sono comunque anche
altri, il valore di \var{n} non sarà nullo ed il controllo sarà
eseguito. Prima di entrare nel ciclo di veridica comunque si inizializza
(\texttt{\small 25}) il valore della variabile \var{i}, che useremo come
indice nella tabella \var{fd\_open}, al valore minimo, corrispondente al file
descriptor del socket in ascolto.

Il primo passo (\texttt{\small 27}) nella verifica è incrementare il valore
dell'indice \var{i} per posizionarsi sul primo valore possibile per un file
descriptor associato ad un eventuale socket connesso, dopo di che si controlla
(\texttt{\small 28}) se questo è nella tabella dei socket connessi, chiedendo
la ripetizione del ciclo in caso contrario. Altrimenti si passa a verificare
(\texttt{\small 29}) se il file descriptor corrisponde ad uno di quelli
attivi, e nel caso si esegue (\texttt{\small 30}) una lettura, uscendo con un
messaggio in caso di errore (\texttt{\small 31--35}).

Se (\texttt{\small 36}) il numero di byte letti \var{nread} è nullo si è in
presenza di una \textit{end-of-file}, indice che una connessione che si è
chiusa, che deve essere trattata (\texttt{\small 36--45}) opportunamente.  Il
primo passo è chiudere (\texttt{\small 37}) anche il proprio capo del socket e
rimuovere (\texttt{\small 38}) il file descriptor dalla tabella di quelli
aperti, inoltre occorre verificare (\texttt{\small 39}) se il file descriptor
chiuso è quello con il valore più alto, nel qual caso occorre trovare
(\texttt{\small 39--43}) il nuovo massimo, altrimenti (\texttt{\small 44}) si
può ripetere il ciclo da capo per esaminare (se ne restano) ulteriori file
descriptor attivi.

Se però è stato chiuso il file descriptor più alto, dato che la scansione dei
file descriptor attivi viene fatta a partire dal valore più basso, questo
significa che siamo anche arrivati alla fine della scansione, per questo
possiamo utilizzare direttamente il valore dell'indice \var{i} con un ciclo
all'indietro (\texttt{\small 40}) che trova il primo valore per cui la tabella
presenta un file descriptor aperto, e lo imposta (\texttt{\small 41}) come
nuovo massimo, per poi tornare (\texttt{\small 42}) al ciclo principale con un
\code{break}, e rieseguire \func{select}.

Se infine si sono effettivamente letti dei dati dal socket (ultimo caso
rimasto) si potrà invocare immediatamente (\texttt{\small 46})
\func{FullWrite} per riscriverli indietro sul socket stesso, avendo cura di
uscire con un messaggio in caso di errore (\texttt{\small 47--50}). Si noti
che nel ciclo si esegue una sola lettura, contrariamente a quanto fatto con la
precedente versione (si riveda il codice di fig.~\ref{fig:TCP_ServEcho_second})
in cui si continuava a leggere fintanto che non si riceveva un
\textit{end-of-file}, questo perché usando l'\textit{I/O multiplexing} non si
vuole essere bloccati in lettura.  

L'uso di \func{select} ci permette di trattare automaticamente anche il caso
in cui la \func{read} non è stata in grado di leggere tutti i dati presenti
sul socket, dato che alla iterazione successiva \func{select} ritornerà
immediatamente segnalando l'ulteriore disponibilità.

Il nostro server comunque soffre di una vulnerabilità per un attacco di tipo
\textit{Denial of Service}. Il problema è che in caso di blocco di una
qualunque delle funzioni di I/O, non avendo usato processi separati, tutto il
server si ferma e non risponde più a nessuna richiesta. Abbiamo scongiurato
questa evenienza per l'I/O in ingresso con l'uso di \func{select}, ma non vale
altrettanto per l'I/O in uscita. Il problema pertanto può sorgere qualora una
delle chiamate a \func{write} effettuate da \func{FullWrite} si blocchi. 

Con il funzionamento normale questo non accade in quanto il server si limita a
scrivere quanto riceve in ingresso, ma qualora venga utilizzato un client
malevolo che esegua solo scritture e non legga mai indietro l'\textsl{eco} del
server, si potrebbe giungere alla saturazione del buffer di scrittura, ed al
conseguente blocco del server su di una \func{write}.

Le possibili soluzioni in questo caso sono quelle di ritornare ad eseguire il
ciclo di risposta alle richieste all'interno di processi separati, utilizzare
un timeout per le operazioni di scrittura, o eseguire queste ultime in
modalità non bloccante, concludendo le operazioni qualora non vadano a buon
fine.


\subsection{\textit{I/O multiplexing} con \func{poll}}
\label{sec:TCP_serv_poll}

Finora abbiamo trattato le problematiche risolubili con l'\textit{I/O
  multiplexing} impiegando la funzione \func{select}. Questo è quello che
avviene nella maggior parte dei casi, in quanto essa è nata sotto BSD proprio
per affrontare queste problematiche con i socket.  Abbiamo però visto in
sez.~\ref{sec:file_multiplexing} come la funzione \func{poll} possa costituire
una alternativa a \func{select}, con alcuni vantaggi, non soffrendo delle
limitazioni dovute all'uso dei \textit{file descriptor set}.

Ancora una volta in sez.~\ref{sec:file_poll} abbiamo trattato la funzione in
maniera generica, parlando di file descriptor, ma come per \func{select}
quando si ha a che fare con dei socket, il concetto di essere \textsl{pronti}
per l'I/O deve essere specificato nei dettagli per tener conto delle
condizioni della rete. Inoltre deve essere specificato come viene classificato
il traffico nella suddivisione fra dati normali e prioritari. In generale
pertanto:
\begin{itemize}
\item i dati inviati su un socket vengono considerati traffico normale,
  pertanto vengono rilevati alla loro ricezione sull'altro capo da una
  selezione effettuata con \const{POLLIN} o \const{POLLRDNORM};
\item i dati urgenti \textit{out-of-band} (vedi
  sez.~\ref{sec:TCP_urgent_data}) su un socket TCP vengono considerati
  traffico prioritario e vengono rilevati da una condizione \const{POLLIN},
  \const{POLLPRI} o \const{POLLRDBAND}.
\item la chiusura di una connessione (cioè la ricezione di un segmento FIN)
  viene considerato traffico normale, pertanto viene rilevato da una
  condizione \const{POLLIN} o \const{POLLRDNORM}, ma una conseguente chiamata
  a \func{read} restituirà 0.
\item la disponibilità di spazio sul socket per la scrittura di dati viene
  segnalata con una condizione \const{POLLOUT}.
\item quando uno dei due capi del socket chiude un suo lato della connessione
  con \func{shutdown} si riceve una condizione di \const{POLLHUP}.
\item la presenza di un errore sul socket (sia dovuta ad un segmento RST che a
  timeout) viene considerata traffico normale, ma viene segnalata anche dalla
  condizione \const{POLLERR}.
\item la presenza di una nuova connessione su un socket in ascolto può essere
  considerata sia traffico normale che prioritario, nel caso di Linux la
  realizzazione dello \textit{stack TCP} la classifica come normale.
\end{itemize}

Come esempio dell'uso di \func{poll} proviamo allora a riscrivere il server
\textit{echo} secondo lo schema di fig.~\ref{fig:TCP_echo_multiplex} usando
\func{poll} al posto di \func{select}. In questo caso dovremo fare qualche
modifica, per tenere conto della diversa sintassi delle due funzioni, ma la
struttura del programma resta sostanzialmente la stessa.


\begin{figure}[!htbp]
  \footnotesize \centering
  \begin{minipage}[c]{\codesamplewidth}
    \includecodesample{listati/poll_echod.c}
  \end{minipage} 
  \normalsize
  \caption{La sezione principale della nuova versione di server
    \textit{echo} basato sull'uso della funzione \func{poll}.}
  \label{fig:TCP_PollEchod}
\end{figure}

In fig.~\ref{fig:TCP_PollEchod} è riportata la sezione principale della nuova
versione del server, la versione completa del codice è riportata nel file
\texttt{poll\_echod.c} dei sorgenti allegati alla guida. Al solito nella
figura si sono tralasciate la gestione delle opzioni, la creazione del socket
in ascolto, la cessione dei privilegi e le operazioni necessarie a far
funzionare il programma come demone, privilegiando la sezione principale del
programma.

Come per il precedente server basato su \func{select} il primo passo
(\texttt{\small 2--8}) è quello di inizializzare le variabili necessarie. Dato
che in questo caso dovremo usare un vettore di strutture occorre anzitutto
(\texttt{\small 2}) allocare la memoria necessaria utilizzando il numero
massimo \var{n} di socket osservabili, che viene impostato attraverso
l'opzione \texttt{-n} ed ha un valore di default di 256. 

Dopo di che si preimposta (\texttt{\small 3}) il valore \var{max\_fd} del file
descriptor aperto con valore più alto a quello del socket in ascolto (al
momento l'unico), e si provvede (\texttt{\small 4--7}) ad inizializzare le
strutture, disabilitando l'osservazione (\texttt{\small 5}) con un valore
negativo del campo \var{fd}, ma predisponendo (\texttt{\small 6}) il campo
\var{events} per l'osservazione dei dati normali con \const{POLLRDNORM}.
Infine (\texttt{\small 8}) si attiva l'osservazione del socket in ascolto
inizializzando la corrispondente struttura. Questo metodo comporta, in
modalità interattiva, lo spreco di tre strutture (quelle relative a
\textit{standard input}, \textit{standard output} e \textit{standard error})
che non vengono mai utilizzate in quanto la prima è sempre quella relativa al
socket in ascolto.

Una volta completata l'inizializzazione tutto il lavoro viene svolto
all'interno del ciclo principale \texttt{\small 9--53}) che ha una struttura
sostanzialmente identica a quello usato per il precedente esempio basato su
\func{select}. La prima istruzione (\texttt{\small 10}) è quella di eseguire
\func{poll} all'interno di un ciclo che la ripete qualora venisse interrotta
da un segnale, da cui si esce soltanto quando la funzione ritorna restituendo
nella variabile \var{n} il numero di file descriptor trovati attivi.  Qualora
invece si sia ottenuto un errore si procede (\texttt{\small 11--14}) alla
terminazione immediata del processo provvedendo a stampare una descrizione
dello stesso.

Una volta ottenuta dell'attività su un file descriptor si hanno di nuovo due
possibilità. La prima è che ci sia attività sul socket in ascolto, indice di
una nuova connessione, nel qual caso si controlla (\texttt{\small 17}) se il
campo \var{revents} della relativa struttura è attivo; se è così si provvede
(\texttt{\small 16}) a decrementare la variabile \var{n} (che assume il
significato di numero di file descriptor attivi rimasti da controllare) per
poi (\texttt{\small 17--21}) effettuare la chiamata ad \func{accept},
terminando il processo in caso di errore. Se la chiamata ad \func{accept} ha
successo si procede attivando (\texttt{\small 22}) la struttura relativa al
nuovo file descriptor da essa ottenuto, modificando (\texttt{\small 23})
infine quando necessario il valore massimo dei file descriptor aperti
mantenuto in \var{max\_fd}.

La seconda possibilità è che vi sia dell'attività su uno dei socket aperti in
precedenza, nel qual caso si inizializza (\texttt{\small 25}) l'indice \var{i}
del vettore delle strutture \struct{pollfd} al valore del socket in ascolto,
dato che gli ulteriori socket aperti avranno comunque un valore superiore.  Il
ciclo (\texttt{\small 26--52}) prosegue fintanto che il numero di file
descriptor attivi, mantenuto nella variabile \var{n}, è diverso da zero. Se
pertanto ci sono ancora socket attivi da individuare si comincia con
l'incrementare (\texttt{\small 27}) l'indice e controllare (\texttt{\small
  28}) se corrisponde ad un file descriptor in uso analizzando il valore del
campo \var{fd} della relativa struttura e chiudendo immediatamente il ciclo
qualora non lo sia. Se invece il file descriptor è in uso si verifica
(\texttt{\small 29}) se c'è stata attività controllando il campo
\var{revents}.

Di nuovo se non si verifica la presenza di attività il ciclo si chiude subito,
altrimenti si provvederà (\texttt{\small 30}) a decrementare il numero \var{n}
di file descriptor attivi da controllare e ad eseguire (\texttt{\small 31}) la
lettura, ed in caso di errore (\texttt{\small 32--35}) al solito lo si
notificherà uscendo immediatamente. Qualora invece si ottenga una condizione
di \textit{end-of-file} (\texttt{\small 36--45}) si provvederà a chiudere
(\texttt{\small 37}) anche il nostro capo del socket e a marcarlo
(\texttt{\small 38}) come inutilizzato nella struttura ad esso associata.
Infine dovrà essere ricalcolato (\texttt{\small 39--43}) un eventuale nuovo
valore di \var{max\_fd}. L'ultimo passo è chiudere (\texttt{\small 44}) il
ciclo in quanto in questo caso non c'è più niente da riscrivere all'indietro
sul socket.

Se invece si sono letti dei dati si provvede (\texttt{\small 46}) ad
effettuarne la riscrittura all'indietro, con il solito controllo ed eventuale
uscita e notifica in caso di errore (\texttt{\small 47--51}).

Come si può notare la logica del programma è identica a quella vista in
fig.~\ref{fig:TCP_SelectEchod} per l'analogo server basato su \func{select};
la sola differenza significativa è che in questo caso non c'è bisogno di
rigenerare i \textit{file descriptor set} in quanto l'uscita è indipendente
dai dati in ingresso. Si applicano comunque anche a questo server le
considerazioni finali di sez.~\ref{sec:TCP_serv_select}.


%\subsection{\textit{I/O multiplexing} con \textit{epoll}}
%\label{sec:TCP_serv_epoll}

%Da fare.

% TODO fare esempio con epoll



% LocalWords:  socket TCP client dell'I multiplexing stream three way handshake
% LocalWords:  header stack kernel SYN ACK URG syncronize sez bind listen fig
% LocalWords:  accept connect active acknowledge l'acknowledge nell'header MSS
% LocalWords:  sequence number l'acknowledgement dell'header options l'header
% LocalWords:  option MMS segment size MAXSEG window advertised Mbit sec nell'
% LocalWords:  timestamp RFC long fat close of l'end l'ACK half shutdown CLOSED
% LocalWords:  netstat SENT ESTABLISHED WAIT IPv Ethernet piggybacking UDP MSL
% LocalWords:  l'overhead Stevens Lifetime router hop limit TTL to live RST SSH
% LocalWords:  routing dell'MSL l'IP multitasking well known port ephemeral BSD
% LocalWords:  ports dall' IANA Assigned Authority like glibc netinet IPPORT AF
% LocalWords:  RESERVED USERRESERVED rsh rlogin pair socketpair Local Address
% LocalWords:  Foreing DNS caching INADDR ANY multihoming loopback ssh fuser ip
% LocalWords:  lsof SOCK sys int sockfd const struct sockaddr serv addr socklen
% LocalWords:  addrlen errno EBADF descriptor EINVAL ENOTSOCK EACCES EADDRINUSE
% LocalWords:  EADDRNOTAVAIL EFAULT ENOTDIR ENOENT ENOMEM ELOOP ENOSR EROFS RPC
% LocalWords:  portmapper htonl tab endianness BROADCAST broadcast any extern fd
% LocalWords:  ADRR INIT DGRAM SEQPACKET servaddr ECONNREFUSED ETIMEDOUT EAGAIN
% LocalWords:  ENETUNREACH EINPROGRESS EALREADY EAFNOSUPPORT EPERM EISCONN proc
% LocalWords:  sysctl filesystem syn retries reset ICMP backlog EOPNOTSUPP RECV
% LocalWords:  connection queue dell'ACK flood spoofing syncookies SOMAXCONN CR
% LocalWords:  RDM EWOULDBLOCK firewall ENOBUFS EINTR EMFILE ECONNABORTED NULL
% LocalWords:  ESOCKTNOSUPPORT EPROTONOSUPPORT ERESTARTSYS connected listening
% LocalWords:  DECnet read write NONBLOCK fcntl getsockname getpeername name ps
% LocalWords:  namelen namlen ENOTCONN exec inetd POSIX daytime FullRead count
% LocalWords:  BUF FullWrite system call INET perror htons inet pton ctime FTP
% LocalWords:  fputs carriage return line feed superdemone daytimed sleep fork
% LocalWords:  daemon cunc logging list conn sock exit snprintf ntop ntohs echo
% LocalWords:  crash superserver L'RFC first ClientEcho stdin stdout fgets main
% LocalWords:  MAXLINE initd echod ServEcho setgid short nogroup nobody setuid
% LocalWords:  demonize PrintErr syslog wrapper log error root RTT EOF ctrl ack
% LocalWords:  while SIGCHLD Signal RESTART sigaction SignalRestart SigHand win
% LocalWords:  flags select recvfrom debug second compat waiting Nsec ENETDOWN
% LocalWords:  EPROTO ENOPROTOOPT EHOSTDOWN ENONET EHOSTUNREACH LINGER tcpdump
% LocalWords:  ECONNRESET advertising PSH SIGTERM strace SIGPIPE gets tcp ARP
% LocalWords:  cache anarres destination unreachable l'I low watermark RCVLOWAT
% LocalWords:  SNDLOWAT third fset maxfd fileno ISSET closed how SHUT RD WR eof
% LocalWords:  RDWR fifo Trip ping fourth CLR sull'I SETSIZE nread break Denial
% LocalWords:  Service poll POLLIN POLLRDNORM POLLPRI POLLRDBAND POLLOUT events
% LocalWords:  POLLHUP POLLERR revents pollfd Di scaling SYNCNT DoS

%%% Local Variables: 
%%% mode: latex
%%% TeX-master: "gapil"
%%% End: 

